\documentclass[11pt,a4paper,oneside]{report}
% Layout adapted from the supplied SGU Assignment_report.tex.
% XeLaTeX supports native Vietnamese text and modern Unicode fonts.
\usepackage[a4paper,left=24mm,right=22mm,top=33mm,bottom=28mm,headheight=40pt,headsep=9mm,footskip=13mm]{geometry}
\usepackage{fontspec}
\setmainfont{Tinos}
\setsansfont{Noto Sans}
\setmonofont{DejaVu Sans Mono}[Scale=0.85]
\usepackage{amsmath,amssymb}
\usepackage{longtable,booktabs,array,calc,tabularx,multirow}
\usepackage{graphicx,pdfpages,rotating,pdflscape}
\usepackage[table]{xcolor}
\usepackage{float}
\usepackage{enumitem}
\usepackage{fancyhdr,lastpage}
\usepackage{setspace}
\usepackage{etoolbox}
\usepackage{caption}
\usepackage{hyperref}
\definecolor{sgublue}{HTML}{173FAB}
\hypersetup{colorlinks=true,linkcolor=black,urlcolor=sgublue,citecolor=black,pdftitle={BrewLite - Assignment 1},pdfauthor={Nhóm 3 - Trường Đại học Sài Gòn}}
\setcounter{secnumdepth}{0} % Chapter numbers only; preserve source's own section numbering.
\setcounter{tocdepth}{2}
\setlength{\parindent}{0pt}
\setlength{\parskip}{5pt}
\setlength{\tabcolsep}{3pt}
\renewcommand{\arraystretch}{1.10}
\setlength{\LTpre}{5pt}
\setlength{\LTpost}{10pt}
\AtBeginEnvironment{longtable}{\small}
\emergencystretch=3em
\sloppy
\setlist{itemsep=1pt,parsep=0pt,topsep=4pt}
\newcommand{\tightlist}{\setlength{\itemsep}{0pt}\setlength{\parskip}{0pt}}
\pagestyle{fancy}
\fancyhf{}
\fancyhead[L]{%
 \begin{tabular}{@{}cl@{}}
 \includegraphics[width=9mm]{assets/branding/logoITSGUsmall.png} &
 \begin{tabular}{@{}l@{}}{\scriptsize\sffamily\bfseries Trường Đại học Sài Gòn}\\[-1pt]{\scriptsize\sffamily\bfseries Khoa Công nghệ Thông tin}\end{tabular}
 \end{tabular}}
\fancyfoot[L]{{\scriptsize\sffamily Báo cáo Công nghệ phần mềm -- BrewLite | Nhóm 3}}
\fancyfoot[R]{{\scriptsize\sffamily Trang \thepage /\pageref{LastPage}}}
\renewcommand{\headrulewidth}{0.3pt}
\renewcommand{\footrulewidth}{0.3pt}
\fancypagestyle{plain}{\fancyhf{}\fancyfoot[L]{\scriptsize\sffamily BrewLite -- Assignment 1 | Nhóm 3}\fancyfoot[R]{\scriptsize\sffamily Trang \thepage /\pageref{LastPage}}\renewcommand{\headrulewidth}{0pt}\renewcommand{\footrulewidth}{0.3pt}}
\captionsetup{font=small,labelfont=bf}

\renewcommand{\contentsname}{Mục lục}

\begin{document}
\begin{titlepage}
\begin{center}
{\bfseries\small TRƯỜNG ĐẠI HỌC SÀI GÒN}\\[2pt]
{\bfseries\small KHOA CÔNG NGHỆ THÔNG TIN}\\[16mm]
\includegraphics[width=33mm]{assets/branding/logoITSGU.png}\\[17mm]
{\color{sgublue}\LARGE\bfseries CÔNG NGHỆ PHẦN MỀM}\\[7mm]
{\color{sgublue}\rule{0.9\textwidth}{0.8pt}}\\[4mm]
{\color{sgublue}\large\bfseries TÀI LIỆU ASSIGNMENT 1}\\[3mm]
{\color{sgublue}\Huge\bfseries BrewLite}\\[4mm]
\parbox{0.94\textwidth}{\centering\color{sgublue}\Large\bfseries Requirements Specification}\par\vspace{5mm}
{\color{sgublue}\rule{0.9\textwidth}{0.8pt}}\\[15mm]
\end{center}
\begin{flushright}
\begin{tabular}{@{}rl@{}}
\textbf{Giảng viên:} & Từ Lãng Phiêu \\
\textbf{Nhóm:} & Nhóm 3 \\
\textbf{Môn học:} & Công nghệ phần mềm \\
\end{tabular}
\end{flushright}
\vfill
\begin{center}
{\small\bfseries THÀNH VIÊN THỰC HIỆN}\\[3mm]
\begin{tabular}{@{}ll@{}}
 Nguyễn Châu Hoàng Khang & 3124410143 \\
 Hoàng Văn Lập & 3124410184 \\
 Nguyễn Khải Hoàng & 3124410106 \\
 Nguyễn Đăng Thụy & 3124410352 \\
 Cao Nguyễn Quốc Khánh & 3124410149 \\
 Lâm Chí Khang & 3124410142 \\

\end{tabular}
\vfill
{\small TP. HỒ CHÍ MINH, THÁNG 10/2026}
\end{center}
\end{titlepage}
\pagenumbering{roman}
\chapter*{DANH SÁCH THÀNH VIÊN VÀ PHÂN CÔNG}
\noindent Nhóm 3 -- Trường Đại học Sài Gòn -- Môn Công nghệ phần mềm.
\par\smallskip
\textit{Cột Phân công để Chưa cập nhật khi chưa có phân nhiệm chính thức; vai trò lấy theo Project Charter.}
\vspace{4mm}
\begingroup\small
\begin{longtable}{|>{\centering\arraybackslash}p{8mm}|p{41mm}|p{25mm}|p{36mm}|p{36mm}|}
\hline
\rowcolor{blue!9}\textbf{STT} & \textbf{Họ và tên} & \textbf{MSSV} & \textbf{Vai trò} & \textbf{Phân công} \\ \hline
\endfirsthead
\hline\rowcolor{blue!9}\textbf{STT} & \textbf{Họ và tên} & \textbf{MSSV} & \textbf{Vai trò} & \textbf{Phân công} \\ \hline
\endhead
1 & Nguyễn Châu Hoàng Khang & 3124410143 & Team Leader / Developer & \textit{Chưa cập nhật} \\ \hline
2 & Hoàng Văn Lập & 3124410184 & Developer & \textit{Chưa cập nhật} \\ \hline
3 & Nguyễn Khải Hoàng & 3124410106 & Developer & \textit{Chưa cập nhật} \\ \hline
4 & Nguyễn Đăng Thụy & 3124410352 & Developer & \textit{Chưa cập nhật} \\ \hline
5 & Cao Nguyễn Quốc Khánh & 3124410149 & Developer & \textit{Chưa cập nhật} \\ \hline
6 & Lâm Chí Khang & 3124410142 & Developer & \textit{Chưa cập nhật} \\ \hline
\end{longtable}
\endgroup
\clearpage
\tableofcontents
\clearpage
\pagenumbering{arabic}
\chapter{Requirements Specification}
% Auto-converted from BrewLite_Requirements_Specification.docx via pandoc 3.1.11.1.
% Original wording and table data are retained.
REQUIREMENTS SPECIFICATION

BrewLite \textbar{} Phân tích \& Đặc tả Yêu cầu Hệ thống

Phiên bản: v3.0 \textbar{} Ngày lập: 09/10/2026 \textbar{} Trạng thái: Đã chốt

\textbf{Mục tiêu}

Mô tả đầy đủ hành vi nghiệp vụ, permission, lifecycle, failure handling và tiêu chí kiểm chứng để toàn bộ Use Case, ERD, State, Sequence, API, UI và Test có cùng một nguồn yêu cầu

\textbf{Document Control}

\begin{longtable}[]{@{}
  >{\raggedright\arraybackslash}p{(\columnwidth - 2\tabcolsep) * \real{0.2500}}
  >{\raggedright\arraybackslash}p{(\columnwidth - 2\tabcolsep) * \real{0.7500}}@{}}
\toprule\noalign{}
\begin{minipage}[b]{\linewidth}\raggedright
\textbf{Thông tin}
\end{minipage} & \begin{minipage}[b]{\linewidth}\raggedright
\textbf{Giá trị}
\end{minipage} \\
\midrule\noalign{}
\endhead
\bottomrule\noalign{}
\endlastfoot
\textbf{Tên tài liệu} & Tài liệu Phân tích \& Đặc tả Yêu cầu Hệ thống \\
\textbf{Dự án} & BrewLite - Ứng dụng đặt cà phê và thanh toán không tiền mặt \\
\textbf{Phiên bản} & v3.0 \\
\textbf{Ngày cập nhật} & 14/10/2026 \\
\textbf{Trạng thái} & Đã chốt \\
\textbf{Phạm vi sử dụng} & Nguồn yêu cầu cho phân tích / thiết kế / phát triển / kiểm thử \\
\end{longtable}

\textbf{1. Mục đích và nguyên tắc sử dụng tài liệu}

Tài liệu không khóa cứng tên hàm hay tên biến, nhưng mô tả đủ rõ hành vi mong muốn để các thành viên (frontend/backend) không tự phát minh ra các business rule (quy tắc nghiệp vụ) trong lúc code.

Tài liệu định nghĩa rõ các hành vi cốt lõi: luồng khách hàng (Customer), quản lý giỏ hàng, nghiệp vụ xử lý đơn hàng (state machine), thanh toán an toàn (idempotency), và quản lý tồn kho đồng thời (concurrency).

\textbf{2. Quy ước yêu cầu}

\begin{longtable}[]{@{}
  >{\raggedright\arraybackslash}p{(\columnwidth - 2\tabcolsep) * \real{0.2500}}
  >{\raggedright\arraybackslash}p{(\columnwidth - 2\tabcolsep) * \real{0.7500}}@{}}
\toprule\noalign{}
\begin{minipage}[b]{\linewidth}\raggedright
\textbf{Ký hiệu}
\end{minipage} & \begin{minipage}[b]{\linewidth}\raggedright
\textbf{Ý nghĩa}
\end{minipage} \\
\midrule\noalign{}
\endhead
\bottomrule\noalign{}
\endlastfoot
\textbf{MUST} & Bắt buộc cho core flow của hệ thống MVP (10 task bắt buộc của môn học). \\
\textbf{SHOULD} & Nên có để đảm bảo tính hoàn thiện của một ứng dụng thực tế. \\
\textbf{MOCK} & Tích hợp giả lập, không chạy giao dịch thật (đặc biệt cho cổng thanh toán). \\
\textbf{OUT-OF-SCOPE} & Nằm ngoài phạm vi phiên bản môn học. \\
\textbf{CONFIRMED} & Yêu cầu bắt nguồn trực tiếp từ đề bài của giảng viên (ví dụ: 10 task bắt buộc) hoặc đã chốt trong Project Charter. \\
\textbf{DERIVED} & Yêu cầu do nhóm tự suy luận ra từ logic nghiệp vụ để hệ thống không bị lỗi (ví dụ: backend phải tính lại tổng tiền giỏ hàng). \\
\end{longtable}

Mỗi yêu cầu chức năng có: actor, precondition/điều kiện, phân tích nghiệp vụ, vấn đề/rủi ro, hướng xử lý, ảnh hưởng tới thiết kế và tiêu chí kiểm chứng. Đây là cấu trúc dùng trực tiếp để viết Use Case Specification và Kịch bản kiểm thử.

\textbf{3. Bối cảnh và ranh giới hệ thống}

BrewLite là ứng dụng đặt cà phê theo mô hình Single-Vendor (một cửa hàng duy nhất).

\textbf{Web Client (Next.js):} Giao diện phục vụ Guest và Customer. Chỉ đóng vai trò hiển thị và gửi yêu cầu, không phải nguồn quyết định cuối cùng về giá hay số lượng tồn kho.

\textbf{Backend API (NestJS):} Nguồn dữ liệu authoritative (chính chủ) quản lý tập trung Product, User, Order, Payment, Loyalty trong PostgreSQL (truy cập qua Prisma); hệ thống chạy bằng Docker Compose. Backend là nơi quyết định toàn bộ logic nghiệp vụ.

\textbf{Mock Payment:} Hệ thống thanh toán giả lập nội bộ, không tích hợp VNPay/Momo thật. Ranh giới quan trọng: Dù là mock, backend vẫn phải kiểm tra Idempotency-Key và không tin tưởng tuyệt đối vào kết quả do Client gửi lên.

\textbf{Hệ thống không bao gồm:} Kênh Seller/Admin phức tạp, giao diện Barista, vận chuyển (Last-Mile), Microservices.

\textbf{4. Phạm vi nghiệp vụ}

\textbf{4.1. Trong phạm vi}

\begin{itemize}
\item
  \textbf{Identity:} Đăng ký, đăng nhập JWT, băm mật khẩu (bcrypt).
\item
  \textbf{Menu \& Catalog:} Xem danh sách đồ uống, chi tiết sản phẩm, tùy chọn Size (S/M/L) và Topping, tính giá động.
\item
  \textbf{Cart:} thêm, cập nhật số lượng, xóa, tự tính tổng tạm; server tính lại khi tạo đơn.
\item
  \textbf{Order:} Tạo đơn từ giỏ hàng, kiểm tra hợp lệ, chống bán vượt tồn kho (oversell), lưu trạng thái PENDING và trả mã đơn.
\item
  \textbf{Payment:} Thanh toán giả lập qua Ví/Thẻ, xử lý Idempotency-Key để tránh giao dịch lặp, xử lý thành công/thất bại (PAID / PAYMENT\_FAILED), giữ giỏ khi lỗi; xem xác nhận đơn và lịch sử đơn.
\item
  \textbf{Nghiệp vụ nâng cao:} State machine cho Order; coupon/khuyến mãi; cộng điểm tích lũy sau khi thanh toán thành công.
\item
  \textbf{Hạ tầng và bàn giao:} kiểm thử, Docker Compose, README, demo nghiệm thu.
\end{itemize}

\textbf{4.2 Ngoài phạm vi}

\begin{itemize}
\item
  Tích hợp giao dịch tiền thật với ví điện tử/thẻ (chỉ dùng mock payment gateway).
\item
  Hệ thống vận chuyển, định vị thời gian thực.
\item
  Microservices, multi-vendor, kênh Seller/Admin phức tạp, ứng dụng mobile native.
\item
  Giao diện Barista nhận/pha chế đơn: phần mở rộng, không thuộc MVP; tuy nhiên vòng đời trạng thái đơn vẫn được mô hình hóa theo Task 10.
\item
  Hệ thống ERP/POS quản lý nguyên liệu đầy đủ; chỉ triển khai cơ chế tồn kho phục vụ chống oversell.
\end{itemize}

\textbf{5. Stakeholder và actor}

\begin{longtable}[]{@{}
  >{\raggedright\arraybackslash}p{(\columnwidth - 6\tabcolsep) * \real{0.1706}}
  >{\raggedright\arraybackslash}p{(\columnwidth - 6\tabcolsep) * \real{0.2044}}
  >{\raggedright\arraybackslash}p{(\columnwidth - 6\tabcolsep) * \real{0.3247}}
  >{\raggedright\arraybackslash}p{(\columnwidth - 6\tabcolsep) * \real{0.3002}}@{}}
\toprule\noalign{}
\begin{minipage}[b]{\linewidth}\raggedright
\textbf{Actor/}

\textbf{Stakeholder}
\end{minipage} & \begin{minipage}[b]{\linewidth}\raggedright
\textbf{Mục tiêu}
\end{minipage} & \begin{minipage}[b]{\linewidth}\raggedright
\textbf{Dữ liệu/khả năng chính}
\end{minipage} & \begin{minipage}[b]{\linewidth}\raggedright
\textbf{Rủi ro phân tích cần chú ý}
\end{minipage} \\
\midrule\noalign{}
\endhead
\bottomrule\noalign{}
\endlastfoot
Guest & Xem menu không cần đăng nhập & Danh sách đồ uống, chi tiết sản phẩm size/topping. & Phải đăng nhập trước khi tạo Order. \\
Customer & Đặt đồ uống, thanh toán mock, tích điểm & Tài khoản, giỏ hàng, đơn, lịch sử đơn, điểm & Stale cart, đặt trùng, trừ tiền mock 2 lần, xem đơn người khác (IDOR(1)), đặt đồng thời vượt tồn kho (oversell(2)). \\
Giảng viên & Đánh giá, nghiệm thu & Luồng E2E, Postman, Docker Compose, README & API chậm, sai luồng state machine (state transition(3)). \\
Mock Payment & Trả kết quả thanh toán giả lập & Payment request/response & Request lặp gây tạo 2 giao dịch (idempotency(4)). \\
Quản lý cửa hàng / PO mô phỏng & Xác định giá trị, ưu tiên backlog, chấp nhận increment & Yêu cầu, thứ tự ưu tiên & Scope creep (Charter R1). \\
Nhân viên quầy / Barista & Nhận và pha chế đơn & Trạng thái đơn (PREPARING, READY, COMPLETED) & Không có giao diện trong MVP; cách kích hoạt chuyển trạng thái được đặc tả ở mục 11 (xem FR-ORD-003). \\
Nhóm phát triển & Phân tích, thiết kế, code, review, test & Source, tài liệu, Jira & Tài liệu lệch nhau (Charter R2), tích hợp FE/BE (Charter R3). \\
\end{longtable}

(1): IDOR (Insecure Direct Object Reference): Customer có thể thay đổi mã định danh trong request (ví dụ orderId) để truy cập đơn hàng của người khác nếu backend không kiểm tra ownership. BrewLite phòng tránh bằng cách luôn lọc theo tài khoản trong JWT (ví dụ /orders/me).

(2): Oversell: bán vượt số lượng tồn kho, thường xảy ra khi nhiều người đặt đồng thời món cuối cùng mà hệ thống chỉ kiểm tra rồi mới trừ tồn.

(3): Quy tắc chuyển đổi trạng thái (state transition): xác định các trạng thái hợp lệ của đơn hàng (PENDING, PAID, PREPARING, READY, COMPLETED, PAYMENT\_FAILED, CANCELLED) và các bước chuyển được phép; mọi bước chuyển khác phải bị backend chặn.

(4): Idempotency: cùng một yêu cầu thanh toán (cùng Idempotency-Key và cùng nội dung) được gửi nhiều lần chỉ tạo ra một kết quả duy nhất, không tạo giao dịch kép.

\textbf{6. Phân rã chức năng cấp cao}

\begin{longtable}[]{@{}
  >{\raggedright\arraybackslash}p{(\columnwidth - 4\tabcolsep) * \real{0.2200}}
  >{\raggedright\arraybackslash}p{(\columnwidth - 4\tabcolsep) * \real{0.3000}}
  >{\raggedright\arraybackslash}p{(\columnwidth - 4\tabcolsep) * \real{0.4800}}@{}}
\toprule\noalign{}
\begin{minipage}[b]{\linewidth}\raggedright
\textbf{Mã nhóm}
\end{minipage} & \begin{minipage}[b]{\linewidth}\raggedright
\textbf{Chức năng}
\end{minipage} & \begin{minipage}[b]{\linewidth}\raggedright
\textbf{Đầu ra phân tích chính}
\end{minipage} \\
\midrule\noalign{}
\endhead
\bottomrule\noalign{}
\endlastfoot
F01 & Identity \& Account & Đăng ký, đăng nhập, JWT verification, bảo vệ API cần xác thực. \\
F02 & Catalog & Danh sách sản phẩm, chi tiết, tùy chọn Size/Topping, tính giá, nạp dữ liệu mẫu (seed). \\
F03 & Cart & Client state (Zustand), thêm/sửa/xóa, tổng tạm tính, giữ trạng thái giỏ, kiểm tra lại khi tạo đơn. \\
F04 & Order & Tạo đơn PENDING, kiểm tra lại giỏ và giá, snapshot dòng hàng, State Machine, lịch sử, hủy đơn, thử lại thanh toán, giới hạn một đơn chưa thanh toán mỗi khách và tự hủy khi quá hạn. \\
F05 & Inventory & Quản lý tồn kho không âm, chống oversell khi đặt đồng thời, giữ tồn khi tạo đơn và qua các lần thanh toán thất bại, hoàn đúng một lần khi đơn chuyển CANCELLED, hỗ trợ đối soát biến động. \\
F06 & Payment & Mock Payment Ví/Thẻ, kiểm soát Idempotency-Key, nhiều lần thử Payment trên một Order, tối đa một SUCCESS, xử lý PAID/PAYMENT\_FAILED, xác nhận đơn và kịch bản giả lập. \\
F07 & Loyalty \& Promo & Kiểm tra mã giảm giá, tính giảm giá, giới hạn lượt dùng, cộng điểm sau thanh toán thành công (không cộng trùng), xem điểm, dữ liệu mã mẫu. \\
\end{longtable}

\textbf{7. Yêu cầu chức năng chi tiết}

\textbf{7.1. Identity \& Account}

\textbf{Actor/phạm vi chính:} Guest; Customer; hệ thống xác thực (NestJS)

\textbf{Artifact thiết kế chịu ảnh hưởng:} Use Case auth, Security Design, API Auth, UI Web

\textbf{FR-AUTH-001 - Đăng ký Customer}

\begin{longtable}[]{@{}
  >{\raggedright\arraybackslash}p{(\columnwidth - 2\tabcolsep) * \real{0.2500}}
  >{\raggedright\arraybackslash}p{(\columnwidth - 2\tabcolsep) * \real{0.7500}}@{}}
\toprule\noalign{}
\begin{minipage}[b]{\linewidth}\raggedright
\textbf{Thuộc tính}
\end{minipage} & \begin{minipage}[b]{\linewidth}\raggedright
\textbf{Nội dung}
\end{minipage} \\
\midrule\noalign{}
\endhead
\bottomrule\noalign{}
\endlastfoot
\textbf{Mức ưu tiên} & MUST \\
\textbf{Trạng thái} & CONFIRMED \\
\textbf{Yêu cầu} & Hệ thống phải cho phép Guest tạo tài khoản Customer bằng email và password có thể kèm họ tên. \\
\textbf{Lý do/Nguồn} & Charter §3.1 (Identity: đăng ký); Task 7. \\
\textbf{Actor chính} & Guest; hệ thống xác thực \\
\end{longtable}

\textbf{Phân tích nghiệp vụ}

Đăng ký là điểm tạo danh tính nội bộ đầu tiên và là điều kiện để tạo đơn hàng. BrewLite chỉ có một vai trò Customer nên request đăng ký không được mang theo bất kỳ thông tin phân quyền nào.

\textbf{Điều kiện dữ liệu và ràng buộc liên quan}

\begin{itemize}
\item
  Input tối thiểu: email, password.
\item
  Email được chuẩn hóa (trim, lowercase) và phải duy nhất.
\item
  Password tối thiểu 8 ký tự; lưu dưới dạng hash bcrypt, không bao giờ lưu hoặc trả về password gốc.
\item
  Điểm tích lũy khởi tạo bằng 0.
\end{itemize}

\textbf{Vấn đề có thể gặp và phân tích tác động}

\begin{longtable}[]{@{}
  >{\raggedright\arraybackslash}p{(\columnwidth - 2\tabcolsep) * \real{0.2630}}
  >{\raggedright\arraybackslash}p{(\columnwidth - 2\tabcolsep) * \real{0.7370}}@{}}
\toprule\noalign{}
\begin{minipage}[b]{\linewidth}\raggedright
\textbf{Vấn đề}
\end{minipage} & \begin{minipage}[b]{\linewidth}\raggedright
\textbf{Tác động/Nguyên nhân}
\end{minipage} \\
\midrule\noalign{}
\endhead
\bottomrule\noalign{}
\endlastfoot
Trùng email & Hai tài khoản cho một người, lịch sử đơn và điểm bị tách. \\
Enumeration account(1) & Thông báo quá chi tiết lộ email nào đã tồn tại. \\
Mass assignment(2) & Client gửi thêm trường như loyaltyPoints hoặc role mà backend nhận mà không lọc. \\
Password yếu hoặc lộ qua log & Chiếm tài khoản. \\
\end{longtable}

(1) Enumeration: thử nhiều email hoặc tài khoản để suy ra cái nào tồn tại dựa trên sự khác biệt trong phản hồi.\\
(2) Mass assignment: backend gán trực tiếp mọi trường từ payload của client vào dữ liệu nội bộ.

\textbf{Hướng xử lý/giải pháp phân tích đề xuất}

\begin{itemize}
\item
  Ràng buộc unique ở tầng DB và kiểm tra ở server.
\item
  DTO whitelist: chỉ nhận email, password (và họ tên); bỏ qua mọi trường khác.
\item
  Hash bằng bcrypt; không log password.
\item
  Thông báo lỗi trùng email đủ rõ để dùng được, không trả thông tin nhạy cảm khác.
\item
  Với đăng ký, nhóm chấp nhận để lộ việc email đã tồn tại (trả 409) vì đây là MVP và phản hồi rõ ràng quan trọng hơn; riêng đăng nhập vẫn dùng thông báo lỗi chung để chống dò tài khoản.
\end{itemize}

\textbf{Chuyển giao sang thiết kế}

\begin{itemize}
\item
  Use Case Đăng ký; ERD thực thể User (unique email); API POST /auth/register; màn hình đăng ký; test case xác thực dữ liệu.
\end{itemize}

\textbf{Tiêu chí kiểm chứng ở mức yêu cầu}

\begin{itemize}
\item
  Đăng ký hợp lệ tạo được tài khoản; password trong DB là hash.
\item
  Email trùng, email sai định dạng, password yếu bị từ chối với lỗi rõ ràng.
\item
  Gửi kèm loyaltyPoints hoặc role không làm thay đổi dữ liệu.
\end{itemize}

\textbf{FR-AUTH-002 - Đăng nhập bằng email/password}

\begin{longtable}[]{@{}
  >{\raggedright\arraybackslash}p{(\columnwidth - 2\tabcolsep) * \real{0.2500}}
  >{\raggedright\arraybackslash}p{(\columnwidth - 2\tabcolsep) * \real{0.7500}}@{}}
\toprule\noalign{}
\begin{minipage}[b]{\linewidth}\raggedright
\textbf{Thuộc tính}
\end{minipage} & \begin{minipage}[b]{\linewidth}\raggedright
\textbf{Nội dung}
\end{minipage} \\
\midrule\noalign{}
\endhead
\bottomrule\noalign{}
\endlastfoot
\textbf{Mức ưu tiên} & MUST \\
\textbf{Trạng thái} & CONFIRMED \\
\textbf{Yêu cầu} & Customer đăng nhập bằng thông tin hợp lệ và nhận JWT; thông tin sai không tạo phiên đăng nhập. Đăng nhập thành công không làm thay đổi giỏ hàng hiện có. \\
\textbf{Lý do/Nguồn} & Charter §3.1 (đăng nhập JWT); Task 7. \\
\textbf{Actor chính} & Customer; hệ thống xác thực \\
\end{longtable}

\textbf{Phân tích nghiệp vụ}

Đăng nhập xác thực credential rồi cấp access token. Đăng nhập thành công không có nghĩa là có quyền với mọi tài nguyên; quyền truy cập vẫn kiểm tra ở từng hành động (xem FR-AUTH-003).

\textbf{Điều kiện dữ liệu và ràng buộc liên quan}

\begin{itemize}
\item
  Identifier: email.
\item
  So sánh password bằng bcrypt với hash đã lưu.
\item
  Token chứa định danh người dùng (sub), có thời hạn hết hạn sau 60 phút.
\end{itemize}

\textbf{Vấn đề có thể gặp và phân tích tác động}

\begin{longtable}[]{@{}
  >{\raggedright\arraybackslash}p{(\columnwidth - 2\tabcolsep) * \real{0.2630}}
  >{\raggedright\arraybackslash}p{(\columnwidth - 2\tabcolsep) * \real{0.7370}}@{}}
\toprule\noalign{}
\begin{minipage}[b]{\linewidth}\raggedright
\textbf{Vấn đề}
\end{minipage} & \begin{minipage}[b]{\linewidth}\raggedright
\textbf{Tác động/Nguyên nhân}
\end{minipage} \\
\midrule\noalign{}
\endhead
\bottomrule\noalign{}
\endlastfoot
Thông báo khác nhau cho ``email không tồn tại'' và ``sai password'' & Tăng rủi ro enumeration. \\
Thử password liên tục & Dò mật khẩu. \\
JWT secret yếu hoặc commit lên repo & Giả mạo token. \\
\end{longtable}

\textbf{Hướng xử lý/giải pháp phân tích đề xuất}

\begin{itemize}
\item
  Dùng thông báo lỗi chung ``Email hoặc mật khẩu không đúng''.
\item
  JWT\_SECRET đặt đủ dài và ngẫu nhiên, lưu trong .env, không commit lên repo.
\item
  (SHOULD) giới hạn số lần thử đăng nhập nếu thời gian cho phép.
\end{itemize}

\textbf{Chuyển giao sang thiết kế}

\begin{itemize}
\item
  Use Case Đăng nhập; API POST /auth/login.
\item
  Security Design: Security Design (cấu trúc và thời hạn JWT), xác thực bằng Passport JWT, so khớp mật khẩu bằng bcrypt.
\item
  Cấu hình: cung cấp .env.example không chứa secret.
\item
  NFR bảo mật.
\end{itemize}

\textbf{Tiêu chí kiểm chứng ở mức yêu cầu}

\begin{itemize}
\item
  Hai trường hợp ``email không tồn tại'' và ``sai password'' cho phản hồi giống nhau.
\item
  Credential đúng nhận được token hợp lệ; credential sai trả lỗi chung và không có token.
\end{itemize}

\textbf{FR-AUTH-003 - Xác thực JWT và bảo vệ API}

\begin{longtable}[]{@{}
  >{\raggedright\arraybackslash}p{(\columnwidth - 2\tabcolsep) * \real{0.2500}}
  >{\raggedright\arraybackslash}p{(\columnwidth - 2\tabcolsep) * \real{0.7500}}@{}}
\toprule\noalign{}
\begin{minipage}[b]{\linewidth}\raggedright
\textbf{Thuộc tính}
\end{minipage} & \begin{minipage}[b]{\linewidth}\raggedright
\textbf{Nội dung}
\end{minipage} \\
\midrule\noalign{}
\endhead
\bottomrule\noalign{}
\endlastfoot
\textbf{Mức ưu tiên} & MUST \\
\textbf{Trạng thái} & CONFIRMED \\
\textbf{Yêu cầu} & Các API nghiệp vụ cần đăng nhập chỉ được truy cập khi có JWT hợp lệ. \\
\textbf{Lý do/Nguồn} & Charter §3.1 (đăng nhập JWT); Task 7; Charter §7 (bảo mật cơ bản). \\
\textbf{Actor chính} & Customer; hệ thống xác thực \\
\end{longtable}

\textbf{Phân tích nghiệp vụ}

Guest được xem menu, nhưng tạo đơn, thanh toán và xem lịch sử đơn yêu cầu xác thực. Danh tính phải lấy từ token đã xác minh, nếu không Customer có thể truy cập dữ liệu của người khác (quyền sở hữu tài nguyên xem FR-AUTHZ-002).

\textbf{Điều kiện dữ liệu và ràng buộc liên quan}

\begin{itemize}
\item
  Public: GET /products, chi tiết sản phẩm, POST /auth/register, POST /auth/login.
\item
  Cần JWT: POST /orders, GET /orders/me, POST /payments.
\item
  Token thiếu, sai chữ ký hoặc hết hạn thì bị từ chối (mã lỗi xem bảng ở FR-AUTHZ-001).
\end{itemize}

\textbf{Vấn đề có thể gặp và phân tích tác động}

\begin{longtable}[]{@{}
  >{\raggedright\arraybackslash}p{(\columnwidth - 2\tabcolsep) * \real{0.2630}}
  >{\raggedright\arraybackslash}p{(\columnwidth - 2\tabcolsep) * \real{0.7370}}@{}}
\toprule\noalign{}
\begin{minipage}[b]{\linewidth}\raggedright
\textbf{Vấn đề}
\end{minipage} & \begin{minipage}[b]{\linewidth}\raggedright
\textbf{Tác động/Nguyên nhân}
\end{minipage} \\
\midrule\noalign{}
\endhead
\bottomrule\noalign{}
\endlastfoot
Lấy userId từ body hoặc query & Tạo hoặc xem đơn thay người khác (IDOR). \\
Quên áp Guard cho một endpoint & Lộ dữ liệu hoặc cho thao tác không xác thực. \\
Token hết hạn nhưng FE không xử lý & Người dùng bị lỗi khó hiểu giữa lúc thanh toán. \\
\end{longtable}

\textbf{Hướng xử lý/giải pháp phân tích đề xuất}

\begin{itemize}
\item
  Dùng JWT Guard của NestJS; lấy userId từ token.
\item
  FE chuyển hướng về trang đăng nhập khi nhận 401 (token hết hạn hoặc không hợp lệ). Trường hợp này không phải đăng xuất chủ động nên giỏ hàng được giữ, để người dùng đăng nhập lại và tiếp tục đặt món.
\end{itemize}

\textbf{Chuyển giao sang thiết kế}

\begin{itemize}
\item
  Security Design; bảng phân quyền endpoint (mục 8); API Auth; xử lý lỗi 401 ở FE; test truy cập trái phép.
\end{itemize}

\textbf{Tiêu chí kiểm chứng ở mức yêu cầu}

\begin{itemize}
\item
  Gọi endpoint cần xác thực không có token, token sai hoặc hết hạn đều trả 401.
\end{itemize}

\textbf{FR-AUTH-004 - Đăng xuất}

\begin{longtable}[]{@{}
  >{\raggedright\arraybackslash}p{(\columnwidth - 2\tabcolsep) * \real{0.2500}}
  >{\raggedright\arraybackslash}p{(\columnwidth - 2\tabcolsep) * \real{0.7500}}@{}}
\toprule\noalign{}
\begin{minipage}[b]{\linewidth}\raggedright
\textbf{Thuộc tính}
\end{minipage} & \begin{minipage}[b]{\linewidth}\raggedright
\textbf{Nội dung}
\end{minipage} \\
\midrule\noalign{}
\endhead
\bottomrule\noalign{}
\endlastfoot
\textbf{Mức ưu tiên} & SHOULD \\
\textbf{Trạng thái} & DERIVED \\
\textbf{Yêu cầu} & Customer có thể đăng xuất khỏi thiết bị hiện tại. Khi đăng xuất, client xóa token \textbf{và xóa toàn bộ giỏ hàng}, rồi trở về trạng thái Guest. \\
\textbf{Lý do/Nguồn} & Cần cho trải nghiệm tài khoản thực tế; Charter chỉ nêu đăng nhập. Xóa giỏ để tránh lộ lựa chọn món trên thiết bị dùng chung. \\
\textbf{Actor chính} & Customer \\
\end{longtable}

\textbf{Phân tích nghiệp vụ}

BrewLite chỉ dùng access token ngắn hạn, không có refresh session phía server, nên đăng xuất là xóa token và dữ liệu phiên ở client. Token cũ vẫn dùng được tới khi hết hạn; nhóm chấp nhận giới hạn này cho MVP. Giỏ hàng nằm ở client (Zustand) nên cũng được dọn trong cùng thao tác đăng xuất.

\textbf{Điều kiện dữ liệu và ràng buộc liên quan}

\begin{itemize}
\item
  Đăng xuất là thao tác chủ động của người dùng. Token hết hạn hoặc nhận 401 không được coi là đăng xuất (xem FR-AUTH-003).
\item
  Giỏ hàng bị xóa kể cả khi giỏ còn món chưa đặt; nên có hộp xác nhận nếu giỏ không rỗng (SHOULD).
\item
  Vòng đời giỏ xem FR-CART-004.
\end{itemize}

\begin{longtable}[]{@{}
  >{\raggedright\arraybackslash}p{(\columnwidth - 4\tabcolsep) * \real{0.5175}}
  >{\raggedright\arraybackslash}p{(\columnwidth - 4\tabcolsep) * \real{0.2402}}
  >{\raggedright\arraybackslash}p{(\columnwidth - 4\tabcolsep) * \real{0.2423}}@{}}
\toprule\noalign{}
\endhead
\bottomrule\noalign{}
\endlastfoot
\begin{minipage}[t]{\linewidth}\raggedright
\begin{quote}
Tình huống
\end{quote}
\end{minipage} & \begin{minipage}[t]{\linewidth}\raggedright
\begin{quote}
Token
\end{quote}
\end{minipage} & \begin{minipage}[t]{\linewidth}\raggedright
\begin{quote}
Giỏ hàng
\end{quote}
\end{minipage} \\
\begin{minipage}[t]{\linewidth}\raggedright
\begin{quote}
Đăng xuất chủ động
\end{quote}
\end{minipage} & \begin{minipage}[t]{\linewidth}\raggedright
\begin{quote}
Xóa
\end{quote}
\end{minipage} & \begin{minipage}[t]{\linewidth}\raggedright
\begin{quote}
Xóa
\end{quote}
\end{minipage} \\
\begin{minipage}[t]{\linewidth}\raggedright
\begin{quote}
Token hết hạn hoặc nhận 401
\end{quote}
\end{minipage} & \begin{minipage}[t]{\linewidth}\raggedright
\begin{quote}
Xóa
\end{quote}
\end{minipage} & \begin{minipage}[t]{\linewidth}\raggedright
\begin{quote}
Giữ
\end{quote}
\end{minipage} \\
\end{longtable}

\textbf{Vấn đề có thể gặp và phân tích tác động}

\begin{longtable}[]{@{}
  >{\raggedright\arraybackslash}p{(\columnwidth - 2\tabcolsep) * \real{0.2630}}
  >{\raggedright\arraybackslash}p{(\columnwidth - 2\tabcolsep) * \real{0.7370}}@{}}
\toprule\noalign{}
\begin{minipage}[b]{\linewidth}\raggedright
\textbf{Vấn đề}
\end{minipage} & \begin{minipage}[b]{\linewidth}\raggedright
\textbf{Tác động/Nguyên nhân}
\end{minipage} \\
\midrule\noalign{}
\endhead
\bottomrule\noalign{}
\endlastfoot
Token bị đánh cắp vẫn dùng được tới khi hết hạn & Không có thu hồi phía server. \\
Người dùng mất giỏ khi bấm nhầm đăng xuất & Phải chọn lại món. \\
Giỏ lưu ở localStorage mà không bị xóa & Người đăng nhập sau trên cùng máy thấy giỏ của người trước. \\
\end{longtable}

\textbf{Hướng xử lý/giải pháp phân tích đề xuất}

\begin{itemize}
\item
  Khi đăng xuất: xóa token, reset toàn bộ state giỏ hàng (kể cả bản persist ở localStorage nếu dùng), chuyển về trang menu.
\item
  Ghi rõ giới hạn ``không thu hồi token phía server'' trong tài liệu.
\end{itemize}

\textbf{Chuyển giao sang thiết kế}

\begin{itemize}
\item
  Auth State và Cart State phía FE; Use Case Đăng xuất.
\end{itemize}

\textbf{Tiêu chí kiểm chứng ở mức yêu cầu}

\begin{itemize}
\item
  Sau đăng xuất, thao tác cần xác thực yêu cầu đăng nhập lại.
\item
  Sau đăng xuất rồi đăng nhập tài khoản khác, giỏ hàng rỗng, kể cả sau khi tải lại trang.
\end{itemize}

\textbf{7.2. Authorization \& Ownership}

\textbf{Actor/phạm vi chính:} Guest; Customer; backend authorization layer

\textbf{Artifact thiết kế chịu ảnh hưởng:} Bảng phân quyền endpoint (mục 8), Security Design, Use Case, Test Strategy

\textbf{FR-AUTHZ-001 - Phân quyền truy cập endpoint và mặc định từ chối}

\begin{longtable}[]{@{}
  >{\raggedright\arraybackslash}p{(\columnwidth - 2\tabcolsep) * \real{0.2500}}
  >{\raggedright\arraybackslash}p{(\columnwidth - 2\tabcolsep) * \real{0.7500}}@{}}
\toprule\noalign{}
\begin{minipage}[b]{\linewidth}\raggedright
\textbf{Thuộc tính}
\end{minipage} & \begin{minipage}[b]{\linewidth}\raggedright
\textbf{Nội dung}
\end{minipage} \\
\midrule\noalign{}
\endhead
\bottomrule\noalign{}
\endlastfoot
\textbf{Mức ưu tiên} & MUST \\
\textbf{Trạng thái} & CONFIRMED \\
\textbf{Yêu cầu} & Hệ thống phân tách rõ endpoint công khai (Guest) và endpoint cần đăng nhập (Customer). Endpoint không được khai báo công khai thì mặc định yêu cầu xác thực; khi không xác định được quyền thì từ chối. \\
\textbf{Lý do/Nguồn} & Charter §3.1 (đăng nhập JWT), §7 (bảo mật cơ bản); nguyên tắc fail closed(1). \\
\textbf{Actor chính} & Guest; Customer; backend authorization layer \\
\end{longtable}

\textbf{Phân tích nghiệp vụ}

BrewLite chỉ có hai nhóm truy cập: Guest (chưa đăng nhập) và Customer (có JWT hợp lệ). Không có Role/Permission chi tiết, nên việc phân quyền là quyết định ``endpoint này có yêu cầu đăng nhập không''. Nếu để mặc định mở, endpoint mới thêm sẽ vô tình công khai.

\textbf{Điều kiện dữ liệu và ràng buộc liên quan}

\begin{itemize}
\item
  Công khai: xem danh sách và chi tiết sản phẩm, đăng ký, đăng nhập.
\item
  Cần JWT: tạo đơn, xem lịch sử đơn, thanh toán, các thao tác trên đơn của chính mình.
\item
  Endpoint nội bộ dùng để demo chuyển trạng thái đơn (PREPARING → READY → COMPLETED) không dùng chung quyền Customer; cách bảo vệ do nhóm chốt ở mục 11.
\item
  Bảng đầy đủ endpoint ↔ quyền nằm ở mục 8.
\end{itemize}

\begin{longtable}[]{@{}
  >{\raggedright\arraybackslash}p{(\columnwidth - 4\tabcolsep) * \real{0.3869}}
  >{\raggedright\arraybackslash}p{(\columnwidth - 4\tabcolsep) * \real{0.1313}}
  >{\raggedright\arraybackslash}p{(\columnwidth - 4\tabcolsep) * \real{0.4818}}@{}}
\toprule\noalign{}
\endhead
\bottomrule\noalign{}
\endlastfoot
\textbf{Tình huống} & \textbf{Mã lỗi} & \textbf{Ghi chú} \\
Không có token, token sai chữ ký hoặc hết hạn & 401 & Thông báo chung, không nêu lý do chi tiết. \\
Truy cập đơn không thuộc về mình hoặc không tồn tại & 404 & Hai trường hợp phải cho phản hồi giống hệt nhau (xem FR-AUTHZ-002). \\
Gọi endpoint nội bộ khi không có khóa & 403 & Chỉ áp dụng cho endpoint nội bộ đổi trạng thái đơn. \\
Đơn của mình nhưng trạng thái không cho phép hành động (ví dụ thanh toán đơn đã PAID) & 409 & Chỉ trả sau khi đã xác nhận đơn thuộc về người gọi; nếu là đơn người khác vẫn trả 404. \\
\end{longtable}

\textbf{Vấn đề có thể gặp và phân tích tác động}

\begin{longtable}[]{@{}
  >{\raggedright\arraybackslash}p{(\columnwidth - 2\tabcolsep) * \real{0.2500}}
  >{\raggedright\arraybackslash}p{(\columnwidth - 2\tabcolsep) * \real{0.7500}}@{}}
\toprule\noalign{}
\begin{minipage}[b]{\linewidth}\raggedright
\textbf{Vấn đề}
\end{minipage} & \begin{minipage}[b]{\linewidth}\raggedright
\textbf{Tác động/Nguyên nhân}
\end{minipage} \\
\midrule\noalign{}
\endhead
\bottomrule\noalign{}
\endlastfoot
Quên áp Guard cho endpoint mới & Lộ dữ liệu hoặc cho thao tác không xác thực. \\
Endpoint nội bộ đổi trạng thái đơn bị Customer gọi được & Customer tự chuyển đơn sang COMPLETED hoặc PAID. \\
Thông báo lỗi chứa chi tiết dữ liệu & Rò thông tin tài nguyên. \\
\end{longtable}

\textbf{Hướng xử lý/giải pháp phân tích đề xuất}

\begin{itemize}
\item
  Áp JWT Guard ở mức toàn ứng dụng, chỉ mở công khai các endpoint đã liệt kê.
\item
  Endpoint nội bộ có cơ chế bảo vệ riêng (ví dụ khóa nội bộ trong cấu hình), không công khai cho Customer.
\item
  Trả mã lỗi theo bảng trên; thông báo chung, không kèm dữ liệu tài nguyên.
\end{itemize}

\textbf{Chuyển giao sang thiết kế}

\begin{itemize}
\item
  Bảng phân quyền endpoint (mục 8); Security Design (NestJS Guard); danh mục endpoint; test truy cập trái phép.
\end{itemize}

\textbf{Tiêu chí kiểm chứng ở mức yêu cầu}

\begin{itemize}
\item
  Endpoint cần xác thực gọi không có token trả 401.
\item
  Endpoint mới thêm mà chưa khai báo công khai không truy cập được khi chưa đăng nhập.
\item
  Customer không gọi được endpoint nội bộ đổi trạng thái đơn.
\item
  Mỗi tình huống trong bảng trả đúng mã lỗi đã quy định.
\end{itemize}

\begin{quote}
(1) Fail closed: nếu hệ thống không thể xác minh đầy đủ quyền truy cập hoặc quyền sở hữu tài nguyên thì thao tác bị từ chối thay vì được cho phép mặc định.
\end{quote}

\textbf{FR-AUTHZ-002 - Kiểm tra quyền sở hữu đơn hàng (ownership)}

\begin{longtable}[]{@{}
  >{\raggedright\arraybackslash}p{(\columnwidth - 2\tabcolsep) * \real{0.2500}}
  >{\raggedright\arraybackslash}p{(\columnwidth - 2\tabcolsep) * \real{0.7500}}@{}}
\toprule\noalign{}
\begin{minipage}[b]{\linewidth}\raggedright
\textbf{Thuộc tính}
\end{minipage} & \begin{minipage}[b]{\linewidth}\raggedright
\textbf{Nội dung}
\end{minipage} \\
\midrule\noalign{}
\endhead
\bottomrule\noalign{}
\endlastfoot
\textbf{Mức ưu tiên} & MUST \\
\textbf{Trạng thái} & DERIVED \\
\textbf{Yêu cầu} & Customer chỉ được xem, thanh toán và hủy đơn thuộc về chính mình. Danh tính lấy từ JWT, không lấy từ dữ liệu client gửi. \\
\textbf{Lý do/Nguồn} & Phòng IDOR (mục 5); Charter §3.1 (lịch sử đơn); Task 9. \\
\textbf{Actor chính} & Customer; backend authorization layer \\
\end{longtable}

\textbf{Phân tích nghiệp vụ}

Có đăng nhập là chưa đủ: một Customer hợp lệ vẫn không được đụng đến đơn của người khác. Đây là rủi ro bảo mật đáng kể nhất của hệ thống một cửa hàng, vì mã đơn có thể đoán hoặc bị thay trên request.

\textbf{Điều kiện dữ liệu và ràng buộc liên quan}

\begin{itemize}
\item
  Mỗi đơn gắn với userId của người tạo.
\item
  Danh sách đơn của tôi được lọc theo userId trong token (/orders/me).
\item
  Với thao tác trên một đơn cụ thể (thanh toán, hủy, xem chi tiết), so sánh userId của đơn với userId trong token.
\item
  Không nhận userId từ body hoặc query để quyết định quyền.
\end{itemize}

\textbf{Vấn đề có thể gặp và phân tích tác động}

\begin{longtable}[]{@{}
  >{\raggedright\arraybackslash}p{(\columnwidth - 2\tabcolsep) * \real{0.2500}}
  >{\raggedright\arraybackslash}p{(\columnwidth - 2\tabcolsep) * \real{0.7500}}@{}}
\toprule\noalign{}
\begin{minipage}[b]{\linewidth}\raggedright
\textbf{Vấn đề}
\end{minipage} & \begin{minipage}[b]{\linewidth}\raggedright
\textbf{Tác động/Nguyên nhân}
\end{minipage} \\
\midrule\noalign{}
\endhead
\bottomrule\noalign{}
\endlastfoot
Đổi orderId trên request (IDOR) & Xem hoặc thanh toán đơn của người khác. \\
Tin userId do client gửi & Giả danh người khác khi tạo đơn. \\
Phản hồi khác nhau cho ``đơn không tồn tại'' và ``đơn của người khác'' & Cho phép dò mã đơn hợp lệ. \\
\end{longtable}

\textbf{Hướng xử lý/giải pháp phân tích đề xuất}

\begin{itemize}
\item
  Truy vấn tài nguyên kèm userId từ token; không tìm thấy hoặc không thuộc sở hữu thì trả cùng một lỗi 404 (bảng ở FR-AUTHZ-001).
\item
  Test âm tính bắt buộc: Customer A truy cập đơn của Customer B.
\item
  Mặc định từ chối khi không xác định được chủ sở hữu.
\end{itemize}

\textbf{Chuyển giao sang thiết kế}

\begin{itemize}
\item
  ERD (Order có khóa ngoại userId); thiết kế truy vấn; bảng phân quyền mục 8.
\end{itemize}

\textbf{Tiêu chí kiểm chứng ở mức yêu cầu}

\begin{itemize}
\item
  Customer A không đọc, thanh toán hoặc hủy được đơn của Customer B dù biết mã đơn; phản hồi không phân biệt `không tồn tại' và `không phải của bạn'.
\item
  GET /orders/me chỉ trả đơn của người dùng trong token.
\end{itemize}

\textbf{FR-AUTHZ-003 - Backend là nơi quyết định quyền truy cập cuối cùng}

\begin{longtable}[]{@{}
  >{\raggedright\arraybackslash}p{(\columnwidth - 2\tabcolsep) * \real{0.2500}}
  >{\raggedright\arraybackslash}p{(\columnwidth - 2\tabcolsep) * \real{0.7500}}@{}}
\toprule\noalign{}
\begin{minipage}[b]{\linewidth}\raggedright
\textbf{Thuộc tính}
\end{minipage} & \begin{minipage}[b]{\linewidth}\raggedright
\textbf{Nội dung}
\end{minipage} \\
\midrule\noalign{}
\endhead
\bottomrule\noalign{}
\endlastfoot
\textbf{Mức ưu tiên} & MUST \\
\textbf{Trạng thái} & DERIVED \\
\textbf{Yêu cầu} & Việc Web Client ẩn nút hoặc màn hình chỉ phục vụ UX, không phải cơ chế bảo mật. Backend kiểm tra theo thứ tự: xác thực → quyền sở hữu → quy tắc trạng thái nghiệp vụ → cho phép hoặc từ chối. \\
\textbf{Lý do/Nguồn} & Mục 3 (ranh giới hệ thống); Charter §7. \\
\textbf{Actor chính} & Customer; backend authorization layer \\
\end{longtable}

\textbf{Phân tích nghiệp vụ}

Client có thể bị sửa, replay hoặc bị gọi API trực tiếp bằng Postman. Quyền hợp lệ vẫn có thể không đủ nếu đơn đang ở trạng thái không cho phép (ví dụ thanh toán đơn đã PAID), nên kiểm tra trạng thái là lớp cuối cùng.

\textbf{Điều kiện dữ liệu và ràng buộc liên quan}

\begin{itemize}
\item
  Thứ tự kiểm tra: JWT hợp lệ → tài nguyên thuộc người dùng → trạng thái cho phép hành động → thực hiện.
\item
  Trạng thái chỉ đổi qua các hành động hợp lệ của state machine, không có thao tác gán trạng thái tùy ý.
\end{itemize}

\textbf{Vấn đề có thể gặp và phân tích tác động}

\begin{longtable}[]{@{}
  >{\raggedright\arraybackslash}p{(\columnwidth - 2\tabcolsep) * \real{0.2500}}
  >{\raggedright\arraybackslash}p{(\columnwidth - 2\tabcolsep) * \real{0.7500}}@{}}
\toprule\noalign{}
\begin{minipage}[b]{\linewidth}\raggedright
\textbf{Vấn đề}
\end{minipage} & \begin{minipage}[b]{\linewidth}\raggedright
\textbf{Tác động/Nguyên nhân}
\end{minipage} \\
\midrule\noalign{}
\endhead
\bottomrule\noalign{}
\endlastfoot
Chỉ ẩn nút ở FE & Gọi API trực tiếp vẫn thực hiện được. \\
Kiểm tra trạng thái trước kiểm tra sở hữu & Rò thông tin trạng thái của đơn người khác. \\
Cho gán trạng thái tùy ý & Phá vòng đời Order và Payment. \\
\end{longtable}

\textbf{Hướng xử lý/giải pháp phân tích đề xuất}

\begin{itemize}
\item
  Tập trung kiểm tra ở backend (Guard và service); FE chỉ xử lý hiển thị.
\item
  Kiểm tra sở hữu trước, rồi mới trả lỗi nghiệp vụ về trạng thái.
\item
  Chi tiết các chuyển trạng thái hợp lệ ở mục 11.
\end{itemize}

\textbf{Chuyển giao sang thiết kế}

\begin{itemize}
\item
  API Security; State Diagram; quy ước lỗi; test gọi API độc lập với UI.
\end{itemize}

\textbf{Tiêu chí kiểm chứng ở mức yêu cầu}

\begin{itemize}
\item
  Gọi API trực tiếp (không qua UI) mà không đủ điều kiện vẫn bị từ chối.
\item
  Thanh toán đơn đã PAID hoặc đơn của người khác đều bị từ chối, và thông báo lỗi không lộ trạng thái đơn của người khác.
\end{itemize}

\textbf{7.3. Catalog}

\textbf{Actor/phạm vi chính:} Guest; Customer; Backend API (module Product)

\textbf{Artifact thiết kế chịu ảnh hưởng:} Use Case xem menu, ERD (Product, Topping), API Product, UI trang menu và chi tiết, Test Strategy

\textbf{FR-CAT-001 - Xem danh sách đồ uống (menu)}

\begin{longtable}[]{@{}
  >{\raggedright\arraybackslash}p{(\columnwidth - 2\tabcolsep) * \real{0.2500}}
  >{\raggedright\arraybackslash}p{(\columnwidth - 2\tabcolsep) * \real{0.7500}}@{}}
\toprule\noalign{}
\begin{minipage}[b]{\linewidth}\raggedright
\textbf{Thuộc tính}
\end{minipage} & \begin{minipage}[b]{\linewidth}\raggedright
\textbf{Nội dung}
\end{minipage} \\
\midrule\noalign{}
\endhead
\bottomrule\noalign{}
\endlastfoot
\textbf{Mức ưu tiên} & MUST (danh sách); SHOULD (phân trang) \\
\textbf{Trạng thái} & CONFIRMED (danh sách); DERIVED (phân trang) \\
\textbf{Yêu cầu} & Hệ thống phải cung cấp danh sách đồ uống gồm ảnh, tên, giá cơ bản (cũng là giá size S) và trạng thái còn/hết hàng cho Guest và Customer, không yêu cầu đăng nhập. Hệ thống \textbf{nên} hỗ trợ phân trang tùy chọn để menu lớn vẫn phản hồi nhanh. \\
\textbf{Lý do/Nguồn} & Charter §3.1 (Menu: hiển thị ảnh, tên, giá; API danh sách); Task 2, 3. Phân trang là quyết định của nhóm, áp dụng khi số món vượt ngưỡng, không thuộc yêu cầu của đề. \\
\textbf{Actor chính} & Guest; Customer; Backend API \\
\end{longtable}

\textbf{Phân tích nghiệp vụ}

Menu là điểm vào của toàn bộ luồng đặt món nên phải xem được khi chưa đăng nhập (FR-AUTHZ-001). Danh sách chỉ gồm thông tin cần để hiển thị; size và topping chỉ có ở trang chi tiết. Giá trên menu là giá cơ bản, không phải giá cuối. Khi số món lớn, danh sách có thể chia theo trang để thời gian phản hồi và dung lượng ảnh không tăng theo.

\textbf{Điều kiện, dữ liệu và ràng buộc liên quan}

\begin{itemize}
\item
  Mỗi sản phẩm có: mã, tên, giá cơ bản, đường dẫn ảnh, trạng thái còn/hết hàng.
\item
  Giá lưu bằng số nguyên hoặc Decimal (VND), không dùng số thực.
\item
  Endpoint công khai: GET /products; tham số tùy chọn page, limit.
\item
  Thứ tự sắp xếp cố định (tên rồi mã sản phẩm), để dù có phân trang hay không thì thứ tự vẫn nhất quán.
\item
  \textbf{Điểm quyết định:} nhóm chốt cách hiển thị menu tại Sprint 1, trước khi làm Task 2 và 3. Nếu số món trong dữ liệu mẫu vượt \textbf{20}, frontend gọi API có phân trang (mặc định limit = 12); nếu không, frontend gọi không tham số.
\item
  Không truyền tham số: trả toàn bộ menu. Có truyền: page từ 1, limit từ 1 đến 50; giá trị sai trả 400.
\item
  Cấu trúc response ổn định ở cả hai chế độ: danh sách món kèm total, page, limit, totalPages. Không phân trang thì page = 1, limit = total, totalPages = 1.
\item
  Nếu chỉ truyền một trong hai tham số, tham số còn lại lấy mặc định (page = 1, limit = 12).
\item
  Giới hạn limit từ 1 đến 50 chỉ áp dụng cho request; khi không phân trang, limit trong phản hồi bằng total.
\end{itemize}

\textbf{Vấn đề có thể gặp và phân tích tác động}

\begin{longtable}[]{@{}
  >{\raggedright\arraybackslash}p{(\columnwidth - 2\tabcolsep) * \real{0.2500}}
  >{\raggedright\arraybackslash}p{(\columnwidth - 2\tabcolsep) * \real{0.7500}}@{}}
\toprule\noalign{}
\begin{minipage}[b]{\linewidth}\raggedright
\textbf{Vấn đề}
\end{minipage} & \begin{minipage}[b]{\linewidth}\raggedright
\textbf{Tác động/Nguyên nhân}
\end{minipage} \\
\midrule\noalign{}
\endhead
\bottomrule\noalign{}
\endlastfoot
Danh sách rỗng (chưa nạp dữ liệu) & Trang menu trống, demo thất bại. \\
Ảnh lỗi hoặc thiếu đường dẫn & Giao diện vỡ. \\
Thứ tự không ổn định giữa các lần gọi & Món bị trùng hoặc sót khi chuyển trang. \\
page hoặc limit sai (âm, chữ, quá lớn) & Lỗi không kiểm soát hoặc truy vấn nặng. \\
Trang vượt quá số trang hiện có & Màn hình lỗi dù dữ liệu hợp lệ. \\
Trả thừa dữ liệu nội bộ (số lượng tồn chính xác, trường nội bộ) & Lộ thông tin không cần thiết cho khách. \\
Menu có nhiều món nhưng không phân trang & Phản hồi chậm, vượt mục tiêu 500 ms, tải nhiều ảnh cùng lúc. \\
Bật phân trang sau mà response đổi cấu trúc & Frontend phải sửa lại, dễ lệch hợp đồng API giữa FE và BE. \\
\end{longtable}

\textbf{Hướng xử lý/giải pháp phân tích đề xuất}.

\begin{itemize}
\item
  Validate page và limit bằng DTO; sai thì trả 400 với thông báo rõ ràng.
\item
  Hỗ trợ cả hai chế độ ngay từ đầu trong backend: có tham số thì phân trang, không có thì trả hết; cấu trúc response không đổi.
\item
  Sắp xếp luôn có khóa phụ (mã sản phẩm) để thứ tự ổn định.
\item
  Trang vượt quá số trang trả 200 và danh sách rỗng kèm total, totalPages; UI hiển thị ``không có món ở trang này'' và nút về trang 1.
\item
  Trả danh sách qua DTO chỉ gồm trường hiển thị, không trả trực tiếp bản ghi database.
\item
  Frontend có trạng thái loading, empty, error (Techstack §5); ảnh lỗi dùng ảnh thay thế.
\item
  Nạp dữ liệu mẫu bằng seed (FR-CAT-005).
\end{itemize}

\textbf{Chuyển giao sang thiết kế}

\begin{itemize}
\item
  Use Case Xem menu; ERD Product; API GET /products (query page, limit, response kèm thông tin phân trang); UI trang menu (phân trang, loading/empty/error); test API danh sách và phân trang.
\end{itemize}

\textbf{Tiêu chí kiểm chứng ở mức yêu cầu}

\begin{itemize}
\item
  Gọi GET /products không tham số trả toàn bộ menu, đúng cấu trúc, Guest không cần token.
\item
  Mỗi phần tử có đủ ảnh, tên, giá cơ bản, trạng thái còn/hết hàng.
\item
  (SHOULD) Có page và limit: kết quả theo trang; duyệt hết các trang không trùng, không sót, tổng bằng total.
\item
  (SHOULD) page = 0, limit = 0 hoặc limit = 100 trả 400; page vượt quá totalPages trả 200 với danh sách rỗng.
\item
  Thời gian phản hồi với dữ liệu mẫu \textless{} 500 ms ở cả hai chế độ.
\end{itemize}

\textbf{FR-CAT-002 - Xem chi tiết sản phẩm và tùy chọn}

\begin{longtable}[]{@{}
  >{\raggedright\arraybackslash}p{(\columnwidth - 2\tabcolsep) * \real{0.2500}}
  >{\raggedright\arraybackslash}p{(\columnwidth - 2\tabcolsep) * \real{0.7500}}@{}}
\toprule\noalign{}
\begin{minipage}[b]{\linewidth}\raggedright
\textbf{Thuộc tính}
\end{minipage} & \begin{minipage}[b]{\linewidth}\raggedright
\textbf{Nội dung}
\end{minipage} \\
\midrule\noalign{}
\endhead
\bottomrule\noalign{}
\endlastfoot
\textbf{Mức ưu tiên} & MUST \\
\textbf{Trạng thái} & CONFIRMED \\
\textbf{Yêu cầu} & Hệ thống phải cung cấp chi tiết một sản phẩm gồm thông tin hiển thị, ba size S/M/L kèm hệ số giá, và danh sách topping kèm giá. Size S được chọn mặc định. \\
\textbf{Lý do/Nguồn} & Charter §3.1 (Chi tiết sản phẩm: size S/M/L, topping); Task 2, 4. \\
\textbf{Actor chính} & Guest; Customer; Backend API \\
\end{longtable}

\textbf{Phân tích nghiệp vụ}

Trang chi tiết là nơi khách chọn size, topping và thấy giá thay đổi theo lựa chọn. Backend trả đủ dữ liệu giá (hệ số size, giá topping) để frontend hiển thị giá tạm tính, nhưng chỉ backend quyết định giá cuối (FR-CAT-003).

\textbf{Điều kiện, dữ liệu và ràng buộc liên quan}

\begin{itemize}
\item
  Endpoint công khai: GET /products/:id.
\item
  Size hợp lệ: S (100\%), M (120\%), L (140\%); mặc định S.
\item
  Topping dùng chung cho toàn menu; mỗi topping có mã, tên và giá cố định. Tối đa 3 topping mỗi ly, mỗi loại chọn tối đa 1 lần.
\item
  Mã sản phẩm không tồn tại thì trả 404.
\end{itemize}

\textbf{Vấn đề có thể gặp và phân tích tác động}

\begin{longtable}[]{@{}
  >{\raggedright\arraybackslash}p{(\columnwidth - 2\tabcolsep) * \real{0.2500}}
  >{\raggedright\arraybackslash}p{(\columnwidth - 2\tabcolsep) * \real{0.7500}}@{}}
\toprule\noalign{}
\begin{minipage}[b]{\linewidth}\raggedright
\textbf{Vấn đề}
\end{minipage} & \begin{minipage}[b]{\linewidth}\raggedright
\textbf{Tác động/Nguyên nhân}
\end{minipage} \\
\midrule\noalign{}
\endhead
\bottomrule\noalign{}
\endlastfoot
Mã sản phẩm sai hoặc không tồn tại & Trang chi tiết lỗi, khách không biết xử lý. \\
Hệ số size hoặc giá topping lệch giữa FE và BE & Giá hiển thị khác giá tính khi tạo đơn. \\
Cho chọn quá 3 topping hoặc chọn trùng topping & Dữ liệu đơn không hợp lệ. \\
\end{longtable}

\textbf{Hướng xử lý/giải pháp phân tích đề xuất}

\begin{itemize}
\item
  Trả hệ số size và bảng topping trong cùng response chi tiết để FE không tự hard-code giá.
\item
  Mã không tồn tại trả 404 với thông báo chung; FE hiển thị trang ``không tìm thấy món''.
\item
  Backend kiểm tra lại size, số topping và trùng lặp khi tạo đơn.
\end{itemize}

\textbf{Chuyển giao sang thiết kế}

\begin{itemize}
\item
  Use Case Xem chi tiết; ERD Product, Topping; API GET /products/:id; UI trang chi tiết (size mặc định S); test API chi tiết.
\end{itemize}

\textbf{Tiêu chí kiểm chứng ở mức yêu cầu}

\begin{itemize}
\item
  Chi tiết trả đủ ba size với hệ số và danh sách topping kèm giá.
\item
  Mở trang chi tiết thì size S được chọn sẵn.
\item
  Mã không tồn tại trả 404 và UI xử lý được.
\item
  Guest xem được chi tiết mà không cần đăng nhập.
\end{itemize}

\textbf{FR-CAT-003 - Tính giá theo size và topping}

\begin{longtable}[]{@{}
  >{\raggedright\arraybackslash}p{(\columnwidth - 2\tabcolsep) * \real{0.2500}}
  >{\raggedright\arraybackslash}p{(\columnwidth - 2\tabcolsep) * \real{0.7500}}@{}}
\toprule\noalign{}
\begin{minipage}[b]{\linewidth}\raggedright
\textbf{Thuộc tính}
\end{minipage} & \begin{minipage}[b]{\linewidth}\raggedright
\textbf{Nội dung}
\end{minipage} \\
\midrule\noalign{}
\endhead
\bottomrule\noalign{}
\endlastfoot
\textbf{Mức ưu tiên} & MUST \\
\textbf{Trạng thái} & CONFIRMED \\
\textbf{Yêu cầu} & Giá của một dòng hàng được tính theo công thức: \textbf{Đơn giá = làm tròn(giá cơ bản × hệ số size ÷ 100) + tổng giá topping; Thành tiền dòng = Đơn giá × số lượng.} Frontend hiển thị giá tạm tính; backend tính lại và quyết định giá cuối khi tạo đơn. \\
\textbf{Lý do/Nguồn} & Charter §3.1 (tính giá theo tùy chọn); Task 4; Techstack §10 (unit test giá size/topping). Cách tính bằng hệ số là quyết định của nhóm. \\
\textbf{Actor chính} & Customer; Guest (xem giá tạm tính); Backend API \\
\end{longtable}

\textbf{Phân tích nghiệp vụ}

Đây là quy tắc tính tiền cốt lõi và là đầu vào cho giỏ hàng, đơn hàng và thanh toán. Hệ số size chỉ nhân với giá cơ bản của món; topping là số tiền cố định và không bị nhân. Nếu client tự quyết định giá, khách có thể sửa request để mua rẻ hơn, nên mọi giá trị tiền trong đơn đều do backend tính lại từ catalog.

\textbf{Điều kiện, dữ liệu và ràng buộc liên quan}

\begin{itemize}
\item
  Hệ số size: S = 100, M = 120, L = 140 (số nguyên phần trăm), cấu hình cố định ở backend.
\item
  Làm tròn đến đồng gần nhất (half up), thực hiện trên đơn giá một ly trước khi nhân số lượng.
\item
  Giá cơ bản trong dữ liệu mẫu là bội số của 5.000đ.
\item
  Số lượng nguyên từ 1 đến 10 mỗi dòng; tối đa 3 topping, không chọn trùng.
\item
  Tiền lưu bằng số nguyên/Decimal (VND), không dùng số thực dấu phẩy động.
\item
  Tổng đơn và giảm giá thuộc FR-ORD-001 và FR-PRO.
\end{itemize}

\textbf{Vấn đề có thể gặp và phân tích tác động}

\begin{longtable}[]{@{}
  >{\raggedright\arraybackslash}p{(\columnwidth - 2\tabcolsep) * \real{0.2500}}
  >{\raggedright\arraybackslash}p{(\columnwidth - 2\tabcolsep) * \real{0.7500}}@{}}
\toprule\noalign{}
\begin{minipage}[b]{\linewidth}\raggedright
\textbf{Vấn đề}
\end{minipage} & \begin{minipage}[b]{\linewidth}\raggedright
\textbf{Tác động/Nguyên nhân}
\end{minipage} \\
\midrule\noalign{}
\endhead
\bottomrule\noalign{}
\endlastfoot
Client gửi giá hoặc tổng tiền tự tính & Thất thoát doanh thu, giá sai. \\
Dùng số thực cho hệ số hoặc tiền & Sai số làm tròn, tổng lệch. \\
Công thức khác nhau ở FE và BE & Tổng hiển thị khác tổng thanh toán. \\
Nhân hệ số lên cả topping & Giá cao hơn quy định. \\
Giá catalog đổi giữa lúc khách xem và lúc đặt & Khách thấy giá cũ nhưng hệ thống tính giá mới. \\
Size, topping hoặc số lượng ngoài quy định & Giá sai hoặc đơn không hợp lệ. \\
\end{longtable}

\textbf{Hướng xử lý/giải pháp phân tích đề xuất}

\begin{itemize}
\item
  Backend luôn dùng giá catalog tại thời điểm tạo đơn; giá cuối được trả trong phản hồi tạo đơn và UI hiển thị lại trước bước thanh toán.
\item
  Backend tính toàn bộ giá từ catalog; bỏ qua mọi trường giá do client gửi.
\item
  Một hàm tính giá dùng chung ở backend, có unit test; FE dùng cùng công thức chỉ để hiển thị tạm tính.
\item
  Sau khi tạo đơn thì giá được chốt trong snapshot (FR-ORD-002).
\item
  Từ chối request ngoài quy định (size, topping, số lượng) với mã lỗi 400.
\end{itemize}

\textbf{Chuyển giao sang thiết kế}

\begin{itemize}
\item
  ERD (giá cơ bản, bảng Topping, OrderItem lưu size, đơn giá một ly, topping đã chọn kèm giá tại thời điểm đặt); module tính giá ở backend; Sequence tạo đơn; unit test giá; UI giá tạm tính.
\end{itemize}

\textbf{Tiêu chí kiểm chứng ở mức yêu cầu}

\begin{itemize}
\item
  Món giá cơ bản 45.000đ, size S, không topping, số lượng 1 → thành tiền 45.000đ.
\item
  Món 45.000đ, size L → đơn giá 63.000đ.
\item
  Món 45.000đ, size M, hai topping 5.000đ và 7.000đ, số lượng 2 → đơn giá 66.000đ (54.000 + 12.000), thành tiền 132.000đ.
\item
  Gửi giá thấp hơn trong request không làm giảm giá cuối.
\item
  Size lạ, quá 3 topping, topping trùng hoặc số lượng ngoài 1--10 bị từ chối với lỗi 400.
\item
  Tổng hiển thị ở UI khớp tổng backend tính khi tạo đơn.
\end{itemize}

\textbf{FR-CAT-004 - Hiển thị trạng thái còn/hết hàng}

\begin{longtable}[]{@{}
  >{\raggedright\arraybackslash}p{(\columnwidth - 2\tabcolsep) * \real{0.2500}}
  >{\raggedright\arraybackslash}p{(\columnwidth - 2\tabcolsep) * \real{0.7500}}@{}}
\toprule\noalign{}
\begin{minipage}[b]{\linewidth}\raggedright
\textbf{Thuộc tính}
\end{minipage} & \begin{minipage}[b]{\linewidth}\raggedright
\textbf{Nội dung}
\end{minipage} \\
\midrule\noalign{}
\endhead
\bottomrule\noalign{}
\endlastfoot
\textbf{Mức ưu tiên} & SHOULD \\
\textbf{Trạng thái} & DERIVED \\
\textbf{Yêu cầu} & Menu và trang chi tiết nên hiển thị sản phẩm còn hay hết hàng dựa trên tồn kho từ server; sản phẩm hết hàng không cho thêm vào giỏ. \\
\textbf{Lý do/Nguồn} & Techstack §6 (Product có stock); Charter §3.1 (tồn kho phục vụ Task 10). Charter không yêu cầu hiển thị trên menu. \\
\textbf{Actor chính} & Guest; Customer; Backend API \\
\end{longtable}

\textbf{Phân tích nghiệp vụ}

Hiển thị trạng thái hàng giúp khách tránh chọn món đã hết. Đây chỉ là thông tin tham khảo tại thời điểm xem: việc \textbf{chặn bán vượt tồn kho} vẫn do backend thực hiện khi tạo đơn (FR-INV-002), vì nhiều người có thể đặt cùng lúc.

\textbf{Điều kiện, dữ liệu và ràng buộc liên quan}

\begin{itemize}
\item
  Trạng thái ``còn hàng'' khi tồn kho \textgreater{} 0, ``hết hàng'' khi tồn kho = 0.
\item
  Chỉ trả trạng thái còn/hết, không bắt buộc hiển thị số lượng tồn chính xác.
\item
  Tồn kho được quản lý theo sản phẩm, không theo size hoặc topping.
\end{itemize}

\textbf{Vấn đề có thể gặp và phân tích tác động}

\begin{longtable}[]{@{}
  >{\raggedright\arraybackslash}p{(\columnwidth - 2\tabcolsep) * \real{0.2500}}
  >{\raggedright\arraybackslash}p{(\columnwidth - 2\tabcolsep) * \real{0.7500}}@{}}
\toprule\noalign{}
\begin{minipage}[b]{\linewidth}\raggedright
\textbf{Vấn đề}
\end{minipage} & \begin{minipage}[b]{\linewidth}\raggedright
\textbf{Tác động/Nguyên nhân}
\end{minipage} \\
\midrule\noalign{}
\endhead
\bottomrule\noalign{}
\endlastfoot
Menu hiển thị ``còn hàng'' nhưng khi đặt đã hết & Đơn bị từ chối dù khách tưởng còn. \\
Dữ liệu trạng thái bị cache cũ ở FE & Khách thấy thông tin lệch với server. \\
Chỉ dựa vào trạng thái ở FE để chặn mua & Gọi API trực tiếp vẫn đặt được món hết hàng. \\
\end{longtable}

\textbf{Hướng xử lý/giải pháp phân tích đề xuất}

\begin{itemize}
\item
  Trạng thái lấy từ server mỗi lần tải menu hoặc chi tiết.
\item
  Ghi rõ trong UI và tài liệu: trạng thái chỉ mang tính tham khảo; kết quả cuối cùng là kết quả kiểm tra khi tạo đơn.
\item
  Khi đặt mà hết hàng, backend trả lỗi rõ ràng (xem FR-INV-002) để UI báo khách.
\end{itemize}

\textbf{Chuyển giao sang thiết kế}

\begin{itemize}
\item
  API GET /products và /products/:id (trường trạng thái hàng); UI trạng thái hết hàng; liên kết với FR-INV-001, FR-INV-002.
\end{itemize}

\textbf{Tiêu chí kiểm chứng ở mức yêu cầu}

\begin{itemize}
\item
  Sản phẩm có tồn kho 0 hiển thị ``hết hàng'' và không thêm vào giỏ được qua UI.
\item
  Gọi trực tiếp API tạo đơn với món hết hàng vẫn bị từ chối (mã lỗi 409, xem FR-INV-002).
\end{itemize}

\textbf{FR-CAT-005 - Nạp dữ liệu catalog bằng dữ liệu mẫu (seed)}

\begin{longtable}[]{@{}
  >{\raggedright\arraybackslash}p{(\columnwidth - 2\tabcolsep) * \real{0.2500}}
  >{\raggedright\arraybackslash}p{(\columnwidth - 2\tabcolsep) * \real{0.7500}}@{}}
\toprule\noalign{}
\begin{minipage}[b]{\linewidth}\raggedright
\textbf{Thuộc tính}
\end{minipage} & \begin{minipage}[b]{\linewidth}\raggedright
\textbf{Nội dung}
\end{minipage} \\
\midrule\noalign{}
\endhead
\bottomrule\noalign{}
\endlastfoot
\textbf{Mức ưu tiên} & MUST \\
\textbf{Trạng thái} & DERIVED \\
\textbf{Yêu cầu} & Hệ thống phải có cơ chế nạp dữ liệu mẫu cho sản phẩm, topping và tồn kho (seed/migration) để chạy được demo và kiểm thử, vì MVP không có giao diện quản trị sản phẩm. \\
\textbf{Lý do/Nguồn} & Charter §3.2 (không có Seller/Admin), §6.1 (dữ liệu sản phẩm được mô phỏng); Techstack §8, §10. \\
\textbf{Actor chính} & Nhóm phát triển; Backend API \\
\end{longtable}

\textbf{Phân tích nghiệp vụ}

Vì không có kênh quản lý sản phẩm, toàn bộ dữ liệu catalog phải được nạp từ script khi khởi tạo. Seed cũng là điều kiện để các kiểm thử nghiệp vụ (đặt đồng thời, hết hàng, giá size/topping) chạy được với dữ liệu có thể lặp lại.

\textbf{Điều kiện, dữ liệu và ràng buộc liên quan}

\begin{itemize}
\item
  Dữ liệu mẫu đủ để phủ các tình huống: món bình thường, món hết hàng, món tồn kho thấp (để test oversell), nhiều topping.
\item
  Chạy seed nhiều lần không tạo dữ liệu trùng (có tính lặp lại an toàn).
\item
  Hướng dẫn chạy seed nằm trong README; môi trường Docker Compose tự có dữ liệu mẫu hoặc có lệnh nạp đơn giản.
\item
  Giá cơ bản của dữ liệu mẫu là bội số của 5.000đ; bộ topping dùng chung cho toàn menu.
\end{itemize}

\textbf{Vấn đề có thể gặp và phân tích tác động}

\begin{longtable}[]{@{}
  >{\raggedright\arraybackslash}p{(\columnwidth - 2\tabcolsep) * \real{0.2500}}
  >{\raggedright\arraybackslash}p{(\columnwidth - 2\tabcolsep) * \real{0.7500}}@{}}
\toprule\noalign{}
\begin{minipage}[b]{\linewidth}\raggedright
\textbf{Vấn đề}
\end{minipage} & \begin{minipage}[b]{\linewidth}\raggedright
\textbf{Tác động/Nguyên nhân}
\end{minipage} \\
\midrule\noalign{}
\endhead
\bottomrule\noalign{}
\endlastfoot
Không có dữ liệu mẫu & Demo và kiểm thử không chạy được. \\
Seed chạy nhiều lần tạo trùng dữ liệu & Menu hiển thị món lặp, tồn kho sai. \\
Dữ liệu mẫu không phủ tình huống biên & Không kiểm thử được hết hàng và đặt đồng thời. \\
\end{longtable}

\textbf{Hướng xử lý/giải pháp phân tích đề xuất}

\begin{itemize}
\item
  Viết script seed (qua Prisma) có thể chạy lại an toàn.
\item
  Chuẩn bị bộ dữ liệu mẫu có chủ đích cho kịch bản test.
\item
  Ghi cách nạp dữ liệu vào README.
\end{itemize}

\textbf{Chuyển giao sang thiết kế}

\begin{itemize}
\item
  Script seed/migration; README; cấu hình Docker Compose; dữ liệu cho test.
\end{itemize}

\textbf{Tiêu chí kiểm chứng ở mức yêu cầu}

\begin{itemize}
\item
  Sau khi chạy hệ thống theo README, menu có sản phẩm hiển thị được.
\item
  Chạy seed hai lần không làm tăng số bản ghi.
\item
  Có ít nhất một sản phẩm hết hàng và một sản phẩm tồn kho thấp dùng cho kiểm thử.
\end{itemize}

\textbf{7.4. Cart}

\textbf{Actor/phạm vi chính:} Guest; Customer; Web Client (Zustand)

\textbf{Artifact thiết kế chịu ảnh hưởng:} Use Case quản lý giỏ hàng, Cart State phía FE, UI giỏ hàng và thanh toán, Test Strategy

\textbf{FR-CART-001 - Thêm món vào giỏ}

\begin{longtable}[]{@{}
  >{\raggedright\arraybackslash}p{(\columnwidth - 2\tabcolsep) * \real{0.2500}}
  >{\raggedright\arraybackslash}p{(\columnwidth - 2\tabcolsep) * \real{0.7500}}@{}}
\toprule\noalign{}
\begin{minipage}[b]{\linewidth}\raggedright
\textbf{Thuộc tính}
\end{minipage} & \begin{minipage}[b]{\linewidth}\raggedright
\textbf{Nội dung}
\end{minipage} \\
\midrule\noalign{}
\endhead
\bottomrule\noalign{}
\endlastfoot
\textbf{Mức ưu tiên} & MUST \\
\textbf{Trạng thái} & CONFIRMED (thêm món); DERIVED (gộp dòng, giới hạn dòng và số lượng). \\
\textbf{Yêu cầu} & Guest và Customer có thể thêm một đồ uống vào giỏ kèm size, danh sách topping và số lượng. Cùng một cấu hình (sản phẩm, size, tập topping) thì gộp vào một dòng và cộng dồn số lượng. \\
\textbf{Lý do/Nguồn} & Charter §3.1 (Giỏ hàng: thêm); Task 5; Techstack §3 (Zustand). Các giới hạn và quy tắc gộp, xóa là quyết định của nhóm. \\
\textbf{Actor chính} & Guest; Customer; Web Client \\
\end{longtable}

\textbf{Phân tích nghiệp vụ}

Giỏ hàng là bước trung gian giữa xem menu và tạo đơn. Giỏ nằm hoàn toàn ở client (Zustand), không có bảng giỏ hàng hay API giỏ hàng ở backend, nên Guest vẫn thêm món được và chỉ phải đăng nhập khi tạo đơn (FR-AUTHZ-001). Mỗi dòng giỏ lưu cấu hình mà khách chọn, không lưu giá làm căn cứ; giá cuối luôn do backend tính (FR-CAT-003).

\textbf{Điều kiện, dữ liệu và ràng buộc liên quan}

\begin{itemize}
\item
  Mỗi dòng giỏ gồm: mã sản phẩm, size, danh sách mã topping, số lượng (mã dòng do client tạo).
\item
  Size thuộc \{S, M, L\}, mặc định S; tối đa 3 topping, không trùng; số lượng 1 đến 10 mỗi dòng.
\item
  Hai dòng có cùng sản phẩm, cùng size và cùng tập topping (không phân biệt thứ tự chọn) được coi là một cấu hình và gộp lại.
\item
  Giỏ có tối đa 20 dòng.
\item
  Sản phẩm hết hàng (FR-CAT-004) không thêm được qua UI.
\item
  Giá hiển thị trong giỏ chỉ là giá tạm tính tính từ dữ liệu catalog đã tải.
\end{itemize}

\textbf{Vấn đề có thể gặp và phân tích tác động}

\begin{longtable}[]{@{}
  >{\raggedright\arraybackslash}p{(\columnwidth - 2\tabcolsep) * \real{0.2500}}
  >{\raggedright\arraybackslash}p{(\columnwidth - 2\tabcolsep) * \real{0.7500}}@{}}
\toprule\noalign{}
\begin{minipage}[b]{\linewidth}\raggedright
\textbf{Vấn đề}
\end{minipage} & \begin{minipage}[b]{\linewidth}\raggedright
\textbf{Tác động/Nguyên nhân}
\end{minipage} \\
\midrule\noalign{}
\endhead
\bottomrule\noalign{}
\endlastfoot
Thêm cùng một cấu hình nhiều lần tạo nhiều dòng giống nhau & Giỏ lộn xộn, tổng khó kiểm tra. \\
Số lượng cộng dồn vượt 10 & Vi phạm giới hạn của FR-CAT-003, đơn bị từ chối. \\
Thêm món hết hàng hoặc món đã bị gỡ & Đặt đơn thất bại ở bước sau. \\
Lưu giá vào giỏ rồi dùng như giá cuối & Giá lệch khi catalog thay đổi hoặc bị sửa ở client. \\
Giỏ quá lớn & UI chậm, request tạo đơn quá nặng. \\
\end{longtable}

\textbf{Hướng xử lý/giải pháp phân tích đề xuất}

\begin{itemize}
\item
  Khi thêm, kiểm tra cấu hình hợp lệ ở client theo cùng quy tắc với FR-CAT-003, rồi gộp dòng nếu trùng cấu hình.
\item
  Nếu cộng dồn vượt 10, giữ ở 10 và thông báo rõ cho khách.
\item
  Chặn thêm khi giỏ đã đủ 20 dòng, có thông báo.
\item
  Giỏ chỉ lưu mã và lựa chọn; giá hiển thị tính lại từ catalog mỗi lần hiển thị.
\item
  Các kiểm tra ở client chỉ phục vụ UX; backend kiểm tra lại khi tạo đơn (FR-CART-005).
\end{itemize}

\textbf{Chuyển giao sang thiết kế}

\begin{itemize}
\item
  Use Case Thêm món vào giỏ; Cart State (Zustand store); UI nút thêm và thông báo lỗi; unit test logic gộp dòng.
\end{itemize}

\textbf{Tiêu chí kiểm chứng ở mức yêu cầu}

\begin{itemize}
\item
  Guest chưa đăng nhập thêm được món vào giỏ.
\item
  Thêm hai lần cùng sản phẩm, size, tập topping (thứ tự chọn khác nhau) cho một dòng với số lượng cộng dồn.
\item
  Cộng dồn vượt 10 bị giữ ở 10 và có thông báo.
\item
  Thêm món có size hoặc topping khác nhau tạo dòng riêng.
\item
  Món hết hàng không thêm được qua UI.
\item
  Giỏ đã 20 dòng thì thêm cấu hình mới bị từ chối có thông báo.
\end{itemize}

\textbf{FR-CART-002 - Cập nhật số lượng và xóa món}

\begin{longtable}[]{@{}
  >{\raggedright\arraybackslash}p{(\columnwidth - 2\tabcolsep) * \real{0.2500}}
  >{\raggedright\arraybackslash}p{(\columnwidth - 2\tabcolsep) * \real{0.7500}}@{}}
\toprule\noalign{}
\begin{minipage}[b]{\linewidth}\raggedright
\textbf{Thuộc tính}
\end{minipage} & \begin{minipage}[b]{\linewidth}\raggedright
\textbf{Nội dung}
\end{minipage} \\
\midrule\noalign{}
\endhead
\bottomrule\noalign{}
\endlastfoot
\textbf{Mức ưu tiên} & MUST \\
\textbf{Trạng thái} & CONFIRMED (cập nhật, xóa); DERIVED (giảm về 0 thì xóa). \\
\textbf{Yêu cầu} & Khách có thể tăng hoặc giảm số lượng của từng dòng trong giỏ và xóa dòng khỏi giỏ. Giảm số lượng khi dòng đang ở 1 có tác dụng như xóa dòng đó. Các giới hạn và quy tắc gộp, xóa là quyết định của nhóm. \\
\textbf{Lý do/Nguồn} & Charter §3.1 (Giỏ hàng: cập nhật số lượng, xóa); Task 5. \\
\textbf{Actor chính} & Guest; Customer; Web Client \\
\end{longtable}

\textbf{Phân tích nghiệp vụ}

Giỏ phải chỉnh được trước khi đặt đơn. Số lượng luôn nằm trong khoảng hợp lệ. Khi khách giảm từ 1 xuống 0, hệ thống hiểu là khách muốn bỏ món và xóa dòng, thay vì giữ một dòng có số lượng 0.

\textbf{Điều kiện, dữ liệu và ràng buộc liên quan}

\begin{itemize}
\item
  Số lượng của một dòng từ 1 đến 10; nút tăng dừng ở 10.
\item
  Bấm giảm khi số lượng bằng 1 xóa dòng khỏi giỏ (tương đương nút xóa); số lượng 0 không bao giờ được lưu trong giỏ.
\item
  Xóa một dòng không ảnh hưởng các dòng khác.
\item
  Khi giỏ trống, UI hiển thị trạng thái giỏ trống và không cho tạo đơn.
\item
  Thay đổi giỏ cập nhật ngay tổng tạm tính (FR-CART-003).
\end{itemize}

\textbf{Vấn đề có thể gặp và phân tích tác động}

\begin{longtable}[]{@{}
  >{\raggedright\arraybackslash}p{(\columnwidth - 2\tabcolsep) * \real{0.2500}}
  >{\raggedright\arraybackslash}p{(\columnwidth - 2\tabcolsep) * \real{0.7500}}@{}}
\toprule\noalign{}
\begin{minipage}[b]{\linewidth}\raggedright
\textbf{Vấn đề}
\end{minipage} & \begin{minipage}[b]{\linewidth}\raggedright
\textbf{Tác động/Nguyên nhân}
\end{minipage} \\
\midrule\noalign{}
\endhead
\bottomrule\noalign{}
\endlastfoot
Nhập tay số lượng âm hoặc quá lớn & Dữ liệu giỏ không hợp lệ, tổng sai. \\
Giảm từ 1 xuống 0 làm mất món ngoài ý muốn & Khách bấm giảm nhầm rồi phải chọn lại món. \\
Cập nhật trên bản sao state cũ (race trong UI) & Số lượng hiển thị lệch dữ liệu thực. \\
Giỏ trống nhưng vẫn bấm đặt & Tạo đơn rỗng. \\
\end{longtable}

\textbf{Hướng xử lý/giải pháp phân tích đề xuất}

\begin{itemize}
\item
  Giới hạn số lượng trong store, không chỉ ở UI. Hành động giảm số lượng ở mức 1 gọi thao tác xóa dòng; giá trị âm hoặc quá 10 bị chuẩn hóa hoặc từ chối.
\item
  Hành động trên store dùng id dòng, không dùng vị trí trong mảng.
\item
  Nút tạo đơn bị vô hiệu khi giỏ trống.
\item
  (SHOULD) Khi một dòng bị xóa (bằng nút xóa hoặc giảm về 0), hiển thị thông báo ngắn kèm nút ``Hoàn tác'' trong vài giây, thay cho hộp xác nhận để không cản trở thao tác.
\end{itemize}

\textbf{Chuyển giao sang thiết kế}

\begin{itemize}
\item
  Use Case Cập nhật giỏ; Cart State; UI giỏ hàng (tăng, giảm, xóa, giỏ trống); unit test store.
\end{itemize}

\textbf{Tiêu chí kiểm chứng ở mức yêu cầu}

\begin{itemize}
\item
  Tăng số lượng dừng ở 10.
\item
  Bấm giảm khi số lượng bằng 1 xóa dòng khỏi giỏ, giống nút xóa; sau đó tổng tạm tính cập nhật.
\item
  Giảm từ 3 xuống 2 chỉ đổi số lượng, không xóa dòng.
\item
  Không có dòng nào trong giỏ có số lượng nhỏ hơn 1.
\item
  Xóa dòng làm dòng biến mất và tổng tạm tính cập nhật.
\item
  Giỏ trống hiển thị trạng thái giỏ trống và không cho tạo đơn.
\item
  Thao tác trên một dòng không làm đổi dòng khác.
\end{itemize}

\textbf{FR-CART-003 - Tự tính tổng tạm tính}

\begin{longtable}[]{@{}
  >{\raggedright\arraybackslash}p{(\columnwidth - 2\tabcolsep) * \real{0.2500}}
  >{\raggedright\arraybackslash}p{(\columnwidth - 2\tabcolsep) * \real{0.7500}}@{}}
\toprule\noalign{}
\begin{minipage}[b]{\linewidth}\raggedright
\textbf{Thuộc tính}
\end{minipage} & \begin{minipage}[b]{\linewidth}\raggedright
\textbf{Nội dung}
\end{minipage} \\
\midrule\noalign{}
\endhead
\bottomrule\noalign{}
\endlastfoot
\textbf{Mức ưu tiên} & MUST \\
\textbf{Trạng thái} & CONFIRMED \\
\textbf{Yêu cầu} & Giỏ hàng phải tự tính và hiển thị thành tiền từng dòng và tổng tạm tính của cả giỏ, theo đúng công thức ở FR-CAT-003. \\
\textbf{Lý do/Nguồn} & Charter §3.1 (Giỏ hàng: tự tính tổng); Task 5. \\
\textbf{Actor chính} & Guest; Customer; Web Client \\
\end{longtable}

\textbf{Phân tích nghiệp vụ}

Khách cần thấy ngay số tiền dự kiến khi chỉnh giỏ. Đây là số \textbf{tạm tính}: chưa gồm giảm giá (coupon xử lý ở FR-PRO) và không phải giá cuối. Tổng cuối do backend quyết định khi tạo đơn, được trả về cho UI hiển thị lại trước bước thanh toán.

\textbf{Điều kiện, dữ liệu và ràng buộc liên quan}

\begin{itemize}
\item
  Thành tiền dòng = đơn giá × số lượng; đơn giá tính theo FR-CAT-003 (làm tròn half up trên đơn giá một ly).
\item
  Tổng tạm tính = tổng thành tiền các dòng.
\item
  Tiền hiển thị theo định dạng VND; tính bằng số nguyên, không dùng số thực dấu phẩy động.
\item
  UI ghi rõ nhãn ``Tạm tính''.
\end{itemize}

\textbf{Vấn đề có thể gặp và phân tích tác động}

\begin{longtable}[]{@{}
  >{\raggedright\arraybackslash}p{(\columnwidth - 2\tabcolsep) * \real{0.2500}}
  >{\raggedright\arraybackslash}p{(\columnwidth - 2\tabcolsep) * \real{0.7500}}@{}}
\toprule\noalign{}
\begin{minipage}[b]{\linewidth}\raggedright
\textbf{Vấn đề}
\end{minipage} & \begin{minipage}[b]{\linewidth}\raggedright
\textbf{Tác động/Nguyên nhân}
\end{minipage} \\
\midrule\noalign{}
\endhead
\bottomrule\noalign{}
\endlastfoot
Công thức ở FE khác BE & Tổng trong giỏ khác tổng khi thanh toán. \\
Sai số làm tròn do dùng số thực & Tổng lệch vài đồng. \\
Hiển thị tạm tính như thể đã là giá cuối & Khách hiểu sai khi tổng cuối khác. \\
Giá catalog đã đổi sau khi khách xem & Tạm tính khác tổng backend. \\
\end{longtable}

\textbf{Hướng xử lý/giải pháp phân tích đề xuất}

\begin{itemize}
\item
  Dùng chung một hàm tính giá ở FE theo cùng quy tắc với backend; kiểm tra bằng cùng bộ ví dụ test.
\item
  Tính bằng số nguyên VND.
\item
  Ghi nhãn ``Tạm tính'' và thông báo giá cuối được xác nhận khi đặt đơn.
\item
  Sau khi tạo đơn, UI hiển thị tổng từ phản hồi của backend, không dùng tổng tạm tính của giỏ.
\end{itemize}

\textbf{Chuyển giao sang thiết kế}

\begin{itemize}
\item
  Cart State (selector tính tổng); UI giỏ hàng (nhãn ``Tạm tính''); unit test cùng bộ ví dụ với FR-CAT-003.
\end{itemize}

\textbf{Tiêu chí kiểm chứng ở mức yêu cầu}

\begin{itemize}
\item
  Giỏ có một dòng 66.000đ × 2 (món 45.000đ, size M, hai topping 5.000đ và 7.000đ) và một dòng 45.000đ × 1 cho tổng tạm tính 177.000đ.
\item
  Thay đổi số lượng, size hoặc topping làm tổng cập nhật đúng.
\item
  Tổng tạm tính của FE khớp tổng backend tính cho cùng giỏ (khi catalog không đổi).
\end{itemize}

\textbf{FR-CART-004 - Giữ trạng thái giỏ và vòng đời giỏ}

\begin{longtable}[]{@{}
  >{\raggedright\arraybackslash}p{(\columnwidth - 2\tabcolsep) * \real{0.2500}}
  >{\raggedright\arraybackslash}p{(\columnwidth - 2\tabcolsep) * \real{0.7500}}@{}}
\toprule\noalign{}
\begin{minipage}[b]{\linewidth}\raggedright
\textbf{Thuộc tính}
\end{minipage} & \begin{minipage}[b]{\linewidth}\raggedright
\textbf{Nội dung}
\end{minipage} \\
\midrule\noalign{}
\endhead
\bottomrule\noalign{}
\endlastfoot
\textbf{Mức ưu tiên} & MUST \\
\textbf{Trạng thái} & CONFIRMED (giữ trạng thái cơ bản); DERIVED (các quy tắc về vòng đời) \\
\textbf{Yêu cầu} & Giỏ hàng phải được giữ khi khách tải lại trang hoặc chuyển trang. Giỏ được giữ khi đăng nhập, khi token hết hạn, khi đã tạo đơn nhưng chưa thanh toán (PENDING) và khi thanh toán thất bại; giỏ bị xóa khi đăng xuất chủ động và khi đơn thanh toán thành công (PAID). \\
\textbf{Lý do/Nguồn} & Charter §3.1 (giữ trạng thái cơ bản), §7 (PAYMENT\_FAILED giữ giỏ); FR-AUTH-004 (đăng xuất xóa giỏ). Các quy tắc còn lại là quyết định của nhóm. \\
\textbf{Actor chính} & Guest; Customer; Web Client \\
\end{longtable}

\textbf{Phân tích nghiệp vụ}

Khách không muốn mất giỏ chỉ vì tải lại trang hay vì thanh toán lỗi, nhưng giỏ cũng không nên tồn tại mãi sau khi đã mua hoặc trên thiết bị người khác. Vòng đời được quy định theo sự kiện:

\begin{longtable}[]{@{}
  >{\raggedright\arraybackslash}p{(\columnwidth - 2\tabcolsep) * \real{0.6005}}
  >{\raggedright\arraybackslash}p{(\columnwidth - 2\tabcolsep) * \real{0.3995}}@{}}
\toprule\noalign{}
\endhead
\bottomrule\noalign{}
\endlastfoot
\textbf{Sự kiện} & \textbf{Giỏ hàng} \\
Tải lại trang, chuyển trang & Giữ \\
Guest đăng nhập thành công & Giữ \\
Token hết hạn hoặc nhận 401 & Giữ \\
Tạo đơn thành công nhưng chưa thanh toán (PENDING) & Giữ \\
Thanh toán thất bại (PAYMENT\_FAILED) & Giữ \\
Thanh toán thành công (PAID) & \textbf{Xóa} \\
Đăng xuất chủ động & \textbf{Xóa} (FR-AUTH-004) \\
\end{longtable}

\textbf{Điều kiện, dữ liệu và ràng buộc liên quan}

\begin{itemize}
\item
  Giỏ được lưu bền ở trình duyệt (Zustand với persist ra localStorage), chỉ chứa mã sản phẩm, size, mã topping và số lượng; không chứa token hay thông tin cá nhân.
\item
  Khi không dùng được localStorage (chế độ riêng tư), giỏ vẫn hoạt động trong bộ nhớ nhưng mất khi tải lại, và UI không báo lỗi.
\item
  Giỏ chỉ bị xóa bởi hai sự kiện: thanh toán thành công (PAID) và đăng xuất chủ động; không có xóa ngầm.
\item
  Dữ liệu giỏ đọc từ localStorage coi là không tin cậy: phải kiểm tra lại cấu trúc và giới hạn khi khôi phục.
\end{itemize}

\textbf{Vấn đề có thể gặp và phân tích tác động}

\begin{longtable}[]{@{}
  >{\raggedright\arraybackslash}p{(\columnwidth - 2\tabcolsep) * \real{0.2500}}
  >{\raggedright\arraybackslash}p{(\columnwidth - 2\tabcolsep) * \real{0.7500}}@{}}
\toprule\noalign{}
\begin{minipage}[b]{\linewidth}\raggedright
\textbf{Vấn đề}
\end{minipage} & \begin{minipage}[b]{\linewidth}\raggedright
\textbf{Tác động/Nguyên nhân}
\end{minipage} \\
\midrule\noalign{}
\endhead
\bottomrule\noalign{}
\endlastfoot
Không xóa giỏ sau khi đã thanh toán & Khách đặt lại đúng giỏ cũ ngoài ý muốn. \\
Xóa giỏ khi thanh toán thất bại & Mất giỏ đúng lúc khách cần thử lại (trái Charter §7). \\
Giỏ còn trên máy dùng chung sau khi đăng xuất & Người sau thấy lựa chọn của người trước. \\
Dữ liệu localStorage bị sửa tay hoặc hỏng & Giỏ có cấu hình không hợp lệ. \\
Token hết hạn nhưng người khác đăng nhập trên cùng trình duyệt & Người sau thừa hưởng giỏ của người trước (hạn chế đã biết của MVP). \\
Khách đóng tab hoặc mất mạng sau khi thanh toán thành công, trước khi nhận phản hồi & Đơn đã PAID nhưng giỏ vẫn còn; khách có thể đặt lại giỏ cũ ngoài ý muốn. \\
\end{longtable}

\textbf{Hướng xử lý/giải pháp phân tích đề xuất}

\begin{itemize}
\item
  Bật persist cho store giỏ hàng; phiên bản hóa dữ liệu lưu để bỏ giá trị cũ không còn hợp lệ.
\item
  Khi khôi phục, kiểm tra từng dòng theo giới hạn ở FR-CART-001; dòng sai bị loại bỏ.
\item
  Xóa giỏ ở đúng hai điểm: ngay khi nhận kết quả thanh toán PAID và trong thao tác đăng xuất.
\item
  Ghi rõ hạn chế ``giỏ không gắn với tài khoản'' trong tài liệu.
\item
  Chấp nhận hạn chế này trong MVP; khách xác nhận qua lịch sử đơn (/orders/me) và có thể xóa giỏ thủ công.
\end{itemize}

\textbf{Chuyển giao sang thiết kế}

\begin{itemize}
\item
  Cart State (persist, version, hành động clear); luồng Auth State liên quan; Sequence thanh toán (bước xóa giỏ); test vòng đời giỏ.
\end{itemize}

\textbf{Tiêu chí kiểm chứng ở mức yêu cầu}

\begin{itemize}
\item
  Tải lại trang giữ nguyên giỏ.
\item
  Guest có giỏ đăng nhập xong vẫn còn đủ giỏ.
\item
  Nhận 401 do token hết hạn không làm mất giỏ.
\item
  Thanh toán thất bại (PAYMENT\_FAILED) giữ nguyên giỏ và khách thử lại được.
\item
  Thanh toán thành công (PAID) làm giỏ trống.
\item
  Đăng xuất chủ động làm giỏ trống, kể cả sau khi tải lại trang.
\item
  Sửa tay dữ liệu giỏ trong localStorage thành giá trị sai thì dòng sai bị loại bỏ khi tải lại.
\item
  Dòng có số lượng 0 hoặc âm trong dữ liệu localStorage bị loại bỏ khi khôi phục giỏ.
\end{itemize}

\textbf{FR-CART-005 - Giỏ chỉ là dữ liệu tạm; kiểm tra lại khi tạo đơn}

\begin{longtable}[]{@{}
  >{\raggedright\arraybackslash}p{(\columnwidth - 2\tabcolsep) * \real{0.2500}}
  >{\raggedright\arraybackslash}p{(\columnwidth - 2\tabcolsep) * \real{0.7500}}@{}}
\toprule\noalign{}
\begin{minipage}[b]{\linewidth}\raggedright
\textbf{Thuộc tính}
\end{minipage} & \begin{minipage}[b]{\linewidth}\raggedright
\textbf{Nội dung}
\end{minipage} \\
\midrule\noalign{}
\endhead
\bottomrule\noalign{}
\endlastfoot
\textbf{Mức ưu tiên} & MUST \\
\textbf{Trạng thái} & DERIVED \\
\textbf{Yêu cầu} & Giỏ hàng ở client không được coi là nguồn quyết định về giá, tồn kho hay tính hợp lệ của món. Khi tạo đơn, backend phải kiểm tra lại toàn bộ nội dung giỏ gửi lên; giỏ lỗi thời hoặc không hợp lệ phải bị từ chối kèm thông báo đủ rõ để khách sửa. \\
\textbf{Lý do/Nguồn} & Mục 3 (Web Client không phải nguồn quyết định cuối); Charter §3.1 (kiểm tra dữ liệu khi tạo đơn); FR-CAT-003. \\
\textbf{Actor chính} & Customer; Backend API; Web Client \\
\end{longtable}

\textbf{Phân tích nghiệp vụ}

Giỏ có thể cũ vì khách để mở lâu, hoặc bị chỉnh sửa trực tiếp ở client. Vì vậy khi POST /orders, backend phải coi nội dung giỏ chỉ là \emph{ý định mua} gồm mã sản phẩm, size, topping, số lượng; mọi giá, tổng tiền và điều kiện tồn kho đều tính lại từ dữ liệu server. Kết quả kiểm tra phải đủ cụ thể để UI cho khách sửa giỏ mà không phải đoán.

\textbf{Điều kiện, dữ liệu và ràng buộc liên quan}

\begin{itemize}
\item
  Payload tạo đơn chỉ gồm danh sách dòng (mã sản phẩm, size, mã topping, số lượng); không nhận giá, đơn giá hay tổng tiền.
\item
  Backend kiểm tra: giỏ không trống và tối đa 20 dòng; sản phẩm còn tồn tại; size và topping hợp lệ; không trùng topping; tối đa 3 topping; số lượng 1 đến 10; không có hai dòng trùng cấu hình; tồn kho đủ.
\item
  Backend từ chối (400) request có hai dòng cùng sản phẩm, size và tập topping, vì client hợp lệ luôn gộp dòng trước khi gửi.
\item
  Mã lỗi: dữ liệu giỏ không hợp lệ (sản phẩm không còn, size hoặc topping sai, số lượng ngoài khoảng, giỏ trống) trả 400 và chỉ ra dòng lỗi; không đủ tồn kho trả 409 (FR-INV-002).
\item
  Thông tin tồn kho và trạng thái hàng của FE chỉ mang tính tham khảo.
\item
  (SHOULD) UI làm mới dữ liệu catalog khi mở giỏ hoặc vào bước thanh toán, đánh dấu dòng không còn hợp lệ.
\item
  Backend kiểm tra dữ liệu giỏ trước (lỗi thì trả 400), rồi mới kiểm tra tồn kho (thiếu thì trả 409). Giỏ có đồng thời hai loại lỗi thì trả 400.
\end{itemize}

\textbf{Vấn đề có thể gặp và phân tích tác động}

\begin{longtable}[]{@{}
  >{\raggedright\arraybackslash}p{(\columnwidth - 2\tabcolsep) * \real{0.2500}}
  >{\raggedright\arraybackslash}p{(\columnwidth - 2\tabcolsep) * \real{0.7500}}@{}}
\toprule\noalign{}
\begin{minipage}[b]{\linewidth}\raggedright
\textbf{Vấn đề}
\end{minipage} & \begin{minipage}[b]{\linewidth}\raggedright
\textbf{Tác động/Nguyên nhân}
\end{minipage} \\
\midrule\noalign{}
\endhead
\bottomrule\noalign{}
\endlastfoot
Giỏ giữ món đã hết hàng hoặc món không còn trong catalog & Tạo đơn thất bại dù khách tưởng còn. \\
Client gửi giá thấp hơn hoặc tổng tự tính & Thất thoát doanh thu nếu backend tin client. \\
Lỗi trả về chung chung & Khách không biết sửa dòng nào. \\
Bấm tạo đơn nhiều lần liên tiếp & Tạo trùng đơn và giữ tồn kho nhiều lần (xử lý ở FR-ORD-006). \\
Giá catalog đổi so với lúc khách xem & Tổng cuối khác tạm tính. \\
\end{longtable}

\textbf{Hướng xử lý/giải pháp phân tích đề xuất}

\begin{itemize}
\item
  DTO tạo đơn chỉ nhận các trường được liệt kê; trường lạ bị bỏ qua hoặc từ chối.
\item
  Backend tính lại giá từ catalog và trả tổng cuối trong phản hồi; UI hiển thị lại trước bước thanh toán.
\item
  Lỗi trả theo dòng (ví dụ ``dòng 2: hết hàng'') để UI đánh dấu và hướng dẫn sửa.
\item
  Khi lỗi, giỏ ở client được giữ nguyên để khách chỉnh rồi thử lại.
\item
  Cách chống tạo trùng đơn do bấm nhiều lần do FR-ORD-006 quyết định.
\end{itemize}

\textbf{Chuyển giao sang thiết kế}

\begin{itemize}
\item
  API POST /orders (DTO, mã lỗi 400 và 409, định dạng lỗi theo dòng); Sequence tạo đơn; UI giỏ hàng (đánh dấu dòng lỗi); test API với giỏ không hợp lệ.
\end{itemize}

\textbf{Tiêu chí kiểm chứng ở mức yêu cầu}

\begin{itemize}
\item
  Gửi request tạo đơn có thêm trường giá hoặc tổng tiền thì đơn tạo ra vẫn dùng giá từ catalog.
\item
  Giỏ có món không tồn tại, size lạ, số lượng ngoài 1 đến 10 hoặc giỏ trống bị từ chối 400, chỉ rõ dòng lỗi.
\item
  Giỏ có món vượt tồn kho bị từ chối 409.
\item
  Sau khi bị từ chối, giỏ ở client vẫn còn và khách sửa được.
\item
  Gọi API trực tiếp (không qua UI) với giỏ không hợp lệ vẫn bị từ chối.
\item
  Request có hơn 20 dòng hoặc có hai dòng trùng cấu hình bị từ chối 400.
\end{itemize}

\textbf{7.5. Order}

\textbf{Actor/phạm vi chính:} Customer; Backend API (module Order); hệ thống (tác vụ định kỳ)

\textbf{Artifact thiết kế chịu ảnh hưởng:} Use Case tạo đơn và lịch sử đơn, ERD (Order, OrderItem), State Diagram Order, Sequence tạo đơn, API Order, UI giỏ hàng và lịch sử đơn, Test Strategy

\textbf{FR-ORD-001 - Tạo đơn từ giỏ hàng}

\begin{longtable}[]{@{}
  >{\raggedright\arraybackslash}p{(\columnwidth - 2\tabcolsep) * \real{0.2500}}
  >{\raggedright\arraybackslash}p{(\columnwidth - 2\tabcolsep) * \real{0.7500}}@{}}
\toprule\noalign{}
\begin{minipage}[b]{\linewidth}\raggedright
\textbf{Thuộc tính}
\end{minipage} & \begin{minipage}[b]{\linewidth}\raggedright
\textbf{Nội dung}
\end{minipage} \\
\midrule\noalign{}
\endhead
\bottomrule\noalign{}
\endlastfoot
\textbf{Mức ưu tiên} & MUST \\
\textbf{Trạng thái} & CONFIRMED \\
\textbf{Yêu cầu} & Customer đã đăng nhập tạo được đơn từ giỏ hàng. Backend kiểm tra lại giỏ (FR-CART-005), tính lại giá (FR-CAT-003), giữ chỗ tồn kho và lưu đơn ở trạng thái PENDING trong cùng một giao dịch. Phản hồi gồm mã đơn, các dòng hàng và tổng tiền do server tính. \\
\textbf{Lý do/Nguồn} & Charter §3.1 (Đơn hàng: tạo từ giỏ, kiểm tra dữ liệu, lưu PENDING, trả mã đơn); Task 6. \\
\textbf{Actor chính} & Customer; Backend API \\
\end{longtable}

\textbf{Phân tích nghiệp vụ}

Tạo đơn là bước chuyển từ ý định mua sang cam kết nghiệp vụ: giá được chốt và tồn kho được giữ. Vì vậy việc kiểm tra, giữ tồn kho và lưu đơn phải xảy ra trọn vẹn hoặc không xảy ra gì.

\textbf{Điều kiện, dữ liệu và ràng buộc liên quan}

\begin{itemize}
\item
  Endpoint POST /orders, yêu cầu JWT; userId lấy từ token (FR-AUTHZ-002).
\item
  Payload gồm danh sách dòng hàng (và mã giảm giá tùy chọn, xử lý ở FR-PRO); không nhận giá hay tổng tiền.
\item
  Thứ tự kiểm tra: xác thực (401) → dữ liệu giỏ và mã giảm giá (400) → đã có đơn PENDING hoặc PAYMENT\_FAILED (409) → tồn kho (409 INSUFFICIENT\_STOCK) → lượt dùng mã (409 COUPON\_UNAVAILABLE)''; thêm mã máy đọc INVALID\_COUPON, COUPON\_UNAVAILABLE.
\item
  Trừ tồn kho và lưu đơn trong một transaction; lỗi ở bất kỳ bước nào thì không có thay đổi nào được lưu.
\item
  Tổng tiền = tổng thành tiền các dòng − giảm giá; tổng không âm; tiền là số nguyên VND.
\item
  Phản hồi thành công gồm: mã đơn, trạng thái PENDING, các dòng hàng, tạm tính, giảm giá, tổng tiền, thời điểm tạo.
\item
  Lỗi trả kèm mã máy đọc được (INVALID\_CART, INSUFFICIENT\_STOCK, PENDING\_ORDER\_EXISTS) để UI phân biệt.
\item
  Mã đơn duy nhất và không phải số tăng dần dễ đoán.
\item
  Backend hoàn toàn bỏ qua các trường price, lineTotal hoặc total nếu client (giao diện Zustand) cố tình gửi lên trong Payload. Mọi giá trị tiền tệ phải query từ Database (Prisma) ngay tại thời điểm tạo transaction.
\end{itemize}

\textbf{Vấn đề có thể gặp và phân tích tác động}

\begin{longtable}[]{@{}
  >{\raggedright\arraybackslash}p{(\columnwidth - 2\tabcolsep) * \real{0.2500}}
  >{\raggedright\arraybackslash}p{(\columnwidth - 2\tabcolsep) * \real{0.7500}}@{}}
\toprule\noalign{}
\begin{minipage}[b]{\linewidth}\raggedright
\textbf{Vấn đề}
\end{minipage} & \begin{minipage}[b]{\linewidth}\raggedright
\textbf{Tác động/Nguyên nhân}
\end{minipage} \\
\midrule\noalign{}
\endhead
\bottomrule\noalign{}
\endlastfoot
Trừ tồn kho nhưng lưu đơn thất bại (hoặc ngược lại) & Tồn kho lệch, có đơn ``ma'' hoặc mất hàng không có đơn. \\
Hai khách đặt đồng thời món cuối cùng & Bán vượt tồn kho (Charter R5, xử lý ở FR-INV-002). \\
Phản hồi bị thất lạc do mạng & Khách không biết đơn đã tạo, dễ bấm tạo lần nữa (xử lý ở FR-ORD-006). \\
Client gửi kèm giá hoặc tổng & Giá sai nếu backend tin client. \\
\end{longtable}

\textbf{Hướng xử lý/giải pháp phân tích đề xuất}

\begin{itemize}
\item
  Thực hiện kiểm tra, trừ tồn và lưu đơn trong một transaction; lỗi thì rollback toàn bộ.
\item
  Dùng lại hàm tính giá dùng chung của FR-CAT-003; tồn kho theo FR-INV.
\item
  Khi tạo đơn thành công, UI chuyển sang bước thanh toán và vẫn giữ giỏ (FR-CART-004).
\item
  Lỗi trả theo dòng và kèm mã lỗi để UI hướng dẫn khách sửa.
\end{itemize}

\textbf{Chuyển giao sang thiết kế}

\begin{itemize}
\item
  Use Case Tạo đơn; ERD Order, OrderItem; API POST /orders (DTO, mã lỗi); Sequence tạo đơn (kiểm tra → trừ tồn → lưu); unit và integration test, concurrency test.
\end{itemize}

\textbf{Tiêu chí kiểm chứng ở mức yêu cầu}

\begin{itemize}
\item
  Customer có JWT và giỏ hợp lệ tạo được đơn PENDING có mã đơn; tổng khớp tính tay.
\item
  Gọi không có token trả 401.
\item
  Giỏ hai dòng mà dòng thứ hai thiếu tồn kho trả 409; không có đơn nào được tạo và tồn kho dòng thứ nhất không đổi.
\item
  Sau khi tạo đơn thành công, tồn kho giảm đúng bằng số lượng đã đặt.
\item
  Request có thêm trường giá hoặc tổng thì đơn vẫn dùng giá từ catalog.
\item
  Các lỗi trả đúng mã INVALID\_CART, INSUFFICIENT\_STOCK, PENDING\_ORDER\_EXISTS.
\end{itemize}

\textbf{FR-ORD-002 - Lưu bản chụp (snapshot) các dòng hàng}

\begin{longtable}[]{@{}
  >{\raggedright\arraybackslash}p{(\columnwidth - 2\tabcolsep) * \real{0.2500}}
  >{\raggedright\arraybackslash}p{(\columnwidth - 2\tabcolsep) * \real{0.7500}}@{}}
\toprule\noalign{}
\begin{minipage}[b]{\linewidth}\raggedright
\textbf{Thuộc tính}
\end{minipage} & \begin{minipage}[b]{\linewidth}\raggedright
\textbf{Nội dung}
\end{minipage} \\
\midrule\noalign{}
\endhead
\bottomrule\noalign{}
\endlastfoot
\textbf{Mức ưu tiên} & MUST \\
\textbf{Trạng thái} & CONFIRMED (lưu dòng hàng); DERIVED (nội dung snapshot) \\
\textbf{Yêu cầu} & Mỗi đơn lưu bản chụp bất biến của từng dòng hàng tại thời điểm tạo: tên món, size, các topping kèm giá, đơn giá một ly, số lượng và thành tiền dòng. Thay đổi catalog sau đó không làm đổi đơn đã tạo. \\
\textbf{Lý do/Nguồn} & Techstack §6 (OrderItem); Charter §3.1; FR-CAT-003. Phạm vi nội dung snapshot là quyết định của nhóm. \\
\textbf{Actor chính} & Backend API \\
\end{longtable}

\textbf{Phân tích nghiệp vụ}

Lịch sử đơn là bản ghi giao dịch nên phải giữ nguyên những gì khách đã thấy và đã trả. Nếu chỉ lưu khóa ngoại sang sản phẩm và topping, đổi giá hoặc đổi tên món sau này sẽ làm lịch sử và tổng tiền hiển thị sai.

\textbf{Điều kiện, dữ liệu và ràng buộc liên quan}

\begin{itemize}
\item
  Mỗi dòng lưu giá trị cụ thể (tên, size, topping với tên và giá, đơn giá một ly, số lượng, thành tiền); mã sản phẩm vẫn lưu để tham chiếu.
\item
  Tổng đơn = tổng thành tiền các dòng − giảm giá.
\item
  Không có thao tác cập nhật hoặc xóa dòng hàng của đơn qua API.
\end{itemize}

\textbf{Vấn đề có thể gặp và phân tích tác động}

\begin{longtable}[]{@{}
  >{\raggedright\arraybackslash}p{(\columnwidth - 2\tabcolsep) * \real{0.2500}}
  >{\raggedright\arraybackslash}p{(\columnwidth - 2\tabcolsep) * \real{0.7500}}@{}}
\toprule\noalign{}
\begin{minipage}[b]{\linewidth}\raggedright
\textbf{Vấn đề}
\end{minipage} & \begin{minipage}[b]{\linewidth}\raggedright
\textbf{Tác động/Nguyên nhân}
\end{minipage} \\
\midrule\noalign{}
\endhead
\bottomrule\noalign{}
\endlastfoot
Lịch sử lấy giá hiện tại từ catalog & Đơn cũ hiển thị giá sai sau khi giá đổi. \\
Topping chỉ lưu mã & Đổi hoặc xóa topping làm mất thông tin đơn cũ. \\
Tổng đơn không khớp tổng các dòng & Sai lệch khi đối soát và khi thanh toán. \\
\end{longtable}

\textbf{Hướng xử lý/giải pháp phân tích đề xuất}

\begin{itemize}
\item
  Sao chép giá trị vào dòng hàng khi tạo đơn, không chỉ lưu khóa ngoại.
\item
  Kiểm tra bằng test: tổng các dòng trừ giảm giá bằng tổng đơn.
\item
  Không cung cấp API sửa dòng hàng.
\end{itemize}

\textbf{Chuyển giao sang thiết kế}

\begin{itemize}
\item
  ERD OrderItem (đơn giá một ly, topping snapshot kèm giá, tên món); API trả snapshot cho lịch sử; test snapshot.
\end{itemize}

\textbf{Tiêu chí kiểm chứng ở mức yêu cầu}

\begin{itemize}
\item
  Đổi giá và tên món trong database sau khi tạo đơn thì lịch sử vẫn hiển thị giá và tên cũ.
\item
  Đổi hoặc xóa topping không làm đổi dòng hàng của đơn cũ.
\item
  Với mọi đơn, tổng thành tiền các dòng trừ giảm giá bằng tổng đơn.
\end{itemize}

\textbf{FR-ORD-003 - Vòng đời trạng thái đơn (state machine)}

\begin{longtable}[]{@{}
  >{\raggedright\arraybackslash}p{(\columnwidth - 2\tabcolsep) * \real{0.2500}}
  >{\raggedright\arraybackslash}p{(\columnwidth - 2\tabcolsep) * \real{0.7500}}@{}}
\toprule\noalign{}
\begin{minipage}[b]{\linewidth}\raggedright
\textbf{Thuộc tính}
\end{minipage} & \begin{minipage}[b]{\linewidth}\raggedright
\textbf{Nội dung}
\end{minipage} \\
\midrule\noalign{}
\endhead
\bottomrule\noalign{}
\endlastfoot
\textbf{Mức ưu tiên} & MUST \\
\textbf{Trạng thái} & CONFIRMED \\
\textbf{Yêu cầu} & Đơn chỉ được chuyển trạng thái theo các bước hợp lệ ở bảng dưới; mọi chuyển trạng thái khác phải bị backend chặn (409). Các bước chuyển được thực hiện qua hành động nghiệp vụ, không có thao tác gán trạng thái tùy ý. \\
\textbf{Lý do/Nguồn} & Charter §3.2 (Task 10 vẫn phải mô hình hóa trạng thái đơn), §7 (chặn transition sai); Techstack §7; Task 10. \\
\textbf{Actor chính} & Customer; hệ thống thanh toán; hệ thống (tác vụ định kỳ); thao tác nội bộ \\
\end{longtable}

\textbf{Phân tích nghiệp vụ}

Trạng thái đơn quyết định điều gì được phép xảy ra tiếp theo (thanh toán, hủy, chuẩn bị) và khi nào tồn kho được hoàn lại hoặc điểm được cộng. Một chuyển trạng thái sai sẽ làm sai cả tiền, tồn kho và điểm, nên các bước và tác động đi kèm được quy định rõ.

\begin{longtable}[]{@{}
  >{\raggedright\arraybackslash}p{(\columnwidth - 6\tabcolsep) * \real{0.2735}}
  >{\raggedright\arraybackslash}p{(\columnwidth - 6\tabcolsep) * \real{0.2735}}
  >{\raggedright\arraybackslash}p{(\columnwidth - 6\tabcolsep) * \real{0.2283}}
  >{\raggedright\arraybackslash}p{(\columnwidth - 6\tabcolsep) * \real{0.2246}}@{}}
\toprule\noalign{}
\endhead
\bottomrule\noalign{}
\endlastfoot
\textbf{Từ} & \textbf{Đến} & \textbf{Sự kiện kích hoạt} & \textbf{Tác động đi kèm} \\
(tạo mới) & PENDING & Tạo đơn thành công (FR-ORD-001) & Trừ tồn kho, giữ lượt mã giảm giá (nếu có). \\
PENDING & PAID & Thanh toán mock thành công (FR-PAY-002) & Xác nhận lượt mã đã giữ, Cộng điểm tích lũy một lần (FR-PRO-003), xóa giỏ hàng. \\
PENDING & PAYMENT\_FAILED & Thanh toán mock thất bại (FR-PAY-002) & Giữ tồn kho đã đặt trước, giữ lượt mã giảm giá, giữ giỏ \\
PEYMENT\_FAILED & PENDING & Customer yêu cầu thử lại thanh toán (FR-PAY-003) & Không tạo mới đơn, không trừ tồn kho thêm, không giữ lượt mã thêm. \\
PENDING & CANCELLED & Customer hủy (FR-ORD-005) hoặc đơn quá hạn (FR-ORD-007) & Hoàn tồn kho và trả lượt mã đã giữ đúng một lần. \\
PAYMENT\_FAILED & CANCELLED & Customer hủy (FR-ORD-005) hoặc đơn quá hạn (FR-ORD-007) & Hoàn tồn kho và trả lượt mã đã giữ đúng một lần. \\
PAID & PREPARING & Thao tác nội bộ & Không thay đổi tồn kho, lượt mã hay điểm. \\
PREPARING & READY & Thao tác nội bộ & Không. \\
READY & COMPLETED & Thao tác nội bộ & Không. \\
\end{longtable}

Các trạng thái kết thúc của đơn hàng là COMPLETED và CANCELLED. PAYMENT\_FAILED không phải trạng thái kết thúc, khách hàng được phép thử lại thanh toán trên cùng đơn khi đơn chưa hết hạn.

Khi thử lại, đơn chuyển từ PAYMENT\_FAILED về PENDING trước khi thực hiện lần thanh toán mới. Quá trình này không tạo đơn hàng mới, không trừ tồn kho lần thứ hai và không giữ thêm lượt sử dụng mã giảm giá.

Mỗi lần thử thanh toán mới phải sử dụng Idempotency-Key mới. Việc gửi lại cùng Idempotency-Key và cùng nội dung yêu cầu phải trả về kết quả đã lưu, không thực hiện lại giao dịch.

Mọi chuyển trạng thái không được khai báo trong bảng đều bị backend từ chối. Đặc biệt, đơn đã PAID không được hủy hoặc quay lại trạng thái chờ thanh toán trong phạm vi MVP.

\textbf{Điều kiện, dữ liệu và ràng buộc liên quan}

\begin{itemize}
\item
  Chỉ luồng thanh toán được chuyển PENDING sang PAID hoặc PAYMENT\_FAILED; Customer không gọi trực tiếp được.
\item
  Ba bước PAID → PREPARING → READY → COMPLETED do thao tác nội bộ (đại diện Barista, vì MVP không có giao diện Barista) thực hiện: endpoint bảo vệ bằng khóa cấu hình, chỉ nhận hành động ``chuyển sang bước kế tiếp'', không nhận trạng thái đích tùy ý; thiếu khóa trả 403 (FR-AUTHZ-001). Chi tiết ở mục 11.
\item
  Mỗi lần chuyển trạng thái và tác động đi kèm (hoàn tồn kho, cộng điểm) nằm trong cùng một transaction.
\item
  Chuyển trạng thái sai trả 409 với mã INVALID\_TRANSITION.
\item
  Chỉ cho phép PAYMENT\_FAILED → PENDING khi đơn chưa hết hạn và người yêu cầu là chủ sở hữu đơn.
\item
  Mỗi lần thử thanh toán mới phải dùng Idempotency-Key mới; gửi lại cùng khóa cũ phải trả kết quả thanh toán đã được lưu.
\item
  Không trừ tồn kho hoặc giữ lượt mã giảm giá thêm khi thử lại.
\item
  Thanh toán thành công phải chặn mọi lần thử mới.
\item
  PAYMENT\_FAILED phải được xử lý khi đơn hết hạn tương tự PENDING để không giữ tồn kho vô thời hạn.
\end{itemize}

\textbf{Vấn đề có thể gặp và phân tích tác động}

\begin{longtable}[]{@{}
  >{\raggedright\arraybackslash}p{(\columnwidth - 2\tabcolsep) * \real{0.2500}}
  >{\raggedright\arraybackslash}p{(\columnwidth - 2\tabcolsep) * \real{0.7500}}@{}}
\toprule\noalign{}
\begin{minipage}[b]{\linewidth}\raggedright
\textbf{Vấn đề}
\end{minipage} & \begin{minipage}[b]{\linewidth}\raggedright
\textbf{Tác động/Nguyên nhân}
\end{minipage} \\
\midrule\noalign{}
\endhead
\bottomrule\noalign{}
\endlastfoot
Cho gán trạng thái tùy ý & Phá vòng đời Order và Payment, tự đánh dấu đơn đã trả tiền. \\
Hủy và thanh toán xảy ra đồng thời & Đơn vừa CANCELLED (đã hoàn tồn) vừa PAID. \\
Hoàn tồn kho hai lần (retry, double-click) & Tồn kho tăng sai. \\
Cộng điểm khi đơn chưa PAID hoặc cộng nhiều lần & Điểm sai (FR-PRO-003). \\
Customer gọi được thao tác nội bộ & Tự chuyển đơn sang COMPLETED. \\
\end{longtable}

\textbf{Hướng xử lý/giải pháp phân tích đề xuất}

\begin{itemize}
\item
  Mỗi chuyển trạng thái là một hàm riêng kiểm tra trạng thái hiện tại (không có setStatus chung).
\item
  Cập nhật có điều kiện theo trạng thái hiện tại (chỉ ghi khi đơn đang ở đúng trạng thái xuất phát) để hai thao tác đồng thời chỉ một thao tác thành công; thao tác còn lại nhận 409.
\item
  Hoàn tồn kho và cộng điểm chỉ xảy ra khi chuyển trạng thái thật sự thành công, nên mỗi tác động chạy đúng một lần.
\item
  Endpoint nội bộ tách khỏi API của Customer và bảo vệ bằng khóa.
\item
  Mặc dù không có UI cho Barista, Backend cần cung cấp một endpoint nội bộ (ví dụ: PATCH /internal/orders/:id/status) bảo vệ bằng API Key cấu hình trong biến môi trường .env để nhóm dùng Postman gọi chuyển trạng thái khi quay video demo nghiệm thu.
\end{itemize}

\textbf{Chuyển giao sang thiết kế}

\begin{itemize}
\item
  State Diagram Order (mục 11); ERD (cột trạng thái dạng enum); Sequence thanh toán và hủy; test ma trận chuyển trạng thái; concurrency test cho hủy so với thanh toán.
\end{itemize}

\textbf{Tiêu chí kiểm chứng ở mức yêu cầu}

\begin{itemize}
\item
  Mỗi cặp (trạng thái nguồn, trạng thái đích) ngoài bảng đều bị từ chối 409 (kiểm thử toàn bộ ma trận).
\item
  Đơn đi qua chuỗi PENDING → PAID → PREPARING → READY → COMPLETED đúng thứ tự.
\item
  Hủy và thanh toán đồng thời cùng một đơn chỉ một thao tác thành công; tồn kho không bị hoàn sai.
\item
  Customer gọi thao tác nội bộ trả 403.
\item
  Khi Order chuyển sang PAYMENT\_FAILED, tồn kho tiếp tục được giữ và không hoàn lại. Chỉ khi Order chuyển từ PENDING hoặc PAYMENT\_FAILED sang CANCELLED, hệ thống mới hoàn tồn kho đúng một lần trong cùng transaction với thao tác chuyển trạng thái.
\item
  Gọi endpoint nội bộ trên đơn PENDING, PAYMENT\_FAILED hoặc CANCELLED trả 409; gọi trên đơn COMPLETED trả 409.
\end{itemize}

\textbf{FR-ORD-004 - Lịch sử đơn của tôi}

\begin{longtable}[]{@{}
  >{\raggedright\arraybackslash}p{(\columnwidth - 2\tabcolsep) * \real{0.2500}}
  >{\raggedright\arraybackslash}p{(\columnwidth - 2\tabcolsep) * \real{0.7500}}@{}}
\toprule\noalign{}
\begin{minipage}[b]{\linewidth}\raggedright
\textbf{Thuộc tính}
\end{minipage} & \begin{minipage}[b]{\linewidth}\raggedright
\textbf{Nội dung}
\end{minipage} \\
\midrule\noalign{}
\endhead
\bottomrule\noalign{}
\endlastfoot
\textbf{Mức ưu tiên} & MUST \\
\textbf{Trạng thái} & CONFIRMED \\
\textbf{Yêu cầu} & Customer xem được danh sách đơn của chính mình, mới nhất trước, mỗi đơn gồm mã đơn, thời điểm tạo, trạng thái, tổng tiền và các dòng hàng đã chụp. \\
\textbf{Lý do/Nguồn} & Charter §3.1 (lịch sử đơn); Techstack §6 (GET /orders/me); Task 9. \\
\textbf{Actor chính} & Customer; Backend API \\
\end{longtable}

\textbf{Phân tích nghiệp vụ}

Lịch sử đơn giúp khách xác nhận món đã đặt, theo dõi trạng thái và tìm đơn đang chờ thanh toán. Chỉ trả đơn của chính người dùng trong token (FR-AUTHZ-002).

\textbf{Điều kiện, dữ liệu và ràng buộc liên quan}

\begin{itemize}
\item
  Endpoint GET /orders/me, yêu cầu JWT; lọc theo userId trong token.
\item
  Sắp xếp theo thời điểm tạo giảm dần, có khóa phụ là mã đơn để thứ tự ổn định.
\item
  Nội dung dòng hàng lấy từ snapshot (FR-ORD-002).
\item
  (SHOULD) hỗ trợ phân trang tùy chọn theo cùng quy ước của FR-CAT-001; không có tham số thì trả toàn bộ.
\item
  Không trả dữ liệu nội bộ (khóa idempotency, phiên bản tồn kho).
\item
  Chưa có đơn nào thì trả danh sách rỗng (không lỗi).
\end{itemize}

\textbf{Vấn đề có thể gặp và phân tích tác động}

\begin{longtable}[]{@{}
  >{\raggedright\arraybackslash}p{(\columnwidth - 2\tabcolsep) * \real{0.2500}}
  >{\raggedright\arraybackslash}p{(\columnwidth - 2\tabcolsep) * \real{0.7500}}@{}}
\toprule\noalign{}
\begin{minipage}[b]{\linewidth}\raggedright
\textbf{Vấn đề}
\end{minipage} & \begin{minipage}[b]{\linewidth}\raggedright
\textbf{Tác động/Nguyên nhân}
\end{minipage} \\
\midrule\noalign{}
\endhead
\bottomrule\noalign{}
\endlastfoot
Trả đơn của người khác (thiếu lọc theo token) & Lộ dữ liệu (IDOR, FR-AUTHZ-002). \\
Lịch sử lấy giá và tên hiện tại thay vì snapshot & Hiển thị sai đơn cũ. \\
Danh sách rất dài & Phản hồi chậm. \\
Khách mới chưa có đơn & Màn hình trống gây hiểu nhầm là lỗi. \\
\end{longtable}

\textbf{Hướng xử lý/giải pháp phân tích đề xuất}

\begin{itemize}
\item
  Truy vấn kèm userId từ token; trả qua DTO chỉ gồm trường hiển thị.
\item
  UI có trạng thái loading, empty (``chưa có đơn nào''), error; nhãn trạng thái dễ hiểu cho từng trạng thái.
\item
  Đơn PENDING được làm nổi bật kèm nút thanh toán hoặc hủy (FR-ORD-006).
\end{itemize}

\textbf{Chuyển giao sang thiết kế}

\begin{itemize}
\item
  Use Case Xem lịch sử đơn; API GET /orders/me (DTO, sắp xếp, phân trang tùy chọn); UI trang lịch sử và trang xác nhận; test API.
\end{itemize}

\textbf{Tiêu chí kiểm chứng ở mức yêu cầu}

\begin{itemize}
\item
  Customer A chỉ thấy đơn của mình; không thấy đơn của Customer B.
\item
  Gọi không token trả 401.
\item
  Đơn hiển thị mới nhất trước và đúng snapshot (đổi giá catalog không đổi lịch sử).
\item
  Khách chưa có đơn nhận danh sách rỗng và UI hiển thị trạng thái rỗng.
\item
  Thời gian phản hồi với dữ liệu mẫu \textless{} 500 ms.
\end{itemize}

\textbf{FR-ORD-005 - Hủy đơn chờ thanh toán}

\begin{longtable}[]{@{}
  >{\raggedright\arraybackslash}p{(\columnwidth - 2\tabcolsep) * \real{0.2500}}
  >{\raggedright\arraybackslash}p{(\columnwidth - 2\tabcolsep) * \real{0.7500}}@{}}
\toprule\noalign{}
\begin{minipage}[b]{\linewidth}\raggedright
\textbf{Thuộc tính}
\end{minipage} & \begin{minipage}[b]{\linewidth}\raggedright
\textbf{Nội dung}
\end{minipage} \\
\midrule\noalign{}
\endhead
\bottomrule\noalign{}
\endlastfoot
\textbf{Mức ưu tiên} & SHOULD \\
\textbf{Trạng thái} & DERIVED \\
\textbf{Yêu cầu} & Customer được phép hủy Order thuộc sở hữu của mình khi Order ở trạng thái PENDING hoặc PAYMENT\_FAILED.

Khi hủy thành công, Order chuyển sang CANCELLED. Hệ thống hoàn tồn kho và trả lượt sử dụng mã giảm giá đã giữ đúng một lần.

Order từ PAID trở đi hoặc đã CANCELLED không được hủy trong phạm vi MVP. \\
\textbf{Lý do/Nguồn} & Techstack §7 (CANCELLED); Charter không nêu cách hủy. \\
\textbf{Actor chính} & Customer; Backend API \\
\end{longtable}

\textbf{Phân tích nghiệp vụ}

Do mỗi Customer chỉ có một đơn PENDING (FR-ORD-006), khách cần cách bỏ đơn đó để đặt đơn khác. Hủy sau khi đã PAID kéo theo hoàn tiền và thu hồi điểm, ngoài phạm vi MVP nên không hỗ trợ.

\textbf{Điều kiện, dữ liệu và ràng buộc liên quan}

\begin{itemize}
\item
  Endpoint hủy yêu cầu JWT; chỉ chủ đơn hủy được (đơn của người khác hoặc không tồn tại trả 404).
\item
  Endpoint hủy yêu cầu JWT hợp lệ và kiểm tra quyền sở hữu Order.
\item
  Chỉ cho phép hủy Order ở trạng thái PENDING hoặc PAYMENT\_FAILED.
\item
  Việc chuyển sang CANCELLED, hoàn tồn kho và trả lượt mã phải nằm trong cùng transaction.
\item
  Hai yêu cầu hủy đồng thời chỉ được có một yêu cầu thành công.
\item
  Giỏ hàng của Customer được giữ nguyên khi hủy đơn.
\item
  Hoàn tồn kho cùng transaction với việc đổi trạng thái (FR-ORD-003).
\item
  Giỏ ở client không bị xóa khi hủy đơn.
\end{itemize}

\textbf{Vấn đề có thể gặp và phân tích tác động}

\begin{longtable}[]{@{}
  >{\raggedright\arraybackslash}p{(\columnwidth - 2\tabcolsep) * \real{0.2500}}
  >{\raggedright\arraybackslash}p{(\columnwidth - 2\tabcolsep) * \real{0.7500}}@{}}
\toprule\noalign{}
\begin{minipage}[b]{\linewidth}\raggedright
\textbf{Vấn đề}
\end{minipage} & \begin{minipage}[b]{\linewidth}\raggedright
\textbf{Tác động/Nguyên nhân}
\end{minipage} \\
\midrule\noalign{}
\endhead
\bottomrule\noalign{}
\endlastfoot
Hủy đơn của người khác & Phá đơn hàng người khác (IDOR). \\
Hủy đúng lúc đang thanh toán & Đơn vừa hủy vừa thanh toán (xử lý ở FR-ORD-003). \\
Bấm hủy hai lần & Hoàn tồn kho hai lần nếu không chặn. \\
\end{longtable}

\textbf{Hướng xử lý/giải pháp phân tích đề xuất}

\begin{itemize}
\item
  Kiểm tra sở hữu trước, rồi kiểm tra trạng thái (FR-AUTHZ-003).
\item
  Cập nhật có điều kiện theo trạng thái để lần hủy thứ hai nhận 409 và không hoàn tồn lần nữa.
\end{itemize}

\textbf{Chuyển giao sang thiết kế}

\begin{itemize}
\item
  API hủy đơn (bổ sung vào Techstack §6); Use Case Hủy đơn; UI nút hủy ở lịch sử đơn; test hủy.
\end{itemize}

\textbf{Tiêu chí kiểm chứng ở mức yêu cầu}

\begin{itemize}
\item
  Hủy đơn PENDING của mình thành công; tồn kho trở lại như trước khi đặt.
\item
  Hủy đơn PAID hoặc đơn đã hủy trả 409.
\item
  Hủy đơn của người khác trả 404.
\item
  Bấm hủy hai lần liên tiếp chỉ hoàn tồn kho một lần.
\item
  Hủy PAYMENT\_FAILED và kiểm tra không hoàn tồn kho hai lần.
\end{itemize}

\textbf{FR-ORD-006 - Một đơn chờ thanh toán mỗi khách và chống tạo trùng đơn}

\begin{longtable}[]{@{}
  >{\raggedright\arraybackslash}p{(\columnwidth - 2\tabcolsep) * \real{0.2500}}
  >{\raggedright\arraybackslash}p{(\columnwidth - 2\tabcolsep) * \real{0.7500}}@{}}
\toprule\noalign{}
\begin{minipage}[b]{\linewidth}\raggedright
\textbf{Thuộc tính}
\end{minipage} & \begin{minipage}[b]{\linewidth}\raggedright
\textbf{Nội dung}
\end{minipage} \\
\midrule\noalign{}
\endhead
\bottomrule\noalign{}
\endlastfoot
\textbf{Mức ưu tiên} & MUST \\
\textbf{Trạng thái} & DERIVED \\
\textbf{Yêu cầu} & Mỗi Customer chỉ được có tối đa một Order chưa hoàn tất thanh toán tại một thời điểm, bao gồm trạng thái PENDING hoặc PAYMENT\_FAILED.

Khi Customer đã có Order thuộc một trong hai trạng thái trên, hệ thống từ chối yêu cầu tạo Order mới với mã lỗi 409 PENDING\_ORDER\_EXISTS và trả về mã đơn đang chờ.

Customer có thể tiếp tục thanh toán hoặc hủy Order hiện tại. Sau khi Order chuyển sang PAID hoặc CANCELLED, Customer được phép tạo Order mới.

Khi Customer yêu cầu tạo Order mới mà Order PENDING hoặc PAYMENT\_FAILED hiện tại đã quá hạn (expiresAt) nhưng chưa được Scheduler xử lý, backend phải chủ động chuyển Order cũ sang CANCELLED, hoàn tồn kho và giải phóng lượt coupon đã giữ đúng một lần trong transaction trước khi cho phép tạo Order mới. Nếu không thể hoàn tất việc hủy Order cũ, hệ thống từ chối tạo Order mới để bảo đảm tính nhất quán dữ liệu. \\
\textbf{Lý do/Nguồn} & Charter §6 (R4, R5), §3.1 (giỏ giữ khi lỗi); Techstack §7 (không tạo đơn kép). Quy tắc ``một đơn chờ'' là quyết định của nhóm. \\
\textbf{Actor chính} & Customer; Backend API \\
\end{longtable}

\textbf{Phân tích nghiệp vụ}

Vì giỏ vẫn được giữ sau khi tạo đơn, khách dễ bấm tạo đơn nhiều lần (bấm đúp, retry khi mạng chậm, quay lại giỏ rồi đặt tiếp). Nếu cho phép, mỗi lần sẽ giữ thêm tồn kho. Giới hạn một đơn chờ mỗi khách chặn việc tạo trùng và làm lượng tồn kho bị giữ không vượt quá một giỏ mỗi khách.

\textbf{Điều kiện, dữ liệu và ràng buộc liên quan}

\begin{itemize}
\item
  Khi đã có đơn PENDING, khách thanh toán hoặc hủy đơn đó (hoặc chờ hết hạn, FR-ORD-007) trước khi tạo đơn mới.
\item
  Khi bị từ chối, UI đưa khách tới đơn đang chờ để thanh toán hoặc hủy; nếu phản hồi tạo đơn lần trước bị thất lạc thì khách tìm lại được đơn.
\item
  Sau khi Order chuyển sang PAID hoặc CANCELLED, Customer được phép tạo Order mới.
\item
  Nếu Order đang ở trạng thái PAYMENT\_FAILED, Customer không được tạo Order mới mà phải tiếp tục thanh toán trên cùng Order hoặc hủy Order hiện tại. Nếu Order đã hết hạn, hệ thống phải xử lý chuyển Order sang CANCELLED theo FR-ORD-007 trước khi cho phép tạo Order mới.
\item
  Bảo đảm ở tầng dữ liệu, không chỉ kiểm tra trước khi ghi.
\item
  (SHOULD) POST /orders nhận thêm Idempotency-Key tùy chọn: cùng key và cùng nội dung thì trả lại đơn đã tạo; cùng key khác nội dung bị từ chối.
\end{itemize}

\textbf{Vấn đề có thể gặp và phân tích tác động}

\begin{longtable}[]{@{}
  >{\raggedright\arraybackslash}p{(\columnwidth - 2\tabcolsep) * \real{0.2500}}
  >{\raggedright\arraybackslash}p{(\columnwidth - 2\tabcolsep) * \real{0.7500}}@{}}
\toprule\noalign{}
\begin{minipage}[b]{\linewidth}\raggedright
\textbf{Vấn đề}
\end{minipage} & \begin{minipage}[b]{\linewidth}\raggedright
\textbf{Tác động/Nguyên nhân}
\end{minipage} \\
\midrule\noalign{}
\endhead
\bottomrule\noalign{}
\endlastfoot
Hai request tạo đơn đồng thời của cùng khách & Kiểm tra rồi ghi cho cả hai qua, tạo hai đơn và giữ tồn kho hai lần. \\
Khách bị kẹt vì đơn PENDING bỏ dở & Không đặt được đơn mới (xử lý ở FR-ORD-007). \\
Lỗi 409 khó hiểu & Khách không biết phải làm gì. \\
\end{longtable}

\textbf{Hướng xử lý/giải pháp phân tích đề xuất}

\begin{itemize}
\item
  Dùng ràng buộc duy nhất có điều kiện ở tầng dữ liệu đối với Order thuộc trạng thái PENDING hoặc PAYMENT\_FAILED của cùng một Customer, hoặc sử dụng cơ chế khóa và transaction có mức cô lập phù hợp để bảo đảm không tạo trùng đơn khi các request đến đồng thời.
\item
  Phản hồi 409 kèm mã đơn đang chờ; UI hiển thị ``Bạn đang có đơn chờ thanh toán'' kèm nút thanh toán và hủy.
\item
  Vô hiệu nút ``Đặt hàng'' ngay khi bấm để giảm bấm đúp (chỉ phục vụ UX).
\end{itemize}

\textbf{Chuyển giao sang thiết kế}

\begin{itemize}
\item
  ERD Order (ràng buộc tối đa một Order thuộc trạng thái PENDING hoặc PAYMENT\_FAILED trên mỗi Customer); Sequence tạo đơn; cơ chế transaction và kiểm soát đồng thời; integration test chống tạo đơn trùng; UI thông báo và điều hướng đến đơn đang chờ thanh toán.
\end{itemize}

\textbf{Tiêu chí kiểm chứng ở mức yêu cầu}

\begin{itemize}
\item
  Customer đang có Order thuộc trạng thái PENDING hoặc PAYMENT\_FAILED khi tạo Order mới phải nhận 409 PENDING\_ORDER\_EXISTS kèm mã Order đang chờ.
\item
  Hai request POST /orders đồng thời của cùng một Customer chỉ được tạo một Order và chỉ giữ tồn kho, lượt mã giảm giá một lần.
\item
  Gửi hai POST /orders đồng thời cho cùng một khách chỉ tạo đúng một đơn; tồn kho chỉ bị trừ một lần.
\item
  Sau khi đơn chuyển PAID hoặc CANCELLED, khách được tạo đơn mới. Đơn PAYMENT\_FAILED chưa hết hạn vẫn được xem là đơn đang chờ thanh toán.
\item
  Hai khách khác nhau tạo đơn đồng thời không ảnh hưởng nhau.
\end{itemize}

\textbf{FR-ORD-007 - Tự hủy đơn chờ thanh toán quá hạn}

\begin{longtable}[]{@{}
  >{\raggedright\arraybackslash}p{(\columnwidth - 2\tabcolsep) * \real{0.2500}}
  >{\raggedright\arraybackslash}p{(\columnwidth - 2\tabcolsep) * \real{0.7500}}@{}}
\toprule\noalign{}
\begin{minipage}[b]{\linewidth}\raggedright
\textbf{Thuộc tính}
\end{minipage} & \begin{minipage}[b]{\linewidth}\raggedright
\textbf{Nội dung}
\end{minipage} \\
\midrule\noalign{}
\endhead
\bottomrule\noalign{}
\endlastfoot
\textbf{Mức ưu tiên} & SHOULD \\
\textbf{Trạng thái} & DERIVED \\
\textbf{Yêu cầu} & Đơn hàng ở trạng thái PENDING hoặc PAYMENT\_FAILED hết hiệu lực thanh toán sau 15 phút kể từ thời điểm tạo đơn.

Khi đơn hết hạn, hệ thống tự động chuyển đơn sang CANCELLED, hoàn tồn kho và trả lượt sử dụng mã giảm giá đã giữ đúng một lần.

Backend phải từ chối mọi lần thanh toán mới đối với đơn đã hết hạn, kể cả khi tác vụ tự hủy chưa thực hiện. \\
\textbf{Lý do/Nguồn} & Charter §6 (R5, R6); suy ra từ cơ chế giữ tồn kho khi tạo đơn. Mức thời hạn là quyết định của nhóm. \\
\textbf{Actor chính} & Hệ thống (tác vụ định kỳ) \\
\end{longtable}

\textbf{Phân tích nghiệp vụ}

Đơn hàng PENDING và PAYMENT\_FAILED đều có thể đang giữ tồn kho và lượt sử dụng mã giảm giá. Nếu Customer không tiếp tục thanh toán hoặc không chủ động hủy đơn, tài nguyên có thể bị giữ vô thời hạn.

BrewLite quy định thời hạn thanh toán là 15 phút kể từ khi tạo đơn. Quá thời hạn này, Customer không được tiếp tục thanh toán hoặc thử lại thanh toán.

Tác vụ định kỳ có trách nhiệm chuyển những đơn quá hạn sang CANCELLED và giải phóng tài nguyên. Ngoài ra, API thanh toán phải tự kiểm tra thời hạn để ngăn trường hợp tác vụ chưa kịp chạy.

\textbf{Điều kiện, dữ liệu và ràng buộc liên quan}

\begin{itemize}
\item
  Thời hạn thanh toán là 15 phút kể từ thời điểm tạo đơn, được lưu trong expiresAt.
\item
  expiresAt được xác lập khi tạo Order và không thay đổi khi thử lại thanh toán.
\item
  Tác vụ định kỳ thực hiện quét mỗi phút.
\item
  Chỉ tự động hủy Order ở trạng thái PENDING hoặc PAYMENT\_FAILED đã hết hạn.
\item
  API POST /payments phải kiểm tra expiresAt đối với mọi lần thanh toán mới.
\item
  Việc chuyển CANCELLED, hoàn tồn kho và trả lượt mã giảm giá phải thực hiện đúng một lần, bảo đảm tính nguyên tử.
\item
  Đơn đã PAID hoặc chuyển sang các trạng thái xử lý tiếp theo không được tự hủy bởi tác vụ hết hạn.
\item
  Việc xử lý đồng thời giữa thanh toán, hủy đơn và tác vụ hết hạn phải sử dụng khóa hoặc cập nhật có điều kiện để tránh sai lệch dữ liệu.
\end{itemize}

\textbf{Vấn đề có thể gặp và phân tích tác động}

\begin{longtable}[]{@{}
  >{\raggedright\arraybackslash}p{(\columnwidth - 2\tabcolsep) * \real{0.2500}}
  >{\raggedright\arraybackslash}p{(\columnwidth - 2\tabcolsep) * \real{0.7500}}@{}}
\toprule\noalign{}
\begin{minipage}[b]{\linewidth}\raggedright
\textbf{Vấn đề}
\end{minipage} & \begin{minipage}[b]{\linewidth}\raggedright
\textbf{Tác động/Nguyên nhân}
\end{minipage} \\
\midrule\noalign{}
\endhead
\bottomrule\noalign{}
\endlastfoot
Hết hạn đúng lúc khách đang thanh toán & Khách trả tiền cho đơn đã hủy (xử lý bằng cập nhật có điều kiện). \\
Tác vụ không chạy & Tồn kho bị giữ lâu, khách bị kẹt. \\
Tác vụ chạy trùng ở nhiều tiến trình & Hoàn tồn kho hai lần. \\
\end{longtable}

\textbf{Hướng xử lý/giải pháp phân tích đề xuất}

\begin{itemize}
\item
  Tác vụ định kỳ của NestJS (Schedule), mỗi đơn hủy bằng cập nhật có điều kiện theo trạng thái.
\item
  Nếu chưa làm được, ghi nhận hạn chế và khách vẫn tự hủy được bằng FR-ORD-005.
\end{itemize}

\textbf{Chuyển giao sang thiết kế}

\begin{itemize}
\item
  Tác vụ định kỳ; cấu hình thời hạn; Sequence hết hạn; test với thời hạn rút ngắn.
\item
  Tác vụ định kỳ (Cron) chạy mỗi phút, truy vấn các Order có trạng thái thuộc \{PENDING, PAYMENT\_FAILED\} và expiresAt nhỏ hơn hoặc bằng thời điểm hiện tại.
\item
  Truy vấn phải giới hạn số lượng Order trong mỗi batch để tránh gây quá tải hệ thống.
\item
  Mỗi Order quá hạn được chuyển sang CANCELLED bằng cập nhật có điều kiện. Việc hoàn tồn kho và trả lượt sử dụng mã giảm giá phải thực hiện đúng một lần trong cùng transaction.
\item
  Các Order đã PAID, PREPARING, READY, COMPLETED hoặc CANCELLED không được tác vụ này thay đổi.
\end{itemize}

\textbf{Tiêu chí kiểm chứng ở mức yêu cầu}

\begin{itemize}
\item
  Với thời hạn rút ngắn khi kiểm thử, Order ở trạng thái PENDING hoặc PAYMENT\_FAILED khi hết hạn phải chuyển sang CANCELLED; tồn kho và lượt sử dụng mã giảm giá đã giữ được hoàn trả đúng một lần.
\item
  Đơn đã PAID không bị tác vụ hủy.
\item
  Chạy tác vụ hai lần liên tiếp không hoàn tồn kho hai lần.
\end{itemize}

\textbf{7.6. Inventory}

\textbf{Actor/phạm vi chính:} Customer; Backend API (module Order/Product); hệ thống (tác vụ định kỳ)

\textbf{Artifact thiết kế chịu ảnh hưởng:} ERD (Product.stock), Sequence tạo đơn, hủy đơn và hết hạn, State Diagram Order, concurrency test, Test Strategy

\textbf{FR-INV-001 - Quản lý tồn kho theo sản phẩm}

\begin{longtable}[]{@{}
  >{\raggedright\arraybackslash}p{(\columnwidth - 2\tabcolsep) * \real{0.2500}}
  >{\raggedright\arraybackslash}p{(\columnwidth - 2\tabcolsep) * \real{0.7500}}@{}}
\toprule\noalign{}
\begin{minipage}[b]{\linewidth}\raggedright
\textbf{Thuộc tính}
\end{minipage} & \begin{minipage}[b]{\linewidth}\raggedright
\textbf{Nội dung}
\end{minipage} \\
\midrule\noalign{}
\endhead
\bottomrule\noalign{}
\endlastfoot
\textbf{Mức ưu tiên} & MUST \\
\textbf{Trạng thái} & \emph{CONFIRMED (tồn kho theo sản phẩm, không âm); DERIVED (đơn vị là ly, trừ khi tạo đơn)}. \\
\textbf{Yêu cầu} & Hệ thống quản lý số lượng tồn kho khả dụng (stock) ở cấp độ Sản phẩm (Product), là số nguyên không âm, tính theo số ly. Bất kể khách hàng chọn size nào hay topping nào, tồn kho của sản phẩm đó bị trừ đi đúng số lượng ly (qty) được đặt. Tồn kho được trừ khi Order được tạo thành công (FR-ORD-001, FR-INV-002) để giữ số lượng sản phẩm đã đặt.

Khi Order ở trạng thái PENDING hoặc PAYMENT\_FAILED, tồn kho tiếp tục được giữ. Tồn kho chỉ được hoàn lại khi Order chuyển sang CANCELLED (FR-INV-003).

Việc thử lại thanh toán trên cùng Order không làm thay đổi tồn kho. Tồn kho khả dụng không bao giờ được nhỏ hơn 0. \\
\textbf{Lý do/Nguồn} & Charter §3.1 (cơ chế tồn kho phục vụ Task 10), §3.2 (không có ERP/POS quản lý nguyên liệu); Techstack §6 (Product có stock); Task 10. \\
\textbf{Actor chính} & Backend API \\
\end{longtable}

\textbf{Phân tích nghiệp vụ}

BrewLite chỉ cần cơ chế tồn kho đủ để chống bán vượt số lượng, không quản lý nguyên liệu. Vì vậy tồn kho đếm theo \textbf{số ly của sản phẩm}, không tách theo size hay topping. Mỗi ly bán ra, bất kể size nào, trừ một đơn vị của sản phẩm đó.

\textbf{Điều kiện, dữ liệu và ràng buộc liên quan}

\begin{itemize}
\item
  Tồn kho là số nguyên lớn hơn hoặc bằng 0.
\item
  Sản phẩm có stock = 0 là hết hàng (FR-CAT-004).
\item
  Không có giao diện điều chỉnh tồn kho; giá trị ban đầu do dữ liệu mẫu (FR-CAT-005) và có thể được đặt lại bằng script cho mục đích kiểm thử.
\item
  Tồn kho chỉ thay đổi bởi các nghiệp vụ hợp lệ: khởi tạo hoặc điều chỉnh dữ liệu mẫu, tạo Order thành công (trừ tồn kho để giữ hàng), và Order chuyển sang CANCELLED (hoàn tồn kho).
\item
  Order chuyển sang PAYMENT\_FAILED hoặc được thử thanh toán lại không làm thay đổi tồn kho.
\item
  Ràng buộc ``không âm'' phải được bảo đảm ở \textbf{tầng dữ liệu} (ràng buộc kiểm tra hoặc điều kiện cập nhật), không chỉ ở tầng ứng dụng.
\item
  API công khai chỉ cho biết còn hay hết hàng, không trả số lượng chính xác (FR-CAT-001, FR-CAT-004). Ngoại lệ là thông tin trong lỗi INSUFFICIENT\_STOCK ở FR-INV-002.
\end{itemize}

\textbf{Vấn đề có thể gặp và phân tích tác động}

\begin{longtable}[]{@{}
  >{\raggedright\arraybackslash}p{(\columnwidth - 2\tabcolsep) * \real{0.2500}}
  >{\raggedright\arraybackslash}p{(\columnwidth - 2\tabcolsep) * \real{0.7500}}@{}}
\toprule\noalign{}
\begin{minipage}[b]{\linewidth}\raggedright
\textbf{Vấn đề}
\end{minipage} & \begin{minipage}[b]{\linewidth}\raggedright
\textbf{Tác động/Nguyên nhân}
\end{minipage} \\
\midrule\noalign{}
\endhead
\bottomrule\noalign{}
\endlastfoot
Tồn kho âm & Dữ liệu sai, hệ thống đã bán vượt số lượng. \\
Trừ hoặc hoàn tồn không qua một điểm kiểm soát chung & Mỗi nơi tự cập nhật theo kiểu riêng, dễ lệch. \\
Lộ số lượng tồn chính xác cho mọi người & Người ngoài biết tình trạng kho mà không cần thiết. \\
Giá trị tồn do dữ liệu mẫu đặt sai & Không kiểm thử được hết hàng, tồn kho thấp. \\
\end{longtable}

\textbf{Hướng xử lý/giải pháp phân tích đề xuất}

\begin{itemize}
\item
  Tập trung mọi thay đổi tồn kho vào một thành phần dịch vụ duy nhất (trừ và hoàn); không cập nhật tồn kho ở nơi khác.
\item
  Thêm ràng buộc stock \textgreater= 0 ở cơ sở dữ liệu làm lớp bảo vệ cuối cùng.
\item
  DTO danh sách và chi tiết sản phẩm chỉ trả trạng thái còn hoặc hết.
\item
  Dữ liệu mẫu có món tồn kho 0 và món tồn kho thấp (FR-CAT-005).
\end{itemize}

\textbf{Chuyển giao sang thiết kế}

\begin{itemize}
\item
  ERD Product (stock, ràng buộc không âm); module tồn kho; dữ liệu mẫu cho test; unit test dịch vụ tồn kho.
\end{itemize}

\textbf{Tiêu chí kiểm chứng ở mức yêu cầu}

\begin{itemize}
\item
  Cập nhật trực tiếp làm stock âm trong cơ sở dữ liệu bị từ chối bởi ràng buộc; cập nhật có điều kiện của FR-INV-002 không bao giờ làm stock âm.
\item
  API danh sách và chi tiết không trả số lượng tồn chính xác.
\item
  Tồn kho của một sản phẩm giảm đúng tổng số ly của các dòng cùng sản phẩm, bất kể size và topping.
\end{itemize}

\textbf{FR-INV-002 - Chống bán vượt tồn kho khi đặt đồng thời}

\begin{longtable}[]{@{}
  >{\raggedright\arraybackslash}p{(\columnwidth - 2\tabcolsep) * \real{0.2500}}
  >{\raggedright\arraybackslash}p{(\columnwidth - 2\tabcolsep) * \real{0.7500}}@{}}
\toprule\noalign{}
\begin{minipage}[b]{\linewidth}\raggedright
\textbf{Thuộc tính}
\end{minipage} & \begin{minipage}[b]{\linewidth}\raggedright
\textbf{Nội dung}
\end{minipage} \\
\midrule\noalign{}
\endhead
\bottomrule\noalign{}
\endlastfoot
\textbf{Mức ưu tiên} & MUST \\
\textbf{Trạng thái} & CONFIRMED (chống oversell, concurrency); DERIVED (gộp theo sản phẩm, nội dung lỗi) \\
\textbf{Yêu cầu} & Khi tạo đơn, backend phải trừ tồn kho theo cách bảo đảm tổng số lượng bán ra không bao giờ vượt tồn kho, kể cả khi nhiều khách đặt đồng thời. Nếu một sản phẩm không đủ tồn, toàn bộ đơn bị từ chối với 409 INSUFFICIENT\_STOCK và không có thay đổi nào được lưu. \\
\textbf{Lý do/Nguồn} & Charter §3.1 (chống oversell khi đặt đồng thời), §6 (R5), §7 (concurrency test không vượt tồn kho); Techstack §7 (Concurrency); Task 10. \\
\textbf{Actor chính} & Customer; Backend API \\
\end{longtable}

\textbf{Phân tích nghiệp vụ}

Oversell xảy ra khi hai yêu cầu cùng đọc ``còn 1'' rồi cùng ghi ``đã bán'': kiểm tra rồi mới trừ không an toàn khi đồng thời. Giải pháp phải làm cho bước kiểm tra và bước trừ thành \textbf{một thao tác không thể chen ngang}.

Hai điểm nghiệp vụ cần chốt rõ:

\begin{itemize}
\item
  Số lượng cần so với tồn kho là \textbf{tổng các dòng cùng sản phẩm} (ví dụ cùng món, một dòng size S và một dòng size L), vì chúng dùng chung một kho.
\item
  Đơn chỉ có hiệu lực khi \textbf{mọi} sản phẩm trong đơn đều đủ tồn (tất cả hoặc không gì cả).
\end{itemize}

\textbf{Điều kiện, dữ liệu và ràng buộc liên quan}

\begin{itemize}
\item
  Cơ chế nền tảng: cập nhật có điều kiện theo kiểu ``trừ qty nếu stock \textgreater= qty'' và kiểm tra có đúng một bản ghi được cập nhật (cập nhật nguyên tử). Cách dùng cột phiên bản (optimistic locking) là phương án tương đương nếu nhóm chọn (Techstack §7).
\item
  Toàn bộ việc trừ tồn của mọi sản phẩm trong đơn và lưu đơn nằm trong \textbf{một transaction} (FR-ORD-001); thất bại ở sản phẩm nào cũng hoàn tác phần đã trừ.
\item
  Các sản phẩm được trừ theo một thứ tự cố định (ví dụ theo mã sản phẩm tăng dần) để hai đơn chứa cùng các sản phẩm không khóa chéo nhau (deadlock).
\item
  Lỗi 409 INSUFFICIENT\_STOCK chỉ ra tối thiểu một sản phẩm thiếu hàng (mục tiêu SHOULD: chỉ ra tất cả). Với mỗi sản phẩm thiếu hàng, phản hồi cho biết tổng số ly khách yêu cầu, \textbf{số ly còn có thể đặt} (nhỏ hơn số yêu cầu) và các dòng giỏ liên quan; việc chia số còn lại cho dòng nào do khách tự chỉnh giỏ. Số liệu này là hạn chế đã biết về lộ thông tin tồn kho, và chỉ xuất hiện trong lỗi.
\item
  Thứ tự kiểm tra tổng thể: xác thực (401) → dữ liệu giỏ (400) → đã có đơn PENDING (409 PENDING\_ORDER\_EXISTS) → tồn kho (409 INSUFFICIENT\_STOCK) (FR-ORD-001).
\item
  Kiểm tra tồn kho của frontend (FR-CAT-004) chỉ là thông tin tham khảo; quyết định cuối cùng là cập nhật có điều kiện này.
\item
  Nếu dùng cột phiên bản, phải thử lại có giới hạn khi xung đột phiên bản, và chỉ trả INSUFFICIENT\_STOCK khi tồn kho thật sự không đủ. Phương án cập nhật nguyên tử có điều kiện được ưu tiên vì không phát sinh lỗi xung đột giả.
\end{itemize}

\textbf{Vấn đề có thể gặp và phân tích tác động}

\begin{longtable}[]{@{}
  >{\raggedright\arraybackslash}p{(\columnwidth - 2\tabcolsep) * \real{0.2500}}
  >{\raggedright\arraybackslash}p{(\columnwidth - 2\tabcolsep) * \real{0.7500}}@{}}
\toprule\noalign{}
\begin{minipage}[b]{\linewidth}\raggedright
\textbf{Vấn đề}
\end{minipage} & \begin{minipage}[b]{\linewidth}\raggedright
\textbf{Tác động/Nguyên nhân}
\end{minipage} \\
\midrule\noalign{}
\endhead
\bottomrule\noalign{}
\endlastfoot
Đọc tồn kho, so sánh ở tầng ứng dụng rồi ghi lại (kiểm tra rồi mới ghi) & Hai request đồng thời cùng qua, bán vượt số lượng. \\
Chỉ so sánh từng dòng riêng với tồn kho & Hai dòng cùng sản phẩm mỗi dòng đều đủ nhưng tổng vượt tồn. \\
Trừ tồn thành công nhưng lưu đơn thất bại & Mất hàng không có đơn. \\
Trừ tồn theo thứ tự khác nhau giữa các đơn & Hai transaction chờ nhau (deadlock), một đơn lỗi khó giải thích. \\
Thông báo lỗi không nói dòng nào thiếu & Khách không biết sửa gì. \\
Lộ quá nhiều thông tin tồn kho trong lỗi & Người khác dò được số lượng tồn kho thật. \\
\end{longtable}

\textbf{Hướng xử lý/giải pháp phân tích đề xuất}

\begin{itemize}
\item
  Gộp số lượng các dòng theo sản phẩm trước khi trừ; mỗi sản phẩm thực hiện một cập nhật có điều kiện duy nhất.
\item
  Thực hiện việc trừ và lưu đơn trong cùng transaction; khi một sản phẩm không đủ thì hủy toàn bộ transaction.
\item
  Sắp xếp các sản phẩm theo mã trước khi trừ để tránh deadlock.
\item
  Sau khi hủy, đọc lại số lượng còn của các dòng thiếu để dựng phản hồi lỗi (tối thiểu nêu dòng đầu tiên thiếu hàng; mục tiêu là nêu tất cả).
\item
  Frontend khi nhận lỗi giữ nguyên giỏ, đánh dấu dòng thiếu và gợi ý số lượng còn có thể đặt.
\item
  Ràng buộc stock \textgreater= 0 ở cơ sở dữ liệu là lớp bảo vệ thứ hai.
\end{itemize}

\textbf{Chuyển giao sang thiết kế}

\begin{itemize}
\item
  Sequence tạo đơn (gộp số lượng → trừ có điều kiện theo thứ tự → lưu đơn); ERD Product (stock, tùy chọn cột phiên bản); định dạng lỗi INSUFFICIENT\_STOCK; UI báo thiếu hàng theo dòng; concurrency test.
\item
  Hàm tạo đơn phải gom nhóm (group by productId) và cộng dồn qty của các OrderItem trước khi thực thi vòng lặp UPDATE trừ tồn kho.
\end{itemize}

\textbf{Tiêu chí kiểm chứng ở mức yêu cầu}

\begin{itemize}
\item
  Sản phẩm tồn kho 5; 10 khách khác nhau cùng lúc mỗi người đặt 1 ly: đúng 5 đơn thành công, 5 đơn nhận 409 INSUFFICIENT\_STOCK, tồn kho cuối bằng 0 (không âm).
\item
  Sản phẩm tồn kho 5; một đơn có hai dòng cùng sản phẩm (size S × 3 và size L × 3) bị từ chối 409 dù mỗi dòng riêng đều không vượt 5; không có đơn nào được lưu.
\item
  Đơn có hai sản phẩm A và B, A đủ tồn nhưng B thiếu: đơn bị từ chối và tồn kho của A không đổi.
\item
  Hai đơn cùng chứa A và B nhưng khai báo theo thứ tự ngược nhau, gửi đồng thời nhiều lần: không có đơn nào lỗi do deadlock; kết quả đúng với tồn kho.
\item
  Phản hồi lỗi chỉ ra đúng sản phẩm thiếu hàng, tổng số yêu cầu và số còn có thể đặt; khách đặt đúng số lượng đó thì thành công.
\item
  Gọi trực tiếp API (không qua UI) với số lượng vượt tồn kho vẫn bị từ chối.
\end{itemize}

\textbf{FR-INV-003 - Hoàn tồn kho khi đơn không thành công}

\begin{longtable}[]{@{}
  >{\raggedright\arraybackslash}p{(\columnwidth - 2\tabcolsep) * \real{0.2500}}
  >{\raggedright\arraybackslash}p{(\columnwidth - 2\tabcolsep) * \real{0.7500}}@{}}
\toprule\noalign{}
\begin{minipage}[b]{\linewidth}\raggedright
\textbf{Thuộc tính}
\end{minipage} & \begin{minipage}[b]{\linewidth}\raggedright
\textbf{Nội dung}
\end{minipage} \\
\midrule\noalign{}
\endhead
\bottomrule\noalign{}
\endlastfoot
\textbf{Mức ưu tiên} & MUST \\
\textbf{Trạng thái} & DERIVED \\
\textbf{Yêu cầu} & Khi Order chuyển từ PENDING hoặc PAYMENT\_FAILED sang CANCELLED, hệ thống phải hoàn tồn kho đúng bằng số lượng sản phẩm đã giữ khi tạo Order.

Việc hoàn tồn kho chỉ được thực hiện một lần, trong cùng transaction với thao tác chuyển trạng thái Order sang CANCELLED.

Khi Order chuyển sang PAYMENT\_FAILED hoặc từ PAYMENT\_FAILED về PENDING để thử lại thanh toán, tồn kho tiếp tục được giữ và không thực hiện hoàn hoặc trừ thêm. \\
\textbf{Lý do/Nguồn} & Techstack §7; Charter §3.1 (PAYMENT\_FAILED), §6 (R5). Việc trừ tồn khi tạo đơn và hoàn khi thất bại là quyết định của nhóm. \\
\textbf{Actor chính} & Backend API; hệ thống (tác vụ định kỳ) \\
\end{longtable}

\textbf{Phân tích nghiệp vụ}

Vì tồn kho bị trừ ngay khi tạo đơn (giữ chỗ), đơn không hoàn tất phải trả lại số lượng đã giữ, nếu không hàng bị ``mất'' vĩnh viễn và menu hiển thị hết hàng sai. Ngược lại, hoàn hai lần làm tồn kho cao hơn thực tế và dẫn đến bán vượt số lượng thật. Do đó thao tác hoàn phải gắn chặt với việc chuyển trạng thái thành công.

\textbf{Điều kiện, dữ liệu và ràng buộc liên quan}

\begin{itemize}
\item
  Số lượng hoàn lấy từ dòng hàng đã lưu của đơn (snapshot, FR-ORD-002), gộp theo sản phẩm; không lấy từ giỏ ở client.
\item
  Hoàn tồn chỉ xảy ra khi chuyển trạng thái thật sự thành công (cập nhật có điều kiện theo trạng thái xuất phát, FR-ORD-003); thao tác thứ hai nhận 409 và không hoàn lại.
\item
  Có hai nguồn kích hoạt hoàn tồn kho: Customer hủy Order hợp lệ (FR-ORD-005) và hệ thống tự hủy Order quá hạn (FR-ORD-007).
\item
  Cả hai trường hợp chỉ hoàn tồn kho khi Order thực sự chuyển sang CANCELLED. Payment FAILED không phải sự kiện hoàn tồn kho.
\item
  Đơn PAID trở đi không hoàn tồn; không có chuyển PAID → CANCELLED.
\item
  Mọi cập nhật tồn kho đi cùng việc đổi trạng thái đơn trong một transaction.
\end{itemize}

\textbf{Vấn đề có thể gặp và phân tích tác động}

\begin{longtable}[]{@{}
  >{\raggedright\arraybackslash}p{(\columnwidth - 2\tabcolsep) * \real{0.2500}}
  >{\raggedright\arraybackslash}p{(\columnwidth - 2\tabcolsep) * \real{0.7500}}@{}}
\toprule\noalign{}
\begin{minipage}[b]{\linewidth}\raggedright
\textbf{Vấn đề}
\end{minipage} & \begin{minipage}[b]{\linewidth}\raggedright
\textbf{Tác động/Nguyên nhân}
\end{minipage} \\
\midrule\noalign{}
\endhead
\bottomrule\noalign{}
\endlastfoot
Không hoàn tồn kho khi Order bị hủy hoặc hết hạn & Tồn kho bị giữ dù đơn đã CANCELLED, làm giảm tồn kho khả dụng và ảnh hưởng các khách hàng khác. \\
Hoàn tồn kho hai lần khi hủy đơn hoặc tác vụ hết hạn chạy trùng & Tồn kho tăng sai, dẫn đến nguy cơ bán vượt số lượng thực tế. \\
Hoàn bằng số lượng lấy từ client & Số hoàn sai nếu client gửi sai. \\
Hủy và thanh toán xảy ra đồng thời & Vừa hoàn tồn vừa ghi nhận thanh toán thành công (xử lý ở FR-ORD-003). \\
Sản phẩm không còn trong catalog khi hoàn & Không có chỗ để cộng lại tồn kho. \\
Đơn PENDING bỏ dở giữ chỗ tồn kho & Khách khác thấy món hết hàng dù chưa ai trả tiền. Giới hạn bởi thời hạn 15 phút của FR-ORD-007. \\
\end{longtable}

\textbf{Hướng xử lý/giải pháp phân tích đề xuất}

\begin{itemize}
\item
  Gắn hoàn tồn vào hàm chuyển trạng thái tương ứng, chạy chỉ khi cập nhật có điều kiện thành công.
\item
  Lấy số lượng hoàn từ dòng hàng đã lưu, gộp theo sản phẩm, và cộng lại cho từng sản phẩm.
\item
  MVP không xóa sản phẩm, nên hoàn tồn luôn tìm thấy sản phẩm; ghi nhận giả định này trong tài liệu.
\item
  Kiểm thử chạy lại nhiều lần mỗi nguồn kích hoạt để chứng minh ``đúng một lần''.
\end{itemize}

\textbf{Chuyển giao sang thiết kế}

\begin{itemize}
\item
  Sequence thanh toán thất bại, hủy đơn và hết hạn (bước hoàn tồn); ERD (không xóa sản phẩm khỏi catalog, hoặc ràng buộc khóa ngoại); test từng nguồn kích hoạt và test chạy lặp.
\end{itemize}

\textbf{Tiêu chí kiểm chứng ở mức yêu cầu}

\begin{itemize}
\item
  Tạo đơn rồi hủy: tồn kho trở lại đúng giá trị trước khi tạo đơn.
\item
  Tạo đơn rồi thanh toán thất bại: Order chuyển sang PAYMENT\_FAILED, tồn kho tiếp tục được giữ và không được hoàn lại. Customer có thể thử thanh toán lại trên cùng Order nếu chưa hết hạn.
\item
  Thử lại thanh toán trên cùng Order (PAYMENT\_FAILED → PENDING): tồn kho không bị trừ thêm hoặc hoàn lại; số lượng tồn kho giữ nguyên so với sau khi tạo đơn.
\item
  Đơn quá hạn bị tác vụ hủy: tồn kho trở lại đúng giá trị trước khi tạo đơn.
\item
  Bấm hủy hai lần liên tiếp, hoặc chạy tác vụ hết hạn hai lần: tồn kho chỉ được hoàn một lần.
\item
  Hủy và thanh toán đồng thời trên cùng một đơn: tồn kho cuối cùng nhất quán với trạng thái cuối của đơn (đơn CANCELLED thì đã hoàn một lần; đơn PAID thì không hoàn).
\item
  Đơn PAID, PREPARING, READY, COMPLETED không bao giờ làm tồn kho tăng.
\end{itemize}

\textbf{FR-INV-004 - Đối soát biến động tồn kho}

\begin{longtable}[]{@{}
  >{\raggedright\arraybackslash}p{(\columnwidth - 2\tabcolsep) * \real{0.2500}}
  >{\raggedright\arraybackslash}p{(\columnwidth - 2\tabcolsep) * \real{0.7500}}@{}}
\toprule\noalign{}
\begin{minipage}[b]{\linewidth}\raggedright
\textbf{Thuộc tính}
\end{minipage} & \begin{minipage}[b]{\linewidth}\raggedright
\textbf{Nội dung}
\end{minipage} \\
\midrule\noalign{}
\endhead
\bottomrule\noalign{}
\endlastfoot
\textbf{Mức ưu tiên} & SHOULD \\
\textbf{Trạng thái} & DERIVED \\
\textbf{Yêu cầu} & Mỗi lần tồn kho thay đổi do đơn hàng, hệ thống ghi lại một bản ghi gồm sản phẩm, mã đơn, loại biến động (giữ chỗ hoặc hoàn), lý do (đặt hàng, thanh toán thất bại, hủy, quá hạn), số lượng thay đổi và thời điểm. \\
\textbf{Lý do/Nguồn} & Phục vụ kiểm chứng ``đúng một lần'' của FR-INV-003 và Task 10 (kiểm thử); quyết định của nhóm. \\
\textbf{Actor chính} & Backend API \\
\end{longtable}

\textbf{Phân tích nghiệp vụ}

Khi có tranh chấp ``tồn kho lệch'', chỉ nhìn con số hiện tại không cho biết lệch do đâu. Bản ghi biến động cho phép đối soát: tồn kho hiện tại phải bằng tồn ban đầu cộng tổng các biến động. Đây cũng là bằng chứng rõ ràng cho giảng viên rằng hoàn tồn không bị trùng.

\textbf{Điều kiện, dữ liệu và ràng buộc liên quan}

\begin{itemize}
\item
  Ghi cùng transaction với việc đổi tồn kho; không có biến động mà không có bản ghi và ngược lại.
\item
  Bản ghi không bị sửa hoặc xóa qua API.
\item
  Mỗi bộ ba (mã đơn, mã sản phẩm, loại biến động) có tối đa một bản ghi, bảo đảm ở tầng dữ liệu; nhờ vậy mỗi đơn chỉ có thể giữ chỗ một lần và hoàn một lần cho mỗi sản phẩm.
\item
  Chỉ phục vụ nội bộ; không có API công khai.
\end{itemize}

\textbf{Vấn đề có thể gặp và phân tích tác động}

\begin{longtable}[]{@{}
  >{\raggedright\arraybackslash}p{(\columnwidth - 2\tabcolsep) * \real{0.2500}}
  >{\raggedright\arraybackslash}p{(\columnwidth - 2\tabcolsep) * \real{0.7500}}@{}}
\toprule\noalign{}
\begin{minipage}[b]{\linewidth}\raggedright
\textbf{Vấn đề}
\end{minipage} & \begin{minipage}[b]{\linewidth}\raggedright
\textbf{Tác động/Nguyên nhân}
\end{minipage} \\
\midrule\noalign{}
\endhead
\bottomrule\noalign{}
\endlastfoot
Ghi biến động tách khỏi transaction & Số tồn và nhật ký lệch nhau. \\
Thiếu ràng buộc Unique theo (đơn, sản phẩm, loại biến động) & Nhật ký ghi trùng, che mất lỗi hoàn hai lần. \\
Thêm bảng làm tăng khối lượng việc & Rủi ro scope creep (Charter R1). \\
\end{longtable}

\textbf{Hướng xử lý/giải pháp}

\begin{itemize}
\item
  Một bảng biến động đơn giản; chỉ làm khi hoàn thành các FR MUST.
\item
  Dùng làm công cụ kiểm thử: sau mỗi bài test đồng thời, đối chiếu tồn kho với tổng biến động.
\end{itemize}

\textbf{Chuyển giao sang thiết kế}

\begin{itemize}
\item
  ERD bảng biến động tồn kho; Sequence liên quan; test đối soát.
\end{itemize}

\textbf{Tiêu chí kiểm chứng ở mức yêu cầu}

\begin{itemize}
\item
  Sau chuỗi tạo đơn, hủy, thanh toán thất bại, tồn kho hiện tại bằng tồn ban đầu cộng tổng biến động.
\item
  Với mỗi sản phẩm của một đơn, có tối đa một bản ghi giữ chỗ và một bản ghi hoàn.
\end{itemize}

\textbf{7.7. Payment}

\textbf{Actor/phạm vi chính:} Customer; Backend API (module Payment); Mock Payment Service; các module Order, Inventory, Loyalty

\textbf{Artifact thiết kế chịu ảnh hưởng:} Use Case thanh toán và xác nhận đơn, ERD (Payment), State Diagram Order và Payment, Sequence thanh toán, API Payment, UI bước thanh toán và trang xác nhận, concurrency test, Test Strategy

\textbf{FR-PAY-001 - Thanh toán mock cho đơn chờ thanh toán}

\begin{longtable}[]{@{}
  >{\raggedright\arraybackslash}p{(\columnwidth - 2\tabcolsep) * \real{0.2500}}
  >{\raggedright\arraybackslash}p{(\columnwidth - 2\tabcolsep) * \real{0.7500}}@{}}
\toprule\noalign{}
\begin{minipage}[b]{\linewidth}\raggedright
\textbf{Thuộc tính}
\end{minipage} & \begin{minipage}[b]{\linewidth}\raggedright
\textbf{Nội dung}
\end{minipage} \\
\midrule\noalign{}
\endhead
\bottomrule\noalign{}
\endlastfoot
\textbf{Mức ưu tiên} & MUST (MOCK) \\
\textbf{Trạng thái} & CONFIRMED (thanh toán Ví/Thẻ bằng mock gateway); DERIVED (cấu trúc request, thứ tự kiểm tra) \\
\textbf{Yêu cầu} & Customer đã đăng nhập thanh toán đơn PENDING của chính mình bằng Ví hoặc Thẻ qua POST /payments. Request gồm mã đơn, phương thức, dữ liệu test của phương thức và header Idempotency-Key. Số tiền do server lấy từ đơn; kết quả do Mock Payment Service quyết định, không do client gửi. \\
\textbf{Lý do/Nguồn} & Charter §3.1 (Thanh toán: Ví/Thẻ bằng mock gateway), §7 (Thanh toán); Techstack §6 (POST /payments có Idempotency-Key); Task 8. \\
\textbf{Actor chính} & Customer; Backend API; Mock Payment Service \\
\end{longtable}

\textbf{Phân tích nghiệp vụ}

Thanh toán biến đơn PENDING thành PAID, là bước sai sót tốn kém nhất: trừ tiền giả lập hai lần, thanh toán đơn của người khác, hoặc tin kết quả do client gửi. Dù là mock, hành vi phải như thật: server quyết định số tiền, server quyết định kết quả, mỗi đơn được xử lý đúng một lần (mục 3).

\textbf{Điều kiện, dữ liệu và ràng buộc liên quan}

\begin{itemize}
\item
  Endpoint POST /payments yêu cầu JWT (FR-AUTHZ-001).
\item
  Payload gồm: mã đơn, phương thức (WALLET hoặc CARD), dữ liệu test của phương thức (FR-PAY-006). Không có số tiền và không có kết quả; trường lạ bị bỏ qua.
\item
  Header Idempotency-Key bắt buộc (FR-PAY-003); thiếu hoặc sai định dạng trả 400. Phương thức ngoài WALLET, CARD trả 400.
\item
  Thứ tự kiểm tra: xác thực (401) → dữ liệu request (400) → đơn thuộc về người gọi (404, FR-AUTHZ-002) → tra cứu Idempotency-Key (FR-PAY-003) → kiểm tra Order còn hạn và ở trạng thái PENDING hoặc PAYMENT\_FAILED hợp lệ (409 nếu không đáp ứng) → xử lý thanh toán.
\item
  Số tiền thanh toán bằng tổng tiền của đơn đã lưu (FR-ORD-001).
\item
  Order đã PAID hoặc CANCELLED không được thực hiện lần thanh toán mới và phải bị từ chối với HTTP 409 INVALID\_TRANSITION.
\item
  Order đang PAYMENT\_FAILED và chưa hết hạn được phép thử lại thanh toán trên cùng Order bằng Idempotency-Key mới, sau khi chuyển về PENDING theo State Machine.
\item
  Order đã hết hạn không được bắt đầu lần thanh toán mới, dù tác vụ tự hủy chưa chạy.
\item
  Trường hợp phát lại cùng Idempotency-Key và cùng nội dung phải được xử lý theo FR-PAY-003 để trả kết quả đã lưu, không tạo Payment mới.
\item
  Mock Payment Service được gọi đồng bộ bên trong backend, phản hồi trong mục tiêu \textless{} 500 ms (Charter §7).
\item
  Phản hồi gồm: mã đơn, trạng thái đơn, mã giao dịch, trạng thái thanh toán, số tiền, phương thức (đã che thông tin), thời điểm. HTTP 200 cho cả thành công và thất bại nghiệp vụ (kết quả nằm ở trường trạng thái); lỗi request dùng 4xx.
\item
  Đơn ở trạng thái PAYMENT\_FAILED chưa hết hạn được phép thử lại thanh toán thông qua luồng chuyển trạng thái hợp lệ về PENDING.
\item
  Hệ thống phải từ chối thanh toán khi thời điểm hiện đã đạt hoặc vượt expiresAt, kể cả khi tác vụ tự hủy chưa chạy. Kết quả được phát lại bằng Idempotency-Key cũ phải tuần theo quy tắc Idempotency và không thực hiện giao dịch mới.
\end{itemize}

\textbf{Vấn đề có thể gặp và phân tích tác động}

\begin{longtable}[]{@{}
  >{\raggedright\arraybackslash}p{(\columnwidth - 2\tabcolsep) * \real{0.2500}}
  >{\raggedright\arraybackslash}p{(\columnwidth - 2\tabcolsep) * \real{0.7500}}@{}}
\toprule\noalign{}
\begin{minipage}[b]{\linewidth}\raggedright
\textbf{Vấn đề}
\end{minipage} & \begin{minipage}[b]{\linewidth}\raggedright
\textbf{Tác động/Nguyên nhân}
\end{minipage} \\
\midrule\noalign{}
\endhead
\bottomrule\noalign{}
\endlastfoot
Client gửi số tiền hoặc kết quả & Thanh toán rẻ hơn hoặc tự báo thành công. \\
Thanh toán đơn của người khác (IDOR) & Phá đơn hoặc dò mã đơn. \\
Thanh toán đơn đã hết hạn hoặc đã hủy & Đã ``trả tiền'' cho đơn không được xử lý, tồn kho đã hoàn. \\
Hết hạn đơn xảy ra đúng lúc đang thanh toán & Đơn vừa CANCELLED vừa PAID (FR-ORD-003). \\
Lỗi kỹ thuật giữa chừng & Payment đã ghi nhưng đơn chưa đổi trạng thái hoặc ngược lại. \\
Lộ hoặc ghi log dữ liệu thẻ & Rò thông tin nhạy cảm dù là dữ liệu test. \\
\end{longtable}

\textbf{Hướng xử lý/giải pháp phân tích đề xuất}

\begin{itemize}
\item
  DTO chỉ nhận các trường đã nêu; áp dụng đúng thứ tự kiểm tra ở trên.
\item
  Trong một transaction: xác nhận và khóa trạng thái đơn đang PENDING trước, rồi gọi Mock Payment, rồi ghi Payment và đổi trạng thái (FR-PAY-002).
\item
  Lỗi kỹ thuật (không phải kết quả nghiệp vụ): rollback toàn bộ transaction, không tạo Payment mới, không lưu Idempotency-Key của lần thử chưa hoàn tất, không thay đổi tồn kho, lượt mã giảm giá hoặc điểm tích lũy; trả HTTP 5xx.
\item
  Order giữ nguyên trạng thái trước lần thử: PENDING nếu thanh toán lần đầu, PAYMENT\_FAILED nếu đang thử lại sau thất bại.
\item
  Customer có thể gửi lại request với cùng Idempotency-Key sau khi lỗi kỹ thuật được khắc phục (BR-041).
\item
  Không ghi log dữ liệu thẻ hay ví; thông báo lỗi chung.
\item
  UI vô hiệu nút thanh toán ngay khi bấm (chỉ phục vụ UX).
\end{itemize}

\textbf{Chuyển giao sang thiết kế}

\begin{itemize}
\item
  Use Case Thanh toán; API POST /payments (DTO, header, mã lỗi); Sequence thanh toán; ERD Payment; UI bước thanh toán; integration test.
\end{itemize}

\textbf{Tiêu chí kiểm chứng ở mức yêu cầu}

\begin{itemize}
\item
  Customer thanh toán đơn PENDING hợp lệ nhận kết quả có số tiền bằng tổng đơn.
\item
  Gọi không token trả 401; thiếu Idempotency-Key hoặc phương thức lạ trả 400.
\item
  Thanh toán đơn của người khác trả 404, giống đơn không tồn tại.
\item
  Request có thêm trường số tiền hoặc kết quả vẫn dùng số tiền của đơn và kết quả của mock.
\item
  Order đã PAID hoặc CANCELLED phải từ chối yêu cầu thanh toán mới với lỗi 409 INVALID\_TRANSITION.
\item
  Order đang PAYMENT\_FAILED và chưa hết hạn được phép thử thanh toán lại trên cùng Order bằng Idempotency-Key mới.
\item
  Order ở trạng thái PENDING hoặc PAYMENT\_FAILED nhưng đã hết hạn phải từ chối lần thanh toán mới; request phát lại cùng Idempotency-Key và cùng nội dung vẫn trả kết quả đã lưu theo FR-PAY-003.
\item
  Khi cố ý gây lỗi kỹ thuật trong môi trường test: đơn vẫn PENDING, không có Payment, thử lại cùng khóa thành công.
\item
  Thời gian phản hồi với dữ liệu mẫu \textless{} 500 ms.
\end{itemize}

\textbf{FR-PAY-002 - Xử lý kết quả: PAID hoặc PAYMENT\_FAILED}

\begin{longtable}[]{@{}
  >{\raggedright\arraybackslash}p{(\columnwidth - 2\tabcolsep) * \real{0.2500}}
  >{\raggedright\arraybackslash}p{(\columnwidth - 2\tabcolsep) * \real{0.7500}}@{}}
\toprule\noalign{}
\begin{minipage}[b]{\linewidth}\raggedright
\textbf{Thuộc tính}
\end{minipage} & \begin{minipage}[b]{\linewidth}\raggedright
\textbf{Nội dung}
\end{minipage} \\
\midrule\noalign{}
\endhead
\bottomrule\noalign{}
\endlastfoot
\textbf{Mức ưu tiên} & MUST \\
\textbf{Trạng thái} & CONFIRMED \\
\textbf{Yêu cầu} & Khi thanh toán thành công, hệ thống ghi Payment SUCCESS và chuyển Order sang PAID. Tồn kho đã giữ được xác nhận là hàng đã bán, lượt mã giảm giá chuyển sang CONSUMED, điểm tích lũy được cộng đúng một lần và giỏ hàng ở client được xóa.

Khi thanh toán thất bại, hệ thống ghi Payment FAILED và chuyển Order sang PAYMENT\_FAILED. Tồn kho và lượt mã giảm giá tiếp tục được giữ, không cộng điểm tích lũy và không xóa giỏ hàng.

Customer được phép thử lại thanh toán trên cùng Order nếu đơn chưa hết hạn. Mỗi lần thử mới sử dụng Idempotency-Key mới và tạo bản ghi Payment mới, không tạo Order mới hoặc giữ thêm tài nguyên.

Việc ghi Payment, cập nhật trạng thái Order và các thay đổi dữ liệu liên quan phải bảo đảm tính nguyên tử và không được ghi nhận thanh toán thành công trùng lặp. \\
\textbf{Lý do/Nguồn} & Charter §3.1 (xử lý thành công/thất bại), §7 (PAID nếu thành công; PAYMENT\_FAILED và giữ giỏ khi lỗi); Techstack §7; Task 8, 10. \\
\textbf{Actor chính} & Backend API; Mock Payment Service; Customer \\
\end{longtable}

\textbf{Phân tích nghiệp vụ}

Kết quả thanh toán quyết định tiền, tồn kho, điểm và giỏ hàng, nên tác động đi kèm được quy định rõ:

\begin{longtable}[]{@{}
  >{\raggedright\arraybackslash}p{(\columnwidth - 10\tabcolsep) * \real{0.1167}}
  >{\raggedright\arraybackslash}p{(\columnwidth - 10\tabcolsep) * \real{0.1507}}
  >{\raggedright\arraybackslash}p{(\columnwidth - 10\tabcolsep) * \real{0.2674}}
  >{\raggedright\arraybackslash}p{(\columnwidth - 10\tabcolsep) * \real{0.2188}}
  >{\raggedright\arraybackslash}p{(\columnwidth - 10\tabcolsep) * \real{0.1183}}
  >{\raggedright\arraybackslash}p{(\columnwidth - 10\tabcolsep) * \real{0.1280}}@{}}
\toprule\noalign{}
\endhead
\bottomrule\noalign{}
\endlastfoot
\textbf{Kết quả} & \textbf{Payment} & \textbf{Order} & \textbf{Tồn kho} & \textbf{Điểm tích lũy} & \textbf{Giỏ ở client} \\
Thành công & SUCCESS & PAID & Giữ nguyên (đã trừ khi tạo đơn) & Cộng đúng một lần (FR-PRO-003) & Xóa (FR-CART-004) \\
Thất bại & FAILED & PAYMENT\_FAILED & Tiếp tục giữ tồn kho đã trừ khi tạo đơn; không hoàn (FR-INV-003). & Không cộng & Giữ \\
\end{longtable}

\textbf{Điều kiện, dữ liệu và ràng buộc liên quan}

Chuyển trạng thái Order bằng cập nhật có điều kiện theo State Machine (FR-ORD-003).

\begin{itemize}
\item
  Nếu Order đang PENDING và chưa hết hạn, hệ thống được phép xử lý lần thanh toán mới.
\item
  Nếu Order đang PAYMENT\_FAILED và chưa hết hạn, Customer được phép thử lại thanh toán trên cùng Order. Hệ thống chuyển Order về PENDING hợp lệ trước khi xử lý lần thử mới, sử dụng Idempotency-Key mới.
\item
  Nếu Order đã CANCELLED, PAID hoặc hết hạn, hệ thống từ chối lần thanh toán mới với mã lỗi 409 phù hợp và không tạo Payment mới.
\item
  Request phát lại cùng Idempotency-Key và cùng nội dung phải được kiểm tra theo FR-PAY-003 để trả kết quả đã lưu, không thực hiện lại giao dịch.
\item
  Mọi thao tác cập nhật trạng thái và dữ liệu thanh toán phải bảo đảm tính nguyên tử, chống xử lý đồng thời gây thanh toán trùng.
\end{itemize}

\begin{itemize}
\item
  Cộng điểm và ghi nhận việc dùng coupon (FR-PRO) chỉ xảy ra khi đơn chuyển PAID, trong cùng transaction.
\item
  PAYMENT\_FAILED là trạng thái cho phép giữ lại, không phải trạng thái kết thúc.
\item
  Khách hàng có thể thử lại thanh toán cho cùng một đơn hàng trong thời hạn cho phép. Mỗi lần thử lại thì sử dụng một Idempotency-Key mới và tạo một bản ghi Payment mới.
\item
  Nếu lần thanh toán tiếp theo thành công, đơn chuyển PAID, hệ thống cộng điểm đúng một lần và xóa giỏ hàng.
\item
  Nếu đơn đã hết hạn, khách không được thử lại thanh toán. Đơn được hủy và các tài nguyên đã giữ được hoàn trả đúng một lần.
\item
  Frontend xử lý theo trường trạng thái trong phản hồi: PAID thì xóa giỏ; FAILED thì giữ giỏ, hiển thị lý do chung và nút đặt lại.
\item
  Lý do thất bại trả cho client chỉ ở mức chung (ví dụ ``Thanh toán bị từ chối''), không nêu chi tiết nội bộ.
\end{itemize}

\textbf{Vấn đề có thể gặp và phân tích tác động}

\begin{longtable}[]{@{}
  >{\raggedright\arraybackslash}p{(\columnwidth - 2\tabcolsep) * \real{0.2500}}
  >{\raggedright\arraybackslash}p{(\columnwidth - 2\tabcolsep) * \real{0.7500}}@{}}
\toprule\noalign{}
\begin{minipage}[b]{\linewidth}\raggedright
\textbf{Vấn đề}
\end{minipage} & \begin{minipage}[b]{\linewidth}\raggedright
\textbf{Tác động/Nguyên nhân}
\end{minipage} \\
\midrule\noalign{}
\endhead
\bottomrule\noalign{}
\endlastfoot
Ghi Payment thành công nhưng đổi trạng thái đơn thất bại & Tiền đã ``trả'' nhưng đơn vẫn PENDING. \\
Hoàn tồn kho hoặc cộng điểm hai lần khi retry & Tồn kho và điểm sai (FR-INV-003, FR-PRO-003). \\
Hủy hoặc hết hạn và thanh toán xảy ra đồng thời & Trạng thái đơn không nhất quán với tồn kho và điểm. \\
Mất kết nối sau khi PAID, trước khi FE nhận phản hồi & Giỏ chưa xóa (hạn chế đã nêu ở FR-CART-004). \\
Cộng điểm cho đơn thất bại & Điểm sai. \\
\end{longtable}

\textbf{Hướng xử lý/giải pháp phân tích đề xuất}

\begin{itemize}
\item
  Việc tạo Payment, cập nhật trạng thái Order và các tác động nghiệp vụ liên quan phải được xử lý nguyên tử trong transaction.
\item
  Khi Payment SUCCESS, hệ thống chuyển Order sang PAID, xác nhận lượt mã giảm giá đã sử dụng và cộng điểm tích lũy đúng một lần.
\item
  Khi Payment FAILED, hệ thống chuyển Order sang PAYMENT\_FAILED, tiếp tục giữ tồn kho và lượt mã giảm giá, không cộng điểm.
\item
  Nếu xảy ra lỗi kỹ thuật, toàn bộ transaction phải rollback. Order trở về trạng thái trước lần thử thanh toán; không tạo Payment mới và không phát sinh tác động nghiệp vụ một phần.
\item
  Hoàn tồn và cộng điểm chỉ chạy khi cập nhật trạng thái có điều kiện thành công, nên mỗi tác động xảy ra đúng một lần.
\item
  Khi không chắc kết quả (mất mạng), UI hướng khách xem lịch sử đơn; gửi lại cùng Idempotency-Key là an toàn (FR-PAY-003).
\end{itemize}

\textbf{Chuyển giao sang thiết kế}

\begin{itemize}
\item
  State Diagram Order và Payment; Sequence thanh toán (nhánh PAID, nhánh FAILED); UI trạng thái thanh toán; test từng nhánh và concurrency test.
\end{itemize}

\textbf{Tiêu chí kiểm chứng ở mức yêu cầu}

\begin{itemize}
\item
  Thành công: Payment SUCCESS, đơn PAID, tồn kho bằng giá trị sau khi tạo đơn, điểm cộng đúng một lần, giỏ ở FE trống.
\item
  Thanh toán thất bại: Payment FAILED, Order chuyển PAYMENT\_FAILED, tồn kho và lượt mã giảm giá tiếp tục được giữ, không cộng điểm tích lũy, giỏ hàng ở client được giữ nguyên.
\item
  Sau PAYMENT\_FAILED, Customer có thể thử lại thanh toán trên cùng Order nếu chưa hết hạn. Lần thử mới sử dụng Idempotency-Key mới và tạo Payment mới, không trừ thêm tồn kho, không giữ thêm lượt mã giảm giá và không tạo Order mới.
\item
  Thanh toán và hủy (hoặc hết hạn) đồng thời trên một đơn: chỉ một thao tác thành công; trạng thái cuối nhất quán với tồn kho và điểm.
\item
  Gây lỗi giữa chừng trong transaction (môi trường test) không để lại trạng thái nửa vời.
\end{itemize}

\textbf{FR-PAY-003 - Idempotency-Key cho thanh toán}

\begin{longtable}[]{@{}
  >{\raggedright\arraybackslash}p{(\columnwidth - 2\tabcolsep) * \real{0.2500}}
  >{\raggedright\arraybackslash}p{(\columnwidth - 2\tabcolsep) * \real{0.7500}}@{}}
\toprule\noalign{}
\begin{minipage}[b]{\linewidth}\raggedright
\textbf{Thuộc tính}
\end{minipage} & \begin{minipage}[b]{\linewidth}\raggedright
\textbf{Nội dung}
\end{minipage} \\
\midrule\noalign{}
\endhead
\bottomrule\noalign{}
\endlastfoot
\textbf{Mức ưu tiên} & MUST \\
\textbf{Trạng thái} & CONFIRMED \\
\textbf{Yêu cầu} & POST /payments bắt buộc có Idempotency-Key. Cùng khóa và cùng nội dung thì trả lại đúng kết quả đã lưu, không tạo giao dịch mới và không đổi tồn kho hay điểm lần nữa. Cùng khóa nhưng khác nội dung thì bị từ chối an toàn (409 IDEMPOTENCY\_KEY\_CONFLICT) và không thay đổi dữ liệu. Bảo đảm cả khi các request đến đồng thời. \\
\textbf{Lý do/Nguồn} & Charter §6 (R4), §7 (lặp khóa không tạo giao dịch kép); Techstack §7 (Idempotency); Task 10. \\
\textbf{Actor chính} & Customer; Backend API \\
\end{longtable}

\textbf{Phân tích nghiệp vụ}

Khách có thể bấm đúp, mạng chậm hoặc client tự gửi lại, khiến cùng một ý định thanh toán đến nhiều lần. Kết quả nghiệp vụ phải như thể chỉ gửi một lần (định nghĩa ở mục 5).

\textbf{Điều kiện, dữ liệu và ràng buộc liên quan}

\begin{itemize}
\item
  Khóa do client sinh (khuyến nghị UUID v4), dài 8 đến 64 ký tự.
\item
  Khóa có phạm vi theo người dùng: duy nhất theo cặp (người dùng, khóa) ở tầng dữ liệu, nên khóa của người này không ảnh hưởng người khác.
\item
  ``Cùng nội dung'' nghĩa là cùng mã đơn, cùng phương thức và cùng dữ liệu test; hệ thống lưu giá trị băm của nội dung để so sánh, không lưu dữ liệu thẻ đầy đủ.
\item
  Phát lại (replay): trả đúng nội dung phản hồi gốc (kể cả kết quả FAILED), kèm dấu hiệu phát lại (ví dụ header Idempotent-Replayed: true). Phát lại được xử lý \textbf{trước} bước kiểm tra trạng thái đơn, vì đơn lúc đó có thể đã PAID hoặc PAYMENT\_FAILED.
\item
  Chỉ kết quả nghiệp vụ (SUCCESS, FAILED) được lưu và phát lại; lỗi kỹ thuật không lưu, khóa dùng lại được.
\item
  Các request đồng thời cùng khóa: chỉ một request được xử lý; request còn lại nhận lại kết quả đã lưu, hoặc 409 PAYMENT\_IN\_PROGRESS nếu kết quả chưa sẵn sàng để client gửi lại cùng khóa.
\item
  Cùng Idempotency-Key và cùng nội dung yêu cầu: trả lại kết quả đã lưu, không tạo thêm Payment.
\item
  Cùng Idempotency-Key nhưng khác nội dung: từ chối với 409 IDEMPOTENCY\_KEY\_CONFLICT.
\item
  Idempotency-Key mới trên đơn PAYMENT\_FAILED chưa hết hạn: cho phép tạo lần thử thanh toán mới.
\item
  Idempotency-Key mới trên đơn đã PAID, CANCELLED hoặc COMPLETED: từ chối, không tạo Payment mới.
\item
  Một đơn có thể có nhiều Payment FAILED nhưng tối đa một Payment SUCCESS.
\item
  Mỗi Payment phải gắn với đúng một lần thử thanh toán và không thay đổi trạng thái sau khi ghi nhận.
\item
  Khóa được lưu cùng bản ghi Payment và không hết hạn trong MVP.
\item
  Frontend tạo một Idempotency-Key riêng cho mỗi lần thử thanh toán mới.
\item
  Khi gửi lại cùng một lần thử do bấm đúp, lỗi mạng hoặc chưa nhận được phản hồi, frontend phải sử dụng lại Idempotency-Key cũ cùng nội dung request để nhận kết quả đã lưu.
\item
  Khi Payment FAILED và Customer chủ động thử lại thanh toán trên cùng Order, frontend phải tạo Idempotency-Key mới cho lần thử mới.
\item
  Không tạo Idempotency-Key mới chỉ vì gửi lại request của cùng một lần thử.
\end{itemize}

\textbf{Vấn đề có thể gặp và phân tích tác động}

\begin{longtable}[]{@{}
  >{\raggedright\arraybackslash}p{(\columnwidth - 2\tabcolsep) * \real{0.2500}}
  >{\raggedright\arraybackslash}p{(\columnwidth - 2\tabcolsep) * \real{0.7500}}@{}}
\toprule\noalign{}
\begin{minipage}[b]{\linewidth}\raggedright
\textbf{Vấn đề}
\end{minipage} & \begin{minipage}[b]{\linewidth}\raggedright
\textbf{Tác động/Nguyên nhân}
\end{minipage} \\
\midrule\noalign{}
\endhead
\bottomrule\noalign{}
\endlastfoot
Bấm đúp hoặc gửi lại khi mạng chậm & Hai giao dịch thanh toán cho một đơn. \\
Hai request song song cùng khóa (kiểm tra rồi mới ghi) & Cả hai cùng qua, tạo hai Payment. \\
Cùng khóa nhưng khác nội dung & Trả kết quả của giao dịch khác hoặc xử lý nhầm. \\
Khóa mới cho đơn đã thanh toán & Giao dịch thứ hai cho cùng một đơn. \\
Hai người dùng trùng khóa & Lộ hoặc ghi đè kết quả của người khác. \\
Frontend sinh khóa mới ở mỗi lần gửi & Cơ chế idempotency vô tác dụng. \\
\end{longtable}

\textbf{Hướng xử lý/giải pháp phân tích đề xuất}

\begin{itemize}
\item
  Ràng buộc duy nhất (người dùng, khóa) ở cơ sở dữ liệu, xử lý trong transaction để race được giải quyết ở tầng dữ liệu.
\item
  Lưu giá trị băm nội dung và đủ dữ liệu để dựng lại phản hồi gốc.
\item
  Frontend giữ khóa trong bộ nhớ hoặc sessionStorage trong suốt một lần thanh toán.
\end{itemize}

\textbf{Chuyển giao sang thiết kế}

\begin{itemize}
\item
  ERD Payment (khóa, người dùng, giá trị băm nội dung, dữ liệu phản hồi); Sequence thanh toán (nhánh phát lại); API (header, mã lỗi 400 và 409); quản lý khóa ở frontend; concurrency test.
\end{itemize}

\textbf{Tiêu chí kiểm chứng ở mức yêu cầu}

\begin{itemize}
\item
  Gửi hai lần cùng khóa và cùng nội dung: chỉ một Payment; kết quả giống hệt; tồn kho và điểm thay đổi đúng một lần; lần hai có dấu hiệu phát lại.
\item
  Gửi 10 request song song cùng khóa: đúng một Payment, điểm cộng đúng một lần; mọi request nhận cùng kết quả, hoặc nhận PAYMENT\_IN\_PROGRESS rồi lấy lại đúng kết quả khi gửi lại.
\item
  Cùng khóa nhưng khác phương thức hoặc khác mã đơn: 409 IDEMPOTENCY\_KEY\_CONFLICT, dữ liệu không đổi.
\item
  Khóa mới cho đơn đã PAID: 409 INVALID\_TRANSITION, không có Payment thứ hai.
\item
  Phát lại một kết quả FAILED trả FAILED và không hoàn tồn kho lần hai.
\item
  Người dùng B dùng đúng khóa của người dùng A không nhận được kết quả của A.
\item
  Thiếu khóa trả 400.
\end{itemize}

\textbf{FR-PAY-004 - Bản ghi thanh toán và trạng thái thanh toán tách biệt}

\begin{longtable}[]{@{}
  >{\raggedright\arraybackslash}p{(\columnwidth - 2\tabcolsep) * \real{0.2500}}
  >{\raggedright\arraybackslash}p{(\columnwidth - 2\tabcolsep) * \real{0.7500}}@{}}
\toprule\noalign{}
\begin{minipage}[b]{\linewidth}\raggedright
\textbf{Thuộc tính}
\end{minipage} & \begin{minipage}[b]{\linewidth}\raggedright
\textbf{Nội dung}
\end{minipage} \\
\midrule\noalign{}
\endhead
\bottomrule\noalign{}
\endlastfoot
\textbf{Mức ưu tiên} & MUST \\
\textbf{Trạng thái} & DERIVED \\
\textbf{Yêu cầu} & Mỗi lần thử thanh toán có kết quả từ Mock Payment Service được lưu thành một bản ghi Payment riêng, gắn với Order tương ứng.

Một Order có thể có nhiều Payment FAILED nhưng tối đa một Payment SUCCESS. Các bản ghi Payment đã ghi nhận không được sửa hoặc xóa thông qua API và không được chứa thông tin thẻ nhạy cảm. \\
\textbf{Lý do/Nguồn} & Techstack §6 (Payment: mã đơn, khóa idempotency, số tiền, phương thức); Charter §6 (R4).

Việc tách trạng thái Payment khỏi trạng thái Order và cho phép nhiều lần thử thanh toán trên cùng Order là quyết định thiết kế của nhóm nhằm hỗ trợ retry, lưu lịch sử giao dịch và phục vụ đối soát.

Mỗi Order chỉ được có tối đa một Payment SUCCESS để ngăn ghi nhận thanh toán thành công trùng lặp. \\
\textbf{Actor chính} & Backend API \\
\end{longtable}

\textbf{Phân tích nghiệp vụ}

Trạng thái đơn mô tả vòng đời của đơn; trạng thái thanh toán mô tả kết quả của giao dịch. Gộp chung làm khó đối soát và dễ sai khi đơn đã sang PREPARING hoặc COMPLETED. Một Order có thể có nhiều lần thử thanh toán. Mỗi lần thử được lưu thành một Payment riêng để bảo toàn lịch sử giao dịch và phục vụ đối soát.

PAYMENT\_FAILED không phải trạng thái kết thúc của Order. Customer được phép thử lại thanh toán trên cùng Order khi đơn chưa hết hạn.

Một Order có thể có nhiều Payment FAILED nhưng tối đa một Payment SUCCESS. Backend phải ngăn thanh toán thành công trùng lặp, kể cả khi các request sử dụng Idempotency-Key khác nhau và được gửi đồng thời.

\textbf{Điều kiện, dữ liệu và ràng buộc liên quan}

\begin{itemize}
\item
  Các trường tối thiểu: mã giao dịch, mã đơn, người dùng, khóa idempotency, giá trị băm nội dung, số tiền, phương thức, trạng thái thanh toán, mã tham chiếu của mock, thời điểm.
\item
  Số tiền bằng tổng đơn tại thời điểm thanh toán và bất biến sau khi ghi.
\item
  Không lưu số thẻ đầy đủ hay mã bảo mật; nếu cần hiển thị chỉ lưu 4 số cuối; không ghi log các dữ liệu này.
\item
  Order PENDING mới tạo chưa có Payment; Order PENDING sau khi thử lại có thể đã có các Payment FAILED trước đó.
\item
  Order CANCELLED có thể chưa có Payment hoặc có các Payment FAILED từ những lần thử trước khi hủy.
\item
  Payment SUCCESS tương ứng với Order từ trạng thái PAID trở đi. Payment FAILED có thể tồn tại khi Order ở trạng thái PAYMENT\_FAILED, PENDING (sau khi thử lại), CANCELLED hoặc đã PAID sau một lần thử thanh toán thành công tiếp theo.
\item
  Trạng thái Payment phản ánh kết quả của từng lần thử thanh toán, không trực tiếp quyết định trạng thái hiện tại của Order nếu Order đã có các lần thử tiếp theo.
\end{itemize}

\textbf{Vấn đề có thể gặp và phân tích tác động}

\begin{longtable}[]{@{}
  >{\raggedright\arraybackslash}p{(\columnwidth - 2\tabcolsep) * \real{0.2500}}
  >{\raggedright\arraybackslash}p{(\columnwidth - 2\tabcolsep) * \real{0.7500}}@{}}
\toprule\noalign{}
\begin{minipage}[b]{\linewidth}\raggedright
\textbf{Vấn đề}
\end{minipage} & \begin{minipage}[b]{\linewidth}\raggedright
\textbf{Tác động/Nguyên nhân}
\end{minipage} \\
\midrule\noalign{}
\endhead
\bottomrule\noalign{}
\endlastfoot
Dùng chung một cột trạng thái cho đơn và thanh toán & Mất thông tin kết quả giao dịch khi đơn sang trạng thái tiếp theo. \\
Hai Payment SUCCESS trên cùng một Order & Thanh toán thành công trùng lặp, sai dữ liệu đối soát nếu backend không đồng bộ các request thanh toán song song. \\
Lưu dữ liệu thẻ đầy đủ & Rò rỉ dữ liệu nhạy cảm, trái nguyên tắc bảo mật. \\
Sửa hoặc xóa Payment & Mất dấu vết đối soát. \\
\end{longtable}

\textbf{Hướng xử lý/giải pháp phân tích đề xuất}

\begin{itemize}
\item
  Hai enum riêng: OrderStatus do FR-ORD-003 quản lý, PaymentStatus chỉ phản ánh giao dịch.
\item
  Không đặt ràng buộc UNIQUE trên orderId của bảng Payment vì một Order có thể có nhiều lần thử thanh toán.
\item
  Đặt ràng buộc UNIQUE trên (userId, idempotencyKey).
\item
  Bảo đảm tối đa một Payment SUCCESS trên mỗi Order bằng ràng buộc dữ liệu phù hợp và cơ chế xử lý transaction an toàn.
\item
  Không cung cấp API sửa hoặc xóa bản ghi Payment đã hoàn tất.
\item
  Không cung cấp API sửa hoặc xóa Payment.
\end{itemize}

\textbf{Chuyển giao sang thiết kế}

\begin{itemize}
\item
  ERD Payment (các ràng buộc duy nhất); enum trạng thái; test ràng buộc dữ liệu.
\end{itemize}

\textbf{Tiêu chí kiểm chứng ở mức yêu cầu}

\begin{itemize}
\item
  Một Order có thể lưu nhiều Payment FAILED từ những lần thử thanh toán khác nhau.
\item
  Mỗi Order chỉ được có tối đa một Payment SUCCESS, kể cả khi nhiều request thanh toán được gửi đồng thời.
\item
  Cùng Idempotency-Key và cùng nội dung request phải trả kết quả đã lưu, không tạo Payment mới.
\item
  Hai request thanh toán đồng thời sử dụng hai Idempotency-Key khác nhau không thể tạo hai Payment SUCCESS trên cùng Order.
\item
  Payment đã hoàn tất không được sửa hoặc xóa thông qua API.
\item
  Cơ sở dữ liệu không được chứa số thẻ đầy đủ hoặc thông tin thanh toán nhạy cảm.
\item
  Số tiền của mỗi Payment phải bằng tổng tiền đã chốt của Order tương ứng.
\item
  Order CANCELLED có thể không có Payment hoặc có một hay nhiều Payment FAILED trước đó, nhưng không được có Payment SUCCESS.
\end{itemize}

\textbf{FR-PAY-005 - Xác nhận đơn sau thanh toán}

\begin{longtable}[]{@{}
  >{\raggedright\arraybackslash}p{(\columnwidth - 2\tabcolsep) * \real{0.2500}}
  >{\raggedright\arraybackslash}p{(\columnwidth - 2\tabcolsep) * \real{0.7500}}@{}}
\toprule\noalign{}
\begin{minipage}[b]{\linewidth}\raggedright
\textbf{Thuộc tính}
\end{minipage} & \begin{minipage}[b]{\linewidth}\raggedright
\textbf{Nội dung}
\end{minipage} \\
\midrule\noalign{}
\endhead
\bottomrule\noalign{}
\endlastfoot
\textbf{Mức ưu tiên} & MUST \\
\textbf{Trạng thái} & CONFIRMED (xác nhận và mã đơn); DERIVED (cách dựng lại trang) \\
\textbf{Yêu cầu} & Sau khi thanh toán, giao diện hiển thị màn hình xác nhận. Nếu thành công: mã đơn, trạng thái, các món (bản chụp), tổng tiền, phương thức (che thông tin) và thời điểm. Nếu thất bại: thông báo bị từ chối và hướng dẫn đặt lại. Trang xác nhận mở lại được từ lịch sử đơn. \\
\textbf{Lý do/Nguồn} & Charter §2 (nhận mã đơn), §3.1 (xác nhận và lịch sử đơn); Task 9. \\
\textbf{Actor chính} & Customer; Web Client \\
\end{longtable}

\textbf{Phân tích nghiệp vụ}

Khách cần bằng chứng rõ ràng rằng đơn đã được ghi nhận, cùng mã đơn để theo dõi. Nếu mất kết nối hoặc tải lại trang, khách phải xem lại được kết quả mà không vô tình thanh toán lần nữa.

\textbf{Điều kiện, dữ liệu và ràng buộc liên quan}

\begin{itemize}
\item
  Dữ liệu xác nhận lấy từ phản hồi của POST /payments và có thể dựng lại từ GET /orders/me (hệ thống không có GET /orders/:id).
\item
  Tổng tiền hiển thị lấy từ server, không dùng tổng tạm tính của giỏ.
\item
  Nội dung các món lấy từ bản chụp của đơn (FR-ORD-002); phương thức hiển thị dạng che (ví dụ ``Thẻ kết thúc ****0002'').
\item
  Trang xác nhận không thực hiện hành động ghi nào; tải lại không gọi lại thanh toán.
\item
  Thành công: giỏ trống (FR-CART-004). Thất bại: giỏ còn, có nút quay lại giỏ và lý do thất bại ở mức chung.
\item
  Chỉ chủ đơn xem được (FR-AUTHZ-002).
\end{itemize}

\textbf{Vấn đề có thể gặp và phân tích tác động}

\begin{longtable}[]{@{}
  >{\raggedright\arraybackslash}p{(\columnwidth - 2\tabcolsep) * \real{0.2500}}
  >{\raggedright\arraybackslash}p{(\columnwidth - 2\tabcolsep) * \real{0.7500}}@{}}
\toprule\noalign{}
\begin{minipage}[b]{\linewidth}\raggedright
\textbf{Vấn đề}
\end{minipage} & \begin{minipage}[b]{\linewidth}\raggedright
\textbf{Tác động/Nguyên nhân}
\end{minipage} \\
\midrule\noalign{}
\endhead
\bottomrule\noalign{}
\endlastfoot
Tải lại trang xác nhận làm mất dữ liệu phản hồi & Màn hình trống, khách không biết đơn đã thành công chưa. \\
Mất kết nối ngay sau khi gửi thanh toán & Khách không biết đã trả hay chưa và có thể trả lại. \\
Hiển thị tổng tạm tính thay vì tổng của đơn & Số tiền hiển thị khác số tiền đã thanh toán. \\
Hiển thị lỗi chi tiết của nhà cung cấp & Lộ thông tin nội bộ. \\
\end{longtable}

\textbf{Hướng xử lý/giải pháp phân tích đề xuất}

\begin{itemize}
\item
  Trang xác nhận có địa chỉ riêng chứa mã đơn; khi thiếu dữ liệu trong bộ nhớ thì tải từ lịch sử đơn.
\item
  Khi không chắc kết quả, UI hướng khách vào lịch sử đơn; nếu đơn còn PENDING thì thanh toán lại với cùng khóa là an toàn (FR-PAY-003).
\item
  Nhãn trạng thái dễ hiểu cho khách (ví dụ ``Đã thanh toán'', ``Thanh toán thất bại'').
\end{itemize}

\textbf{Chuyển giao sang thiết kế}

\begin{itemize}
\item
  Use Case Xem xác nhận đơn; UI trang xác nhận (thành công, thất bại, đang tải, lỗi); test E2E.
\end{itemize}

\textbf{Tiêu chí kiểm chứng ở mức yêu cầu}

\begin{itemize}
\item
  Sau thanh toán thành công, trang hiển thị đúng mã đơn, tổng tiền và các món khớp với đơn đã lưu.
\item
  Tải lại trang xác nhận vẫn hiển thị đúng thông tin và không tạo giao dịch mới.
\item
  Sau thanh toán thất bại, trang hiển thị thông báo chung và giỏ hàng vẫn còn.
\item
  Đơn PAYMENT\_FAILED hiển thị đúng nhãn trong lịch sử đơn.
\item
  Không mở được trang xác nhận của đơn thuộc người khác.
\end{itemize}

\textbf{FR-PAY-006 - Kịch bản giả lập thanh toán thành công hoặc thất bại}

\begin{longtable}[]{@{}
  >{\raggedright\arraybackslash}p{(\columnwidth - 2\tabcolsep) * \real{0.2500}}
  >{\raggedright\arraybackslash}p{(\columnwidth - 2\tabcolsep) * \real{0.7500}}@{}}
\toprule\noalign{}
\begin{minipage}[b]{\linewidth}\raggedright
\textbf{Thuộc tính}
\end{minipage} & \begin{minipage}[b]{\linewidth}\raggedright
\textbf{Nội dung}
\end{minipage} \\
\midrule\noalign{}
\endhead
\bottomrule\noalign{}
\endlastfoot
\textbf{Mức ưu tiên} & MUST (MOCK) \\
\textbf{Trạng thái} & DERIVED \\
\textbf{Yêu cầu} & Mock Payment Service quyết định kết quả theo quy tắc xác định dựa trên dữ liệu test của phương thức, để cùng đầu vào luôn cho cùng kết quả và nhóm demo được cả hai nhánh mà không phải can thiệp vào server. Client không bao giờ gửi kết quả. \\
\textbf{Lý do/Nguồn} & Charter §3.1 (mock gateway), §7; Techstack §3 (MOCK); mục 3 (không tin kết quả do client gửi). Quy tắc cụ thể là quyết định của nhóm. \\
\textbf{Actor chính} & Mock Payment Service; Nhóm phát triển; Giảng viên (khi nghiệm thu) \\
\end{longtable}

\textbf{Phân tích nghiệp vụ}

Muốn nghiệm thu nhánh PAYMENT\_FAILED thì phải tạo được thất bại theo ý muốn. Nếu mock luôn thành công thì không chứng minh được nhánh lỗi; nếu client gửi cờ ``thất bại'' thì mâu thuẫn với nguyên tắc ở mục 3. Giải pháp là để mock suy ra kết quả từ dữ liệu test, giống cách cổng thanh toán thật có thẻ test riêng.

\textbf{Điều kiện, dữ liệu và ràng buộc liên quan}

\begin{itemize}
\item
  Quy tắc đề xuất (nhóm chốt giá trị cuối cùng):
\end{itemize}

\begin{longtable}[]{@{}
  >{\raggedright\arraybackslash}p{(\columnwidth - 4\tabcolsep) * \real{0.3333}}
  >{\raggedright\arraybackslash}p{(\columnwidth - 4\tabcolsep) * \real{0.3333}}
  >{\raggedright\arraybackslash}p{(\columnwidth - 4\tabcolsep) * \real{0.3334}}@{}}
\toprule\noalign{}
\endhead
\bottomrule\noalign{}
\endlastfoot
\textbf{Phương thức} & \textbf{Dữ liệu test} & \textbf{Kết quả} \\
CARD & Số thẻ 16 chữ số kết thúc bằng 0002 & FAILED (thẻ bị từ chối) \\
CARD & Số thẻ 16 chữ số khác & SUCCESS \\
WALLET & Số điện thoại ví 10 chữ số kết thúc bằng 0000 & FAILED (không đủ số dư) \\
WALLET & Số điện thoại ví 10 chữ số khác & SUCCESS \\
\end{longtable}

\begin{itemize}
\item
  Dữ liệu test sai định dạng (không đủ chữ số hoặc có ký tự khác) trả 400, không phải FAILED.
\item
  Quy tắc xác định, không dùng yếu tố ngẫu nhiên.
\item
  Mock không có số dư hay trạng thái thật; chỉ áp dụng quy tắc trên.
\item
  Danh sách dữ liệu test được ghi vào README và hiển thị gợi ý trên màn hình thanh toán.
\item
  Mock nhận số tiền từ server, không từ client.
\item
  Việc gây lỗi kỹ thuật (hết thời gian chờ, lỗi nội bộ) để kiểm thử chỉ bật được bằng cấu hình môi trường test, không qua API công khai.
\end{itemize}

\textbf{Vấn đề có thể gặp và phân tích tác động}

\begin{longtable}[]{@{}
  >{\raggedright\arraybackslash}p{(\columnwidth - 2\tabcolsep) * \real{0.2500}}
  >{\raggedright\arraybackslash}p{(\columnwidth - 2\tabcolsep) * \real{0.7500}}@{}}
\toprule\noalign{}
\begin{minipage}[b]{\linewidth}\raggedright
\textbf{Vấn đề}
\end{minipage} & \begin{minipage}[b]{\linewidth}\raggedright
\textbf{Tác động/Nguyên nhân}
\end{minipage} \\
\midrule\noalign{}
\endhead
\bottomrule\noalign{}
\endlastfoot
Mock luôn thành công & Không kiểm thử và demo được nhánh PAYMENT\_FAILED. \\
Kết quả ngẫu nhiên & Kiểm thử không lặp lại được. \\
Client điều khiển kết quả & Người dùng tự báo thành công, trái mục 3. \\
Lưu hoặc ghi log số thẻ test đầy đủ & Thói quen xấu và rủi ro rò rỉ dữ liệu. \\
\end{longtable}

\textbf{Hướng xử lý/giải pháp phân tích đề xuất}

\begin{itemize}
\item
  Một module Mock Payment riêng với quy tắc nêu trên; bộ unit test cho từng dòng của bảng.
\item
  Trường lạ do client gửi (ví dụ kết quả mong muốn) bị bỏ qua.
\item
  Ghi quy tắc và dữ liệu test vào README để giảng viên tự thử.
\end{itemize}

\textbf{Chuyển giao sang thiết kế}

\begin{itemize}
\item
  Module Mock Payment; README (bảng dữ liệu test); UI gợi ý dữ liệu test; unit test.
\end{itemize}

\textbf{Tiêu chí kiểm chứng ở mức yêu cầu}

\begin{itemize}
\item
  Thẻ kết thúc 0002 luôn cho FAILED, gửi lặp 10 lần vẫn giống nhau; thẻ khác luôn cho SUCCESS.
\item
  Ví kết thúc 0000 luôn cho FAILED; ví khác luôn cho SUCCESS.
\item
  Dữ liệu test sai định dạng trả 400.
\item
  Request có thêm trường kết quả không làm thay đổi kết quả.
\item
  README liệt kê đủ dữ liệu test cho cả hai nhánh.
\end{itemize}

\textbf{7.8. Loyalty \& Promo}

\textbf{Actor/phạm vi chính:} Customer; Backend API (module Promo/Loyalty, Order, Payment); Nhóm phát triển (nạp dữ liệu mẫu)

\textbf{Artifact thiết kế chịu ảnh hưởng:} Use Case áp mã giảm giá và xem điểm, ERD (Coupon, CouponRedemption, User.loyaltyPoints, Order), Sequence tạo đơn và thanh toán, API Order/Payment, UI giỏ hàng và thanh toán, concurrency test, Test Strategy

\textbf{FR-PRO-001 - Kiểm tra mã giảm giá hợp lệ}

\begin{longtable}[]{@{}
  >{\raggedright\arraybackslash}p{(\columnwidth - 2\tabcolsep) * \real{0.2500}}
  >{\raggedright\arraybackslash}p{(\columnwidth - 2\tabcolsep) * \real{0.7500}}@{}}
\toprule\noalign{}
\begin{minipage}[b]{\linewidth}\raggedright
\textbf{Thuộc tính}
\end{minipage} & \begin{minipage}[b]{\linewidth}\raggedright
\textbf{Nội dung}
\end{minipage} \\
\midrule\noalign{}
\endhead
\bottomrule\noalign{}
\endlastfoot
\textbf{Mức ưu tiên} & MUST \\
\textbf{Trạng thái} & CONFIRMED (kiểm tra mã hợp lệ); DERIVED (các điều kiện cụ thể) \\
\textbf{Yêu cầu} & Customer có thể nhập một mã giảm giá khi tạo đơn. Backend kiểm tra mã hợp lệ (tồn tại, đã đến hạn, chưa hết hạn, đơn đạt giá trị tối thiểu); mã không hợp lệ thì từ chối toàn bộ đơn kèm lý do, không tự bỏ mã để tiếp tục. Mỗi đơn dùng tối đa một mã. \\
\textbf{Lý do/Nguồn} & Charter §3.1 (khuyến mãi theo Task 10); Techstack §7 (kiểm tra mã giảm giá hợp lệ); Task 10. Các điều kiện cụ thể là đề xuất của nhóm. \\
\textbf{Actor chính} & Customer; Backend API \\
\end{longtable}

\textbf{Phân tích nghiệp vụ}

Mã giảm giá thay đổi số tiền khách phải trả nên phải được server xác thực, không tin client (mục 3). Nếu mã sai mà hệ thống lặng lẽ bỏ qua, khách sẽ trả giá khác với kỳ vọng. Vì vậy mã không hợp lệ làm đơn bị từ chối và khách phải sửa hoặc bỏ mã.

\textbf{Điều kiện, dữ liệu và ràng buộc liên quan}

\begin{itemize}
\item
  Mã nằm trong payload của POST /orders (một trường tùy chọn); Guest chưa đăng nhập không dùng được mã vì chưa tạo được đơn.
\item
  Chuẩn hóa mã: bỏ khoảng trắng đầu cuối, không phân biệt hoa thường, tối đa 30 ký tự; mã duy nhất trong hệ thống.
\item
  Mỗi mã gồm: loại (PERCENT hoặc FIXED), giá trị, giá trị đơn tối thiểu, mức giảm tối đa (chỉ với PERCENT), thời điểm bắt đầu và hết hạn, giới hạn tổng số lượt dùng, giới hạn mỗi người (mặc định 1 lượt).
\item
  Đơn tối thiểu so với \textbf{tạm tính} (tổng các dòng, trước giảm giá).
\item
  So sánh thời hạn theo giờ máy chủ (UTC); giao diện hiển thị theo giờ Việt Nam.
\item
  Lỗi trả mã máy đọc được:
\end{itemize}

\begin{longtable}[]{@{}
  >{\raggedright\arraybackslash}p{(\columnwidth - 4\tabcolsep) * \real{0.3204}}
  >{\raggedright\arraybackslash}p{(\columnwidth - 4\tabcolsep) * \real{0.3190}}
  >{\raggedright\arraybackslash}p{(\columnwidth - 4\tabcolsep) * \real{0.3605}}@{}}
\toprule\noalign{}
\endhead
\bottomrule\noalign{}
\endlastfoot
\textbf{Tình huống} & \textbf{Mã lỗi} & \textbf{Mã máy đọc / lý do} \\
Mã không tồn tại, chưa đến hạn, đã hết hạn, hoặc đơn chưa đạt giá trị tối thiểu & 400 & INVALID\_COUPON kèm reason (xem mục Hướng xử lý) \\
Mã hết lượt dùng hoặc người này đã dùng đủ lượt & 409 & COUPON\_UNAVAILABLE (FR-PRO-004) \\
\end{longtable}

\begin{itemize}
\item
  (SHOULD) Có endpoint xem trước, ví dụ POST /coupons/preview (yêu cầu JWT), nhận các dòng giỏ và mã, trả tạm tính, số tiền giảm, tổng dự kiến hoặc lý do không hợp lệ. Endpoint này \textbf{không giữ lượt dùng}.
\end{itemize}

\textbf{Vấn đề có thể gặp và phân tích tác động}

\begin{longtable}[]{@{}
  >{\raggedright\arraybackslash}p{(\columnwidth - 2\tabcolsep) * \real{0.2500}}
  >{\raggedright\arraybackslash}p{(\columnwidth - 2\tabcolsep) * \real{0.7500}}@{}}
\toprule\noalign{}
\begin{minipage}[b]{\linewidth}\raggedright
\textbf{Vấn đề}
\end{minipage} & \begin{minipage}[b]{\linewidth}\raggedright
\textbf{Tác động/Nguyên nhân}
\end{minipage} \\
\midrule\noalign{}
\endhead
\bottomrule\noalign{}
\endlastfoot
Client gửi số tiền giảm hoặc cờ ``mã hợp lệ'' & Khách tự giảm giá. \\
Bỏ qua mã sai mà vẫn tạo đơn & Khách trả khác kỳ vọng, khó giải thích. \\
Dò mã bằng cách thử liên tục & Lộ mã còn hiệu lực. \\
So sánh thời hạn theo giờ máy khách & Mã dùng được sau hạn hoặc bị từ chối sớm. \\
Khác biệt hoa thường, khoảng trắng & Mã đúng nhưng bị từ chối. \\
Khách xem trước hợp lệ nhưng khi đặt lại hết lượt & Lệch giữa xem trước và kết quả thật. \\
\end{longtable}

\textbf{Hướng xử lý/giải pháp phân tích đề xuất}

\begin{itemize}
\item
  Backend nhận mã, chuẩn hóa, tra cứu và kiểm tra theo thứ tự: tồn tại → đến hạn → chưa hết hạn → đơn tối thiểu; bỏ qua mọi trường giảm giá do client gửi.
\item
  Thông điệp cho mã không tồn tại dùng câu chung ``Mã giảm giá không hợp lệ'', không phân biệt với mã đã tắt.
\item
  Xem trước chỉ mang tính tham khảo; kết quả cuối là kết quả khi tạo đơn.
\item
  (SHOULD) Giới hạn tần suất thử mã nếu còn thời gian (chống dò).
\end{itemize}

\textbf{Chuyển giao sang thiết kế}

\begin{itemize}
\item
  ERD Coupon; API POST /orders (trường mã, mã lỗi), POST /coupons/preview (SHOULD); Sequence tạo đơn; UI ô nhập mã và thông báo lỗi; unit test từng điều kiện.
\end{itemize}

\textbf{Tiêu chí kiểm chứng ở mức yêu cầu}

\begin{itemize}
\item
  Mã hợp lệ được chấp nhận; chữ hoa, chữ thường hoặc có khoảng trắng đầu cuối cho cùng kết quả.
\item
  Mã không tồn tại, chưa đến hạn, đã hết hạn, hoặc đơn dưới mức tối thiểu trả 400 INVALID\_COUPON; không có đơn nào được tạo.
\item
  Request có thêm trường số tiền giảm vẫn dùng giá trị server tính.
\item
  Gửi hai mã trong một đơn bị từ chối 400.
\item
  Mã sai không làm thay đổi tồn kho hay tạo đơn.
\end{itemize}

\textbf{FR-PRO-002 - Tính số tiền giảm giá}

\begin{longtable}[]{@{}
  >{\raggedright\arraybackslash}p{(\columnwidth - 2\tabcolsep) * \real{0.2500}}
  >{\raggedright\arraybackslash}p{(\columnwidth - 2\tabcolsep) * \real{0.7500}}@{}}
\toprule\noalign{}
\begin{minipage}[b]{\linewidth}\raggedright
\textbf{Thuộc tính}
\end{minipage} & \begin{minipage}[b]{\linewidth}\raggedright
\textbf{Nội dung}
\end{minipage} \\
\midrule\noalign{}
\endhead
\bottomrule\noalign{}
\endlastfoot
\textbf{Mức ưu tiên} & MUST \\
\textbf{Trạng thái} & CONFIRMED (tính giảm giá); DERIVED (công thức và giới hạn) \\
\textbf{Yêu cầu} & Backend tính số tiền giảm từ mã và tạm tính của đơn theo công thức dưới đây, đảm bảo số tiền giảm không âm, không vượt tạm tính và tổng tiền không âm. Số tiền giảm và mã đã áp dụng được lưu vào đơn và bất biến. \\
\textbf{Lý do/Nguồn} & Charter §3.1 (khuyến mãi theo Task 10); FR-CAT-003 (tổng đơn do server tính); FR-ORD-002. Công thức là đề xuất của nhóm. \\
\textbf{Actor chính} & Backend API \\
\end{longtable}

\textbf{Phân tích nghiệp vụ}

Giảm giá là một phần của giá cuối, nên phải tính ở server với quy tắc làm tròn duy nhất, và phải chốt vào đơn như snapshot để lịch sử không đổi khi mã bị sửa sau này. Công thức:

\begin{longtable}[]{@{}
  >{\raggedright\arraybackslash}p{(\columnwidth - 2\tabcolsep) * \real{0.5000}}
  >{\raggedright\arraybackslash}p{(\columnwidth - 2\tabcolsep) * \real{0.5000}}@{}}
\toprule\noalign{}
\endhead
\bottomrule\noalign{}
\endlastfoot
\textbf{Loại} & \textbf{Công thức} \\
FIXED & Giảm = min(giá trị mã, tạm tính) \\
PERCENT & Giảm = min(làm tròn(tạm tính × giá trị ÷ 100), mức giảm tối đa của mã nếu có, tạm tính) \\
\end{longtable}

Tổng tiền = tạm tính − giảm. Làm tròn half-up đến đồng, cùng quy tắc FR-CAT-003.

\textbf{Điều kiện, dữ liệu và ràng buộc liên quan}

\begin{itemize}
\item
  Giảm giá áp lên \textbf{toàn bộ tạm tính} của đơn (gồm cả topping), không áp theo từng món.
\item
  Mức giảm là số nguyên VND không âm; không có đơn miễn phí vô hạn: tổng tiền tối thiểu bằng 0.
\item
  Order có tổng tiền bằng 0 VND sau khi áp dụng mã giảm giá hợp lệ vẫn được xử lý qua Mock Payment Service theo BR-050. Khi thanh toán thành công, Order chuyển sang PAID và được cộng 0 điểm tích lũy. (nhóm vẫn chưa chốt có thanh toán 0đ không)
\item
  Đơn lưu: mã đã dùng, số tiền giảm, tạm tính, tổng tiền (FR-ORD-002); tổng các dòng trừ giảm giá bằng tổng đơn.
\item
  Không có tổng hợp nhiều mã hoặc cộng dồn với điểm.
\end{itemize}

\textbf{Vấn đề có thể gặp và phân tích tác động}

\begin{longtable}[]{@{}
  >{\raggedright\arraybackslash}p{(\columnwidth - 2\tabcolsep) * \real{0.2500}}
  >{\raggedright\arraybackslash}p{(\columnwidth - 2\tabcolsep) * \real{0.7500}}@{}}
\toprule\noalign{}
\begin{minipage}[b]{\linewidth}\raggedright
\textbf{Vấn đề}
\end{minipage} & \begin{minipage}[b]{\linewidth}\raggedright
\textbf{Tác động/Nguyên nhân}
\end{minipage} \\
\midrule\noalign{}
\endhead
\bottomrule\noalign{}
\endlastfoot
Giảm vượt tạm tính & Tổng âm, tiền thanh toán âm. \\
PERCENT không có mức giảm tối đa & Giảm quá lớn ở đơn lớn. \\
Làm tròn không nhất quán giữa FE hiển thị và BE & Tổng xem trước khác tổng thật. \\
Lưu mã mà không lưu số tiền giảm & Sửa mã sau này làm đổi đơn cũ. \\
Áp giảm giá trên từng dòng rồi cộng lại & Sai số làm tròn, khó kiểm tra. \\
\end{longtable}

\textbf{Hướng xử lý/giải pháp phân tích đề xuất}

\begin{itemize}
\item
  Một hàm tính giảm giá dùng chung ở backend, có unit test cho từng loại và biên; xem trước (SHOULD) dùng đúng hàm này.
\item
  Kẹp kết quả vào khoảng từ 0 đến tạm tính.
\item
  Lưu số tiền giảm vào đơn trong cùng transaction tạo đơn.
\end{itemize}

\textbf{Chuyển giao sang thiết kế}

\begin{itemize}
\item
  ERD Order (mã đã dùng, số tiền giảm, tạm tính, tổng); module tính giảm giá; Sequence tạo đơn; unit test; UI hiển thị giảm giá và tổng.
\end{itemize}

\textbf{Tiêu chí kiểm chứng ở mức yêu cầu}

\begin{itemize}
\item
  Tạm tính 100.000đ, mã FIXED 20.000đ → giảm 20.000đ, tổng 80.000đ.
\item
  Tạm tính 15.000đ, mã FIXED 20.000đ → giảm 15.000đ (không vượt tạm tính), tổng 0đ.
\item
  Tạm tính 100.000đ, mã PERCENT 10\% → giảm 10.000đ, tổng 90.000đ.
\item
  Tạm tính 500.000đ, mã PERCENT 10\% tối đa 30.000đ → giảm 30.000đ.
\item
  Tạm tính 12.345đ, PERCENT 15\% → giảm 1.852đ (1.851,75 làm tròn half-up), tổng 10.493đ.
\item
  Đổi cấu hình mã sau khi tạo đơn không làm đổi số tiền giảm của đơn cũ.
\item
  Với mọi đơn, tạm tính trừ giảm giá bằng tổng đơn, và tổng không âm.
\end{itemize}

\textbf{FR-PRO-003 - Cộng điểm tích lũy sau khi thanh toán thành công}

\begin{longtable}[]{@{}
  >{\raggedright\arraybackslash}p{(\columnwidth - 2\tabcolsep) * \real{0.2500}}
  >{\raggedright\arraybackslash}p{(\columnwidth - 2\tabcolsep) * \real{0.7500}}@{}}
\toprule\noalign{}
\begin{minipage}[b]{\linewidth}\raggedright
\textbf{Thuộc tính}
\end{minipage} & \begin{minipage}[b]{\linewidth}\raggedright
\textbf{Nội dung}
\end{minipage} \\
\midrule\noalign{}
\endhead
\bottomrule\noalign{}
\endlastfoot
\textbf{Mức ưu tiên} & MUST \\
\textbf{Trạng thái} & CONFIRMED (cộng điểm sau PAID, không cộng trùng); DERIVED (tỷ lệ tích điểm) \\
\textbf{Yêu cầu} & Khi đơn chuyển từ PENDING sang PAID, hệ thống cộng điểm cho chủ đơn đúng một lần, bằng phần nguyên của (tổng tiền đã thanh toán ÷ 10.000) điểm. Đơn thất bại, bị hủy hoặc quá hạn không được cộng điểm. Gửi lại thanh toán không cộng thêm. \\
\textbf{Lý do/Nguồn} & Charter §3.1 (điểm tích lũy theo Task 10); Techstack §7 (cộng điểm chỉ sau PAID, không cộng trùng khi retry); Task 10. Tỷ lệ 1 điểm trên 10.000đ là đề xuất của nhóm. \\
\textbf{Actor chính} & Backend API; Customer \\
\end{longtable}

\textbf{Phân tích nghiệp vụ}

Điểm là một giá trị gắn với tiền nên phải đúng tuyệt đối: chỉ cộng khi tiền đã ``thu'', cộng đúng một lần, và cộng đúng số tiền thực trả (sau giảm giá). Cách an toàn là gắn việc cộng điểm vào chính bước chuyển trạng thái đơn sang PAID, để mọi cơ chế chống lặp của thanh toán (FR-PAY-003) và cập nhật có điều kiện (FR-ORD-003) tự bảo vệ điểm.

\textbf{Điều kiện, dữ liệu và ràng buộc liên quan}

\begin{itemize}
\item
  Công thức: điểm = floor(tổng tiền của đơn ÷ 10.000); tổng tiền đã trừ giảm giá; đơn dưới 10.000đ cho 0 điểm.
\item
  Cộng điểm xảy ra trong \textbf{cùng transaction} với việc đổi PENDING → PAID (FR-PAY-002); lỗi ở đâu cũng rollback toàn bộ.
\item
  Cập nhật điểm bằng phép cộng nguyên tử (cộng thêm vào giá trị hiện có), không đọc rồi ghi lại.
\item
  Số điểm đã cộng được lưu vào đơn (ví dụ pointsEarned), để đối soát và làm bằng chứng không cộng trùng.
\item
  Điểm không âm, là số nguyên; chỉ backend thay đổi (FR-AUTH-001 chặn client gửi loyaltyPoints).
\item
  Không có trừ hoặc hoàn điểm trong MVP vì đơn PAID không hủy được (FR-ORD-003).
\item
  Phát lại cùng Idempotency-Key trả kết quả cũ và \textbf{không} cộng lại.
\end{itemize}

\textbf{Vấn đề có thể gặp và phân tích tác động}

\begin{longtable}[]{@{}
  >{\raggedright\arraybackslash}p{(\columnwidth - 2\tabcolsep) * \real{0.2500}}
  >{\raggedright\arraybackslash}p{(\columnwidth - 2\tabcolsep) * \real{0.7500}}@{}}
\toprule\noalign{}
\begin{minipage}[b]{\linewidth}\raggedright
\textbf{Vấn đề}
\end{minipage} & \begin{minipage}[b]{\linewidth}\raggedright
\textbf{Tác động/Nguyên nhân}
\end{minipage} \\
\midrule\noalign{}
\endhead
\bottomrule\noalign{}
\endlastfoot
Cộng điểm khi tạo đơn hoặc khi đơn chưa PAID & Khách nhận điểm cho đơn không thanh toán. \\
Retry hoặc bấm đúp thanh toán cộng điểm hai lần & Điểm sai (R4). \\
Đọc điểm, cộng, ghi lại trong hai request song song & Mất một lần cộng. \\
Tính điểm trên tạm tính thay vì tổng sau giảm & Khách được điểm cho phần tiền không trả. \\
Đơn PAID rồi lỗi trước khi cộng điểm (không cùng transaction) & Đơn PAID nhưng thiếu điểm. \\
Client sửa số điểm & Gian lận điểm. \\
\end{longtable}

\textbf{Hướng xử lý/giải pháp phân tích đề xuất}

\begin{itemize}
\item
  Đặt cộng điểm bên trong hàm chuyển trạng thái PAID; chỉ chạy khi cập nhật có điều kiện theo trạng thái thành công.
\item
  Dùng phép cộng nguyên tử và lưu pointsEarned của đơn.
\item
  Test lặp bằng cách gửi cùng khóa nhiều lần và song song.
\end{itemize}

\textbf{Chuyển giao sang thiết kế}

\begin{itemize}
\item
  ERD User (loyaltyPoints), Order (pointsEarned); Sequence thanh toán (bước cộng điểm); unit test tính điểm; concurrency và idempotency test.
\end{itemize}

\textbf{Tiêu chí kiểm chứng ở mức yêu cầu}

\begin{itemize}
\item
  Đơn tổng 125.000đ thanh toán thành công cộng 12 điểm; pointsEarned của đơn bằng 12.
\item
  Đơn tổng 9.999đ cộng 0 điểm; đơn tổng 10.000đ cộng 1 điểm.
\item
  Thanh toán thất bại, đơn bị hủy hoặc hết hạn: điểm không đổi.
\item
  Gửi cùng khóa hai lần và 10 request song song: điểm chỉ tăng một lần.
\item
  Điểm tính trên tổng sau giảm giá (đơn 100.000đ giảm 20.000đ cộng 8 điểm).
\item
  Gây lỗi giữa transaction trong môi trường test: đơn không PAID và điểm không đổi.
\item
  Gửi kèm loyaltyPoints khi đăng ký hoặc thanh toán không làm đổi điểm.
\end{itemize}

\textbf{FR-PRO-004 - Giới hạn lượt dùng mã và giữ lượt khi đặt đơn}

\begin{longtable}[]{@{}
  >{\raggedright\arraybackslash}p{(\columnwidth - 2\tabcolsep) * \real{0.2500}}
  >{\raggedright\arraybackslash}p{(\columnwidth - 2\tabcolsep) * \real{0.7500}}@{}}
\toprule\noalign{}
\begin{minipage}[b]{\linewidth}\raggedright
\textbf{Thuộc tính}
\end{minipage} & \begin{minipage}[b]{\linewidth}\raggedright
\textbf{Nội dung}
\end{minipage} \\
\midrule\noalign{}
\endhead
\bottomrule\noalign{}
\endlastfoot
\textbf{Mức ưu tiên} & MUST \\
\textbf{Trạng thái} & DERIVED \\
\textbf{Yêu cầu} & Tổng số lượt sử dụng của một mã giảm giá và số lượt sử dụng của mỗi Customer không được vượt giới hạn, kể cả khi nhiều khách hàng đặt đơn đồng thời.

Lượt sử dụng mã giảm giá được giữ (HELD) khi tạo Order và tiếp tục được giữ khi Order ở trạng thái PENDING hoặc PAYMENT\_FAILED.

Khi Order chuyển sang PAID, lượt sử dụng mã được xác nhận (CONSUMED). Khi Order chuyển sang CANCELLED, lượt sử dụng mã đã giữ được trả lại (RELEASED) đúng một lần.

Việc thử lại thanh toán trên cùng Order không được giữ thêm lượt mã hoặc làm tăng số lượt sử dụng.

Nếu mã giảm giá đã hết lượt hoặc vượt giới hạn của Customer, hệ thống trả 409 COUPON\_UNAVAILABLE và không lưu thay đổi nghiệp vụ. \\
\textbf{Lý do/Nguồn} & Charter §6 (R5), §7; Techstack §7; Task 10. Cơ chế giữ và trả lượt là quyết định của nhóm, đồng nhất với FR-INV. \\
\textbf{Actor chính} & Backend API; hệ thống (tác vụ định kỳ) \\
\end{longtable}

\textbf{Phân tích nghiệp vụ}

Mã giảm giá có giới hạn lượt sử dụng là một tài nguyên cần được kiểm soát đồng thời, tương tự tồn kho. Khi nhiều Customer cùng sử dụng những lượt mã cuối cùng, hệ thống phải bảo đảm không vượt giới hạn.

Vì vậy, lượt sử dụng mã được giữ (HELD) ngay khi tạo Order thành công và tiếp tục được giữ khi Order ở trạng thái PENDING hoặc PAYMENT\_FAILED. Khi Customer thử lại thanh toán trên cùng Order, hệ thống không giữ thêm lượt mã.

Nếu thanh toán thành công và Order chuyển sang PAID, lượt mã được xác nhận sử dụng (CONSUMED). Nếu Order chuyển sang CANCELLED do Customer hủy hoặc hết hạn, lượt mã đã giữ được hoàn trả (RELEASED) đúng một lần.

Thời hạn thanh toán 15 phút giúp tránh việc lượt mã bị giữ vô thời hạn khi Customer không hoàn tất thanh toán.

\textbf{Điều kiện, dữ liệu và ràng buộc liên quan}

\begin{itemize}
\item
  Giữ lượt: khi tạo đơn có mã, trong cùng transaction với việc trừ tồn kho (FR-INV-002) và lưu đơn; dùng cập nhật có điều kiện kiểu ``tăng số lượt đã dùng nếu còn dưới giới hạn'' (nguyên tử).
\item
  Mỗi người: giới hạn lượt sử dụng mã giảm giá của từng Customer phải được kiểm tra và cập nhật nguyên tử trong cùng transaction.
\item
  Khi tính số lượt sử dụng của Customer, hệ thống phải tính cả các lượt đang giữ (HELD) trên Order ở trạng thái PENDING hoặc PAYMENT\_FAILED và các lượt đã sử dụng thành công (CONSUMED) trên Order từ PAID trở đi.
\item
  Các lượt đã hoàn trả (RELEASED) do Order chuyển sang CANCELLED không được tính vào giới hạn sử dụng.
\item
  Khi Customer thử lại thanh toán trên cùng Order, hệ thống không đếm hoặc giữ thêm lượt mã. Quy tắc mỗi Customer chỉ có tối đa một Order thuộc trạng thái PENDING hoặc PAYMENT\_FAILED được áp dụng theo FR-ORD-006.
\item
  Trả lượt: chỉ khi Order chuyển từ PENDING hoặc PAYMENT\_FAILED sang CANCELLED, hệ thống mới hoàn trả lượt sử dụng mã giảm giá đã giữ. Thao tác trả lượt phải nằm trong cùng transaction với việc chuyển trạng thái và chỉ được thực hiện đúng một lần.
\item
  Khi Order chuyển sang PAYMENT\_FAILED hoặc từ PAYMENT\_FAILED về PENDING để thử lại thanh toán, lượt sử dụng mã giảm giá tiếp tục được giữ (HELD), không hoàn trả hoặc giữ thêm.
\item
  Đơn PAID giữ lượt vĩnh viễn (không hủy được).
\item
  Lỗi tồn kho và lỗi lượt mã đều gây rollback toàn bộ; khi vừa thiếu tồn kho vừa hết lượt mã thì trả INSUFFICIENT\_STOCK (kiểm tra tồn kho trước).
\item
  Thứ tự kiểm tra trong POST /orders: 401 → dữ liệu giỏ và mã (400) → đã có đơn PENDING (409) → trong transaction: tồn kho (409 INSUFFICIENT\_STOCK) → lượt mã (409 COUPON\_UNAVAILABLE).
\item
  (SHOULD) Mỗi lần giữ hoặc trả lượt ghi một bản ghi (đơn, mã, loại) duy nhất, để đối soát như FR-INV-004.
\end{itemize}

\textbf{Vấn đề có thể gặp và phân tích tác động}

\begin{longtable}[]{@{}
  >{\raggedright\arraybackslash}p{(\columnwidth - 2\tabcolsep) * \real{0.2500}}
  >{\raggedright\arraybackslash}p{(\columnwidth - 2\tabcolsep) * \real{0.7500}}@{}}
\toprule\noalign{}
\begin{minipage}[b]{\linewidth}\raggedright
\textbf{Vấn đề}
\end{minipage} & \begin{minipage}[b]{\linewidth}\raggedright
\textbf{Tác động/Nguyên nhân}
\end{minipage} \\
\midrule\noalign{}
\endhead
\bottomrule\noalign{}
\endlastfoot
Kiểm tra lượt rồi mới tăng ở tầng ứng dụng & Đồng thời cùng qua, vượt giới hạn. \\
Không trả lượt khi đơn thất bại hoặc hủy & Mã bị đốt bởi đơn không thành công. \\
Trả lượt hai lần & Số lượt thực dùng bị đếm ít hơn thật. \\
Chỉ tính lượt khi PAID & Nhiều đơn PENDING cùng dùng lượt cuối. \\
Cộng dồn một mã nhiều lần cho cùng một người & Vượt giới hạn mỗi người. \\
Đơn PENDING bỏ dở giữ lượt & Khách khác tạm thời không dùng được (giới hạn bởi hết hạn 15 phút, FR-ORD-007). \\
\end{longtable}

\textbf{Hướng xử lý/giải pháp phân tích đề xuất}

\begin{itemize}
\item
  Cập nhật có điều kiện nguyên tử cho số lượt đã dùng, cùng transaction với tồn kho và đơn.
\item
  Gắn việc trả lượt vào đúng các hàm chuyển trạng thái đã hoàn tồn kho (FR-INV-003).
\item
  Giữ lượt theo thứ tự cố định giữa tồn kho và mã để tránh deadlock.
\end{itemize}

\textbf{Chuyển giao sang thiết kế}

\begin{itemize}
\item
  ERD Coupon (số lượt đã dùng, giới hạn), bản ghi dùng mã (đơn, người dùng, mã, trạng thái); Sequence tạo đơn, thanh toán thất bại, hủy; concurrency test.
\end{itemize}

\textbf{Tiêu chí kiểm chứng ở mức yêu cầu}

\begin{itemize}
\item
  Mã giới hạn 3 lượt; 10 khách khác nhau đặt đồng thời, mỗi người dùng mã: đúng 3 đơn thành công, 7 đơn nhận 409 COUPON\_UNAVAILABLE; số lượt đã dùng bằng 3.
\item
  Cùng khách dùng mã giới hạn 1 lượt/người cho đơn thứ hai (sau khi đơn đầu đã PAID): 409.
\item
  Khi Order sử dụng mã giảm giá chuyển sang PAYMENT\_FAILED, lượt sử dụng mã vẫn được giữ (HELD); Customer có thể thử lại thanh toán trên cùng Order mà không giữ thêm lượt mã.
\item
  Khi Order chuyển sang CANCELLED do Customer hủy hoặc hết hạn, lượt sử dụng mã đã giữ được hoàn trả (RELEASED) đúng một lần; Customer có thể sử dụng lại mã nếu vẫn đáp ứng các điều kiện hợp lệ.
\item
  Hủy hai lần hoặc tác vụ hết hạn chạy hai lần: lượt chỉ được trả một lần.
\item
  Khi thiếu tồn kho đồng thời với hết lượt mã: trả INSUFFICIENT\_STOCK, không giữ lượt.
\item
  Đơn từ chối vì lỗi bất kỳ không làm đổi số lượt đã dùng.
\end{itemize}

\textbf{FR-PRO-005 - Xem điểm tích lũy và giảm giá trên đơn}

\begin{longtable}[]{@{}
  >{\raggedright\arraybackslash}p{(\columnwidth - 2\tabcolsep) * \real{0.2500}}
  >{\raggedright\arraybackslash}p{(\columnwidth - 2\tabcolsep) * \real{0.7500}}@{}}
\toprule\noalign{}
\begin{minipage}[b]{\linewidth}\raggedright
\textbf{Thuộc tính}
\end{minipage} & \begin{minipage}[b]{\linewidth}\raggedright
\textbf{Nội dung}
\end{minipage} \\
\midrule\noalign{}
\endhead
\bottomrule\noalign{}
\endlastfoot
\textbf{Mức ưu tiên} & SHOULD \\
\textbf{Trạng thái} & DERIVED \\
\textbf{Yêu cầu} & Customer xem được số điểm hiện có của mình. Lịch sử đơn hiển thị mã giảm giá đã dùng, số tiền giảm và số điểm nhận được của từng đơn. \\
\textbf{Lý do/Nguồn} & Charter §2 (tích điểm, xem lịch sử đơn); Techstack §6 (User có loyaltyPoints). Charter không nêu cách hiển thị. \\
\textbf{Actor chính} & Customer; Web Client \\
\end{longtable}

\textbf{Phân tích nghiệp vụ}

Điểm không có tác dụng nếu khách không thấy được. Trong MVP điểm chỉ để \textbf{tích lũy và hiển thị}, chưa dùng để đổi lấy giảm giá (xem mục đối chiếu đề gốc).

\textbf{Điều kiện, dữ liệu và ràng buộc liên quan}

\begin{itemize}
\item
  Số điểm hiện có được trả bởi một endpoint chỉ đọc yêu cầu JWT (ví dụ GET /auth/me) hoặc kèm trong phản hồi đăng nhập; điểm luôn lấy từ server và chỉ của người trong token.
\item
  GET /orders/me trả thêm mã giảm giá đã dùng, số tiền giảm và pointsEarned cho từng đơn (FR-ORD-004).
\item
  Không trả điểm của người khác; không có API sửa điểm.
\item
  Điểm hiển thị có thể chậm một nhịp so với thanh toán mới; UI làm mới khi vào trang xác nhận và trang tài khoản.
\end{itemize}

\textbf{Vấn đề có thể gặp và phân tích tác động}

\begin{longtable}[]{@{}
  >{\raggedright\arraybackslash}p{(\columnwidth - 2\tabcolsep) * \real{0.2500}}
  >{\raggedright\arraybackslash}p{(\columnwidth - 2\tabcolsep) * \real{0.7500}}@{}}
\toprule\noalign{}
\begin{minipage}[b]{\linewidth}\raggedright
\textbf{Vấn đề}
\end{minipage} & \begin{minipage}[b]{\linewidth}\raggedright
\textbf{Tác động/Nguyên nhân}
\end{minipage} \\
\midrule\noalign{}
\endhead
\bottomrule\noalign{}
\endlastfoot
Hiển thị điểm lấy từ bộ nhớ client & Điểm cũ sau khi thanh toán xong. \\
Endpoint đọc điểm không lọc theo token & Lộ điểm người khác (IDOR). \\
Khách tưởng điểm đổi được giảm giá & Hiểu sai chức năng. \\
\end{longtable}

\textbf{Hướng xử lý/giải pháp phân tích đề xuất}

\begin{itemize}
\item
  Lấy điểm từ server theo token; làm mới sau khi thanh toán thành công.
\item
  Ghi rõ trên giao diện: ``Điểm tích lũy'' chưa có chức năng đổi.
\end{itemize}

\textbf{Chuyển giao sang thiết kế}

\begin{itemize}
\item
  API đọc thông tin tài khoản (bổ sung vào Techstack §6); UI hiển thị điểm; mở rộng DTO lịch sử đơn.
\end{itemize}

\textbf{Tiêu chí kiểm chứng ở mức yêu cầu}

\begin{itemize}
\item
  Sau khi thanh toán thành công, điểm hiển thị tăng đúng số điểm của đơn.
\item
  Lịch sử đơn hiển thị đúng mã, số tiền giảm và điểm nhận được.
\item
  Không đọc được điểm của người khác.
\end{itemize}

\textbf{FR-PRO-006 - Nạp dữ liệu mã giảm giá mẫu}

\begin{longtable}[]{@{}
  >{\raggedright\arraybackslash}p{(\columnwidth - 2\tabcolsep) * \real{0.2500}}
  >{\raggedright\arraybackslash}p{(\columnwidth - 2\tabcolsep) * \real{0.7500}}@{}}
\toprule\noalign{}
\begin{minipage}[b]{\linewidth}\raggedright
\textbf{Thuộc tính}
\end{minipage} & \begin{minipage}[b]{\linewidth}\raggedright
\textbf{Nội dung}
\end{minipage} \\
\midrule\noalign{}
\endhead
\bottomrule\noalign{}
\endlastfoot
\textbf{Mức ưu tiên} & MUST \\
\textbf{Trạng thái} & DERIVED \\
\textbf{Yêu cầu} & Hệ thống có dữ liệu mã giảm giá mẫu được nạp cùng dữ liệu catalog (FR-CAT-005), vì MVP không có giao diện quản trị khuyến mãi. Bộ mẫu phải đủ để kiểm thử và demo mọi nhánh của FR-PRO-001, FR-PRO-002 và FR-PRO-004. \\
\textbf{Lý do/Nguồn} & Charter §3.2 (không có Seller/Admin), §6.1 (dữ liệu mô phỏng); Techstack §11 (schema khuyến mãi khóa sau Requirements/ERD). \\
\textbf{Actor chính} & Nhóm phát triển; Backend API \\
\end{longtable}

\textbf{Phân tích nghiệp vụ}

Nếu không có nơi tạo mã thì tính năng khuyến mãi không kiểm thử và không demo được. Seed cũng cần phủ các tình huống biên, nhất là mã sắp hết lượt dùng cho kiểm thử đồng thời.

\textbf{Điều kiện, dữ liệu và ràng buộc liên quan}

\begin{itemize}
\item
  Bộ mẫu tối thiểu (tên mã minh họa): một mã PERCENT có mức giảm tối đa; một mã FIXED; một mã có đơn tối thiểu; một mã đã hết hạn; một mã chưa đến hạn; một mã giới hạn 1 lượt tổng; một mã giới hạn 3 lượt tổng để test đồng thời.
\item
  Chạy seed nhiều lần không tạo mã trùng hay đặt lại số lượt đã dùng ngoài ý muốn (có thao tác đặt lại riêng cho kiểm thử).
\item
  Danh sách mã mẫu và điều kiện được ghi trong README.
\item
  Thời hạn của mã mẫu ``còn hiệu lực'' được tính tương đối so với ngày nạp, để demo không bị hết hạn.
\end{itemize}

\textbf{Vấn đề có thể gặp và phân tích tác động}

\begin{longtable}[]{@{}
  >{\raggedright\arraybackslash}p{(\columnwidth - 2\tabcolsep) * \real{0.2500}}
  >{\raggedright\arraybackslash}p{(\columnwidth - 2\tabcolsep) * \real{0.7500}}@{}}
\toprule\noalign{}
\begin{minipage}[b]{\linewidth}\raggedright
\textbf{Vấn đề}
\end{minipage} & \begin{minipage}[b]{\linewidth}\raggedright
\textbf{Tác động/Nguyên nhân}
\end{minipage} \\
\midrule\noalign{}
\endhead
\bottomrule\noalign{}
\endlastfoot
Không có mã mẫu & Không demo và kiểm thử được khuyến mãi. \\
Mã mẫu ``còn hạn'' bị hết hạn đúng ngày nghiệm thu & Demo thất bại. \\
Seed chạy lại làm đặt lại số lượt đã dùng & Kiểm thử đồng thời sai. \\
Mã mẫu không phủ tình huống biên & Không kiểm thử được hết lượt, hết hạn. \\
\end{longtable}

\textbf{Hướng xử lý/giải pháp phân tích đề xuất}

\begin{itemize}
\item
  Script seed (qua Prisma) chạy lại an toàn; thời hạn tính theo ngày nạp.
\item
  Ghi bảng mã mẫu vào README để giảng viên tự thử.
\end{itemize}

\textbf{Chuyển giao sang thiết kế}

\begin{itemize}
\item
  Script seed/migration; README; dữ liệu cho test.
\end{itemize}

\textbf{Tiêu chí kiểm chứng ở mức yêu cầu}

\begin{itemize}
\item
  Sau khi chạy hệ thống theo README, dùng được mã mẫu hợp lệ để tạo đơn.
\item
  Có đủ mã mẫu cho mỗi nhánh của bảng lỗi ở FR-PRO-001 và FR-PRO-004.
\item
  Chạy seed hai lần không tăng số mã.
\end{itemize}

\textbf{8. Phân tích Authorization, Permission và Ownership chuyên sâu}

\textbf{8.1 Nhóm truy cập}

\begin{longtable}[]{@{}
  >{\raggedright\arraybackslash}p{(\columnwidth - 4\tabcolsep) * \real{0.3333}}
  >{\raggedright\arraybackslash}p{(\columnwidth - 4\tabcolsep) * \real{0.3333}}
  >{\raggedright\arraybackslash}p{(\columnwidth - 4\tabcolsep) * \real{0.3334}}@{}}
\toprule\noalign{}
\endhead
\bottomrule\noalign{}
\endlastfoot
\textbf{Nhóm} & \textbf{Điều kiện} & \textbf{Quyền chính} \\
Guest & Chưa đăng nhập & Xem menu, chi tiết sản phẩm, đăng ký, đăng nhập \\
Customer & Có JWT hợp lệ (còn hạn, đúng chữ ký) & Mọi quyền của Guest, cộng tạo đơn, thanh toán, hủy đơn, xem lịch sử đơn và điểm của chính mình \\
Thao tác nội bộ & Có khóa API (API Key) cấu hình trong .env, không cần JWT & Chuyển đơn đã thanh toán sang bước chuẩn bị kế tiếp (đại diện Barista khi demo) \\
\end{longtable}

Không có role, permission hay Admin. Customer không thể tự nâng quyền vì request đăng ký không nhận role (BR-007).

\textbf{8.2 Nguyên tắc}

\begin{longtable}[]{@{}
  >{\raggedright\arraybackslash}p{(\columnwidth - 4\tabcolsep) * \real{0.2200}}
  >{\raggedright\arraybackslash}p{(\columnwidth - 4\tabcolsep) * \real{0.3000}}
  >{\raggedright\arraybackslash}p{(\columnwidth - 4\tabcolsep) * \real{0.4800}}@{}}
\toprule\noalign{}
\begin{minipage}[b]{\linewidth}\raggedright
\textbf{\#}
\end{minipage} & \begin{minipage}[b]{\linewidth}\raggedright
\textbf{Nguyên tắc}
\end{minipage} & \begin{minipage}[b]{\linewidth}\raggedright
\textbf{Nguồn}
\end{minipage} \\
\midrule\noalign{}
\endhead
\bottomrule\noalign{}
\endlastfoot
1 & Mặc định yêu cầu xác thực. Chỉ các endpoint công khai trong bảng 8.3 mới mở cho Guest; endpoint mới thêm mà chưa khai báo công khai thì phải đăng nhập. & FR-AUTHZ-001; BR-005 \\
2 & Danh tính chỉ lấy từ JWT, không lấy từ body hoặc query. & FR-AUTH-003; BR-009 \\
3 & Có JWT là chưa đủ: thao tác trên một đơn cụ thể phải kiểm tra đơn thuộc về người gọi. Đơn không tồn tại và đơn của người khác cho cùng phản hồi 404. & FR-AUTHZ-002; BR-010 \\
4 & Backend kiểm tra theo thứ tự: xác thực → quyền sở hữu → quy tắc trạng thái nghiệp vụ → thực hiện. Ẩn nút ở Web Client chỉ phục vụ UX. & FR-AUTHZ-003; BR-011 \\
5 & Endpoint nội bộ tách khỏi API của Customer, bảo vệ bằng API Key, không nhận trạng thái đích. & FR-ORD-003; BR-012 \\
6 & Thông báo lỗi chung, không nêu dữ liệu hay trạng thái của tài nguyên người khác. & FR-AUTHZ-001 \\
\end{longtable}

\textbf{8.3. Bảng phân quyền endpoint}

\begin{longtable}[]{@{}
  >{\raggedright\arraybackslash}p{(\columnwidth - 12\tabcolsep) * \real{0.0744}}
  >{\raggedright\arraybackslash}p{(\columnwidth - 12\tabcolsep) * \real{0.2933}}
  >{\raggedright\arraybackslash}p{(\columnwidth - 12\tabcolsep) * \real{0.1776}}
  >{\raggedright\arraybackslash}p{(\columnwidth - 12\tabcolsep) * \real{0.0791}}
  >{\raggedright\arraybackslash}p{(\columnwidth - 12\tabcolsep) * \real{0.0962}}
  >{\raggedright\arraybackslash}p{(\columnwidth - 12\tabcolsep) * \real{0.1447}}
  >{\raggedright\arraybackslash}p{(\columnwidth - 12\tabcolsep) * \real{0.1347}}@{}}
\toprule\noalign{}
\endhead
\bottomrule\noalign{}
\endlastfoot
\textbf{Mã} & \textbf{Method và endpoint} & \textbf{Mục đích} & \textbf{Truy cập} & \textbf{Kiểm tra sở hữu} & \textbf{FR} & \textbf{Trong Techstack §6?} \\
EP-01 & GET /products & Danh sách đồ uống (phân trang tùy chọn) & Công khai & Không & FR-CAT-001 & Có \\
EP-02 & GET /products/:id & Chi tiết sản phẩm, size, topping & Công khai & Không & FR-CAT-002 & Có \\
EP-03 & POST /auth/register & Đăng ký & Công khai & Không & FR-AUTH-001 & Có \\
EP-04 & POST /auth/login & Đăng nhập, cấp JWT & Công khai & Không & FR-AUTH-002 & Có \\
EP-05 & POST /orders & Tạo đơn PENDING từ giỏ & JWT & Chủ đơn là người trong token & FR-ORD-001, FR-CART-005 & Có \\
EP-06 & GET /orders/me & Lịch sử đơn của tôi & JWT & Lọc theo userId trong token & FR-ORD-004 & Có \\
EP-07 & POST /payments & Thanh toán mock (bắt buộc Idempotency-Key) & JWT & Đơn thuộc người gọi & FR-PAY-001 & Có (cần ghi thêm JWT) \\
EP-08 & POST /orders/:id/cancel & Hủy đơn PENDING & JWT & Đơn thuộc người gọi & FR-ORD-005 & \textbf{Chưa, cần thêm} \\
EP-09 & POST /internal/orders/:id/advance & Chuyển đơn sang bước kế tiếp của chuỗi PAID → PREPARING → READY → COMPLETED & API Key nội bộ & Không (không dùng JWT) & FR-ORD-003 & \textbf{Chưa, cần thêm} \\
EP-10 & GET /auth/me & Xem thông tin tài khoản và điểm tích lũy & JWT & Người trong token & FR-PRO-005 (SHOULD) & \textbf{Chưa, cần thêm} \\
EP-11 & POST /coupons/preview & Xem trước giảm giá, không giữ lượt dùng mã & JWT & Không & FR-PRO-001 (SHOULD) & \textbf{Chưa, cần thêm} \\
\end{longtable}

Ghi chú về bảng:

\begin{itemize}
\item
  Đăng xuất không có endpoint: xóa token và giỏ hàng ở client (FR-AUTH-004).
\item
  EP-07: Idempotency-Key là header bắt buộc; thiếu thì 400 (FR-PAY-003). EP-05 chỉ chấp nhận Idempotency-Key tùy chọn (SHOULD, FR-ORD-006).
\item
  EP-08, EP-10, EP-11: tên và đường dẫn là đề xuất. Nếu nhóm bỏ hủy đơn thì xóa EP-08 và rút ``hủy'' khỏi FR-AUTHZ-002.
\item
  EP-09: khóa gửi trong header riêng (ví dụ X-Internal-Key), giá trị đặt trong .env, không commit; .env.example chỉ có giá trị mẫu.
\end{itemize}

\textbf{8.4 Kết quả theo người gọi (các endpoint cần JWT hoặc khóa)}

\begin{longtable}[]{@{}
  >{\raggedright\arraybackslash}p{(\columnwidth - 4\tabcolsep) * \real{0.4609}}
  >{\raggedright\arraybackslash}p{(\columnwidth - 4\tabcolsep) * \real{0.3454}}
  >{\raggedright\arraybackslash}p{(\columnwidth - 4\tabcolsep) * \real{0.1937}}@{}}
\toprule\noalign{}
\endhead
\bottomrule\noalign{}
\endlastfoot
\textbf{Tình huống} & \textbf{Kết quả} & \textbf{Áp dụng cho} \\
Không có token, token sai chữ ký hoặc hết hạn & 401 & EP-05 đến EP-08, EP-10, EP-11 \\
Token hợp lệ, thao tác trên đơn của chính mình & Cho phép, rồi kiểm tra trạng thái & EP-05, EP-07, EP-08 \\
Token hợp lệ, đơn của người khác hoặc không tồn tại & 404 (cùng một phản hồi) & EP-07, EP-08 \\
Đơn của mình nhưng trạng thái không cho phép (ví dụ thanh toán đơn đã PAID, hủy đơn đã PAID) & 409 INVALID\_TRANSITION & EP-07, EP-08 \\
Customer có JWT nhưng không có API Key gọi endpoint nội bộ & 403 & EP-09 \\
Không có API Key hoặc sai khóa & 403 & EP-09 \\
Có API Key, đơn không tồn tại & 404 & EP-09 \\
Có API Key, đơn không ở trạng thái cho phép bước kế tiếp (PENDING, PAYMENT\_FAILED, CANCELLED, COMPLETED) & 409 INVALID\_TRANSITION & EP-09 \\
\end{longtable}

Lưu ý: kiểm tra sở hữu luôn trước kiểm tra trạng thái, nên đơn của người khác không bao giờ nhận 409 (tránh lộ trạng thái đơn người khác). Các mã lỗi nghiệp vụ khác (400 dữ liệu giỏ hoặc mã giảm giá, 409 thiếu tồn kho, đã có đơn chờ, mã hết lượt, trùng khóa idempotency) không thuộc phân quyền và nằm ở mục 15.

\textbf{8.5 Những thứ không có (chủ ý)}

\begin{longtable}[]{@{}
  >{\raggedright\arraybackslash}p{(\columnwidth - 2\tabcolsep) * \real{0.5000}}
  >{\raggedright\arraybackslash}p{(\columnwidth - 2\tabcolsep) * \real{0.5000}}@{}}
\toprule\noalign{}
\endhead
\bottomrule\noalign{}
\endlastfoot
\textbf{Không có} & \textbf{Lý do} \\
GET /orders/:id & Hệ thống dựng lại trang xác nhận từ GET /orders/me (FR-PAY-005). \\
API sửa hoặc xóa đơn, dòng hàng, Payment & Dữ liệu giao dịch bất biến (BR-026, BR-038). \\
API quản trị sản phẩm, tồn kho, mã giảm giá, điểm & Chỉ nạp bằng dữ liệu mẫu (BR-019); điểm chỉ do backend thay đổi (BR-048). \\
API giỏ hàng & Giỏ chỉ ở client (BR-020). \\
Refresh token, đăng xuất phía server & Chỉ dùng access token ngắn hạn (FR-AUTH-004). \\
Role, permission, Admin, Seller & Ngoài phạm vi (Charter §3.2). \\
\end{longtable}

\textbf{8.6 Ma trận kiểm thử truy cập (dùng làm danh sách test)}

\begin{longtable}[]{@{}
  >{\raggedright\arraybackslash}p{(\columnwidth - 12\tabcolsep) * \real{0.1416}}
  >{\raggedright\arraybackslash}p{(\columnwidth - 12\tabcolsep) * \real{0.1416}}
  >{\raggedright\arraybackslash}p{(\columnwidth - 12\tabcolsep) * \real{0.1416}}
  >{\raggedright\arraybackslash}p{(\columnwidth - 12\tabcolsep) * \real{0.1445}}
  >{\raggedright\arraybackslash}p{(\columnwidth - 12\tabcolsep) * \real{0.1445}}
  >{\raggedright\arraybackslash}p{(\columnwidth - 12\tabcolsep) * \real{0.1445}}
  >{\raggedright\arraybackslash}p{(\columnwidth - 12\tabcolsep) * \real{0.1416}}@{}}
\toprule\noalign{}
\endhead
\bottomrule\noalign{}
\endlastfoot
\textbf{Endpoint} & \textbf{Không token} & \textbf{Token sai/hết hạn} & \textbf{Customer A trên dữ liệu của A} & \textbf{Customer A trên dữ liệu của B} & \textbf{Customer có JWT, không khóa nội bộ} & \textbf{Có khóa nội bộ} \\
EP-01, EP-02 & 200 & 200 & 200 & không áp dụng & không áp dụng & không áp dụng \\
EP-03, EP-04 & theo FR-AUTH & theo FR-AUTH & không áp dụng & không áp dụng & không áp dụng & không áp dụng \\
EP-05 & 401 & 401 & thành công & không áp dụng (đơn mới luôn thuộc người gọi) & không áp dụng & không áp dụng \\
EP-06 & 401 & 401 & chỉ đơn của A & không bao giờ thấy đơn của B & không áp dụng & không áp dụng \\
EP-07 & 401 & 401 & theo trạng thái đơn & 404 & không áp dụng & không áp dụng \\
EP-08 & 401 & 401 & theo trạng thái đơn & 404 & không áp dụng & không áp dụng \\
EP-09 & 403 & 403 & không áp dụng & không áp dụng & 403 & theo trạng thái đơn \\
EP-10, EP-11 & 401 & 401 & thành công & không áp dụng & không áp dụng & không áp dụng \\
\end{longtable}

\textbf{9. Business Rules}

\begin{longtable}[]{@{}
  >{\raggedright\arraybackslash}p{(\columnwidth - 6\tabcolsep) * \real{0.0986}}
  >{\raggedright\arraybackslash}p{(\columnwidth - 6\tabcolsep) * \real{0.4212}}
  >{\raggedright\arraybackslash}p{(\columnwidth - 6\tabcolsep) * \real{0.2420}}
  >{\raggedright\arraybackslash}p{(\columnwidth - 6\tabcolsep) * \real{0.2381}}@{}}
\toprule\noalign{}
\endhead
\bottomrule\noalign{}
\endlastfoot
\textbf{ID} & \textbf{Quy tắc} & \textbf{Nguồn} & \textbf{Loại} \\
BR-001 & Backend là nơi quyết định cuối cùng về giá, giảm giá, tổng tiền, tồn kho, kết quả thanh toán, điểm và trạng thái đơn. Web Client chỉ hiển thị và gửi yêu cầu. & Mục 3; FR-AUTHZ-003; FR-CART-005 & CONFIRMED \\
BR-002 & Tiền tệ là VND, lưu và tính bằng số nguyên hoặc Decimal, không dùng số thực dấu phẩy động. Làm tròn half-up đến đồng theo một quy tắc duy nhất; Web Client dùng cùng công thức chỉ để hiển thị tạm tính. & FR-CAT-003; FR-CART-003 & DERIVED \\
BR-003 & Request không được mang giá, đơn giá, tổng tiền, số tiền giảm, kết quả thanh toán, điểm hay role. Backend chỉ nhận các trường đã định nghĩa trong DTO; trường lạ bị bỏ qua. & FR-AUTH-001; FR-CART-005; FR-PAY-001; FR-PAY-006 & DERIVED \\
BR-004 & Một thao tác nghiệp vụ gồm nhiều thay đổi (tạo đơn; chuyển trạng thái kèm hoàn tồn, trả lượt mã, cộng điểm) được thực hiện trong một transaction; lỗi ở bất kỳ bước nào thì hoàn tác toàn bộ. & FR-ORD-001; FR-ORD-003; FR-PAY-002; Techstack §7 & DERIVED \\
BR-005 & Guest xem được danh sách và chi tiết sản phẩm, đăng ký, đăng nhập. Tạo đơn, thanh toán, xem lịch sử đơn cần JWT hợp lệ. Endpoint không khai báo công khai thì mặc định yêu cầu xác thực. & FR-AUTH-003; FR-AUTHZ-001 & CONFIRMED \\
BR-006 & Email được chuẩn hóa (cắt khoảng trắng, chữ thường) và duy nhất. Mật khẩu tối thiểu 8 ký tự, lưu bằng bcrypt, không lưu hoặc trả mật khẩu gốc. & FR-AUTH-001 & CONFIRMED (băm mật khẩu, bcrypt); DERIVED (chuẩn hóa email, tối thiểu 8 ký tự) \\
BR-007 & Đăng ký không nhận role hay điểm từ client; mọi tài khoản mới là Customer với điểm tích lũy bằng 0. & FR-AUTH-001 & DERIVED \\
BR-008 & Đăng nhập thất bại dùng một thông báo chung cho cả email không tồn tại và sai mật khẩu. Riêng đăng ký, email trùng được báo rõ. & FR-AUTH-001; FR-AUTH-002 & DERIVED \\
BR-009 & Danh tính người dùng chỉ lấy từ JWT, không lấy từ body hoặc query. & FR-AUTH-003; FR-AUTHZ-002 & DERIVED \\
BR-010 & Customer chỉ xem, thanh toán và hủy đơn của chính mình. Đơn không tồn tại và đơn của người khác cho cùng một phản hồi. & FR-AUTHZ-002 & DERIVED \\
BR-011 & Backend kiểm tra theo thứ tự: xác thực → quyền sở hữu → quy tắc trạng thái nghiệp vụ → thực hiện. Việc ẩn nút ở Web Client chỉ phục vụ UX. & FR-AUTHZ-003 & DERIVED \\
BR-012 & Các bước PAID → PREPARING → READY → COMPLETED chỉ thực hiện qua endpoint nội bộ bảo vệ bằng API Key. Endpoint không nhận trạng thái đích, mỗi lần gọi chỉ đưa đơn sang bước kế tiếp; Customer không gọi được. & FR-ORD-003; FR-AUTHZ-001 & DERIVED \\
BR-013 & Size gồm S, M, L với hệ số giá 100, 120, 140 (phần trăm), cấu hình cố định ở backend; size mặc định là S. Giá hiển thị trên menu là giá cơ bản, cũng là giá size S. & FR-CAT-001; FR-CAT-002; FR-CAT-003 & CONFIRMED (size); DERIVED (hệ số) \\
BR-014 & Đơn giá một ly = làm tròn(giá cơ bản × hệ số size ÷ 100) + tổng giá topping. Thành tiền dòng = đơn giá × số lượng. Hệ số chỉ nhân với giá cơ bản; topping là số tiền cố định. & FR-CAT-003 & CONFIRMED (tính giá); DERIVED (công thức) \\
BR-015 & Mỗi ly tối đa 3 topping, không chọn trùng. Số lượng mỗi dòng từ 1 đến 10. Topping dùng chung toàn menu và không có tồn kho riêng. & FR-CAT-002; FR-CAT-003 & DERIVED \\
BR-016 & Giá luôn lấy từ catalog tại thời điểm tạo đơn; sau khi tạo đơn giá được chốt trong bản chụp (snapshot) của dòng hàng. & FR-CAT-003; FR-ORD-002 & DERIVED \\
BR-017 & Danh sách sản phẩm sắp xếp ổn định (tên rồi mã sản phẩm). Phân trang là tùy chọn; không truyền tham số thì trả toàn bộ menu; cấu trúc phản hồi như nhau ở cả hai chế độ. & FR-CAT-001 & DERIVED \\
BR-018 & API công khai chỉ cho biết còn hay hết hàng, không trả số lượng tồn chính xác (ngoại lệ: số lượng còn có thể đặt trong lỗi thiếu hàng). & FR-CAT-004; FR-INV-001; FR-INV-002 & DERIVED \\
BR-019 & Dữ liệu sản phẩm, topping, tồn kho và mã giảm giá chỉ do dữ liệu mẫu (seed) nạp; không có giao diện quản trị. & FR-CAT-005; FR-PRO-006; Charter §3.2 & CONFIRMED \\
BR-020 & Giỏ hàng chỉ tồn tại ở client (Zustand, lưu bền ở localStorage), không có bảng hay API giỏ ở backend và không gắn với tài khoản. Giỏ chỉ lưu mã sản phẩm, size, mã topping, số lượng; là ý định mua, không phải nguồn của giá hay tồn kho. & FR-CART-001; FR-CART-004; FR-CART-005 & CONFIRMED (giỏ ở client, Zustand); DERIVED (lưu localStorage, không gắn tài khoản) \\
BR-021 & Các dòng cùng sản phẩm, cùng size, cùng tập topping (không phân biệt thứ tự) được gộp thành một dòng. Cộng dồn vượt 10 thì giữ ở 10. Giỏ tối đa 20 dòng. & FR-CART-001 & DERIVED \\
BR-022 & Giảm số lượng khi đang là 1 thì xóa dòng. Số lượng 0 không bao giờ xuất hiện trong giỏ hay trong request. & FR-CART-002 & DERIVED \\
BR-023 & Giỏ được giữ khi tải lại trang, đăng nhập, token hết hạn hoặc nhận 401, tạo đơn PENDING, thanh toán thất bại (PAYMENT\_FAILED), khách hủy đơn hoặc đơn hết hạn. Giỏ chỉ bị xóa khi thanh toán thành công (PAID) và khi đăng xuất chủ động. & FR-CART-004; FR-ORD-005; FR-AUTH-003; FR-AUTH-004 & CONFIRMED (giữ khi thanh toán lỗi); DERIVED (các trường hợp còn lại) \\
BR-024 & Thứ tự kiểm tra khi tạo đơn: xác thực → dữ liệu giỏ và mã giảm giá → kiểm tra Customer đã có Order ở trạng thái PENDING hoặc PAYMENT\_FAILED → kiểm tra tồn kho → kiểm tra lượt sử dụng mã giảm giá.

Nếu Customer đã có đơn thuộc một trong hai trạng thái trên, hệ thống trả 409 PENDING\_ORDER\_EXISTS kèm mã đơn đang chờ.

Việc trừ tồn kho, giữ lượt sử dụng mã giảm giá và lưu Order phải được thực hiện trong cùng transaction để bảo đảm tính nhất quán dữ liệu. & FR-ORD-001; FR-INV-002; FR-PRO-004 & DERIVED \\
BR-025 & Mỗi Customer chỉ được có tối đa một Order chưa hoàn tất thanh toán tại một thời điểm, bao gồm trạng thái PENDING hoặc PAYMENT\_FAILED.

Ràng buộc này phải được bảo đảm ở tầng dữ liệu hoặc bằng cơ chế transaction và khóa phù hợp, kể cả khi nhiều request tạo đơn đến đồng thời.

Khi đã có Order đang chờ, yêu cầu tạo Order mới bị từ chối với 409 PENDING\_ORDER\_EXISTS.

Customer chỉ được tạo Order mới sau khi Order đang chờ chuyển sang PAID hoặc CANCELLED. & FR-ORD-006 & DERIVED \\
BR-026 & Mỗi dòng hàng lưu bản chụp bất biến (tên món, size, topping kèm giá, đơn giá một ly, số lượng, thành tiền). Tổng đơn = tổng thành tiền các dòng − giảm giá. Không có API sửa hoặc xóa dòng hàng. & FR-ORD-002 & DERIVED \\
BR-027 & Đơn hàng chỉ được chuyển theo các trạng thái hợp lệ: Tạo mới → PENDING.PENDING → PAID, PAYMENT\_FAILED hoặc CANCELLED. PAYMENT\_FAILED → PENDING khi thử lại và đơn chưa hết hạn. PAYMENT\_FAILED → CANCELLED khi khách hủy hoặc đơn hết hạn.PAID → PREPARING → READY → COMPLETED.

COMPLETED và CANCELLED là các trạng thái kết thúc. Mọi chuyển trạng thái khác đều bị chặn. & FR-ORD-003; Techstack §7 & CONFIRMED (chuỗi PENDING đến COMPLETED, có PAYMENT\_FAILED và CANCELLED); DERIVED (các điểm đi vào CANCELLED, trạng thái cuối) \\
BR-028 & Không có thao tác gán trạng thái tùy ý. Mỗi chuyển trạng thái là một hành động riêng, thực hiện bằng cập nhật có điều kiện theo trạng thái xuất phát; hai thao tác đồng thời chỉ một thao tác thành công. Chỉ luồng thanh toán được chuyển PENDING sang PAID hoặc PAYMENT\_FAILED. & FR-ORD-003; FR-PAY-002 & DERIVED \\
BR-029 & Khách hàng chỉ được hủy đơn đang ở PENDING hoặc PAYMENT\_FAILED khi đơn chưa thanh toán thành công.

Hệ thống được tự động hủy đơn PENDING hoặc PAYMENT\_FAILED đã quá thời hạn thanh toán.

Khi hủy thành công, hệ thống hoàn tồn kho và trả lượt sử dụng mã giảm giá đã giữ đúng một lần trong cùng transaction.

Đơn từ PAID trở đi không được hủy theo nghiệp vụ MVP. & FR-ORD-005; FR-ORD-007 & DERIVED \\
BR-030 & Mã đơn là duy nhất và không phải số tăng dần dễ đoán. & FR-ORD-001 & DERIVED \\
BR-031 & Tồn kho đếm theo sản phẩm (đơn vị là ly), là số nguyên không âm, dùng chung cho mọi size và topping. Ràng buộc không âm được bảo đảm ở tầng dữ liệu. & FR-INV-001 & CONFIRMED (theo sản phẩm, không âm); DERIVED (đơn vị ly) \\
BR-032 & Tồn kho bị trừ khi tạo đơn. Số lượng các dòng cùng sản phẩm được cộng gộp trước khi so với tồn kho. Đơn chỉ có hiệu lực khi mọi sản phẩm đủ tồn (tất cả hoặc không gì cả). Việc trừ dùng cập nhật nguyên tử có điều kiện, theo thứ tự mã sản phẩm. & FR-INV-002; Charter §6 (R5) & CONFIRMED (chống oversell); DERIVED (gộp, thứ tự) \\
BR-033 & Tồn kho được giữ khi Order ở trạng thái PENDING hoặc PAYMENT\_FAILED.

Chỉ hoàn tồn kho khi Order chuyển từ PENDING hoặc PAYMENT\_FAILED sang CANCELLED. Số lượng hoàn đúng bằng số ly đã giữ khi tạo Order, lấy từ bản chụp các dòng hàng.

Việc hoàn tồn kho phải được thực hiện đúng một lần, trong cùng transaction với thao tác chuyển Order sang CANCELLED.

Khi Customer thử lại thanh toán (PAYMENT\_FAILED → PENDING), hệ thống không trừ thêm hoặc hoàn tồn kho.

Order từ PAID trở đi không được hoàn tồn kho trong phạm vi MVP. & FR-INV-003 & DERIVED \\
BR-034 & Tồn kho chỉ thay đổi do các nguyên nhân hợp lệ: khởi tạo hoặc điều chỉnh dữ liệu mẫu, tạo Order (trừ tồn kho để giữ hàng), và Order chuyển sang CANCELLED (hoàn tồn kho).

Thanh toán thất bại hoặc thử lại thanh toán không trực tiếp làm thay đổi tồn kho.

Mọi biến động tồn kho phải đi qua thành phần dịch vụ quản lý tồn kho thống nhất và có khả năng đối soát. & FR-INV-001; FR-INV-004 & DERIVED \\
BR-035 & Số tiền thanh toán luôn bằng tổng tiền của đơn do server lưu; kết quả do Mock Payment quyết định. Client không gửi số tiền hoặc kết quả. & FR-PAY-001; Mục 3 & CONFIRMED \\
BR-036 & Chỉ cho phép thanh toán đối với đơn PENDING chính mình và chưa hết hạn. Đối với đơn PAYMENT\_FAILED chưa hết hạn, Customer được phép yêu cầu thử lại thanh toán; hệ thống chuyển đơn về PENDING hợp lệ trước khi xử lý lần thanh toán mới.

Đơn đã PAID, CANCELLED, COMPLETED hoặc hết hạn phải từ chối yêu cầu thanh toán mới. Riêng yêu cầu phát lại cùng Idempotency-Key phải tuân theo BR-037. & FR-PAY-001; FR-AUTHZ-003 & DERIVED \\
BR-037 & Thanh toán bắt buộc có Idempotency-Key, duy nhất theo (người dùng, khóa). Cùng khóa và cùng nội dung thì trả lại đúng kết quả đã lưu, không tạo giao dịch mới và không đổi tồn kho hay điểm lần nữa. Cùng khóa nhưng khác nội dung thì bị từ chối. Các request đồng thời cùng khóa chỉ một request được xử lý. & FR-PAY-003; Charter §6 (R4) & CONFIRMED \\
BR-038 & Mỗi Order có thể có nhiều bản ghi Payment, tương ứng với các lần thử thanh toán khác nhau. Mỗi Payment thuộc duy nhất một Order và có một Idempotency-Key riêng.

Một Order có thể có nhiều Payment FAILED nhưng chỉ được có tối đa một Payment SUCCESS.

Trạng thái Payment (SUCCESS, FAILED) độc lập với trạng thái Order. Các bản ghi Payment đã hoàn tất không được sửa hoặc xóa qua API. & FR-PAY-004 & DERIVED \\
BR-039 & PAYMENT\_FAILED không phải trạng thái kết thúc. Customer được phép thử lại thanh toán trên cùng đơn khi đơn chưa hết hạn.

Khi thử lại, Order chuyển từ PAYMENT\_FAILED về PENDING trước khi xử lý lần thanh toán mới. Mỗi lần thử mới sử dụng một Idempotency-Key mới và tạo một bản ghi Payment mới.

Việc thử lại không tạo Order mới, không trừ tồn kho lần thứ hai và không giữ thêm lượt sử dụng mã giảm giá.

Đơn đã hết hạn hoặc chuyển sang CANCELLED không được phép thử lại thanh toán. & FR-PAY-002; FR-ORD-003 & DERIVED \\
BR-040 & Mock Payment cho kết quả xác định, không ngẫu nhiên: thẻ kết thúc bằng 0002 và ví kết thúc bằng 0000 cho FAILED; các dữ liệu hợp lệ còn lại cho SUCCESS; dữ liệu sai định dạng bị từ chối. & FR-PAY-006 & DERIVED \\
BR-041 & Hệ thống không lưu hoặc ghi log thông tin thẻ thanh toán đầy đủ; thông báo thanh toán thất bại chỉ cung cấp lý do chung, không tiết lộ thông tin nội bộ.

Khi xảy ra lỗi kỹ thuật trong quá trình xử lý thanh toán, toàn bộ transaction phải được rollback: không ghi Payment mới, không lưu Idempotency-Key của lần thử chưa hoàn tất, không cộng điểm và không thay đổi tồn kho hoặc lượt mã giảm giá.

Order phải giữ nguyên trạng thái trước khi bắt đầu lần thử: nếu ban đầu là PENDING thì vẫn PENDING; nếu ban đầu là PAYMENT\_FAILED thì vẫn PAYMENT\_FAILED.

Customer có thể gửi lại request với cùng Idempotency-Key sau khi lỗi kỹ thuật được khắc phục. & FR-PAY-001; FR-PAY-004 & DERIVED \\
BR-042 & Khi thanh toán thành công (SUCCESS), Order chuyển sang PAID, tồn kho và lượt sử dụng mã giảm giá đã giữ được xác nhận, điểm tích lũy được cộng đúng một lần và giỏ hàng được xóa.

Khi thanh toán thất bại (FAILED), Order chuyển sang PAYMENT\_FAILED, hệ thống tiếp tục giữ tồn kho và lượt mã giảm giá để hỗ trợ thử lại, không cộng điểm và giữ nguyên giỏ hàng.

Khi Order chuyển sang CANCELLED do khách hàng hủy hoặc hết hạn, hệ thống phải hoàn tồn kho và trả lượt mã giảm giá đã giữ đúng một lần. & FR-PAY-002; Charter §7 & CONFIRMED \\
BR-043 & Mỗi đơn dùng tối đa một mã giảm giá. Mã không hợp lệ thì toàn bộ đơn bị từ chối; hệ thống không tự bỏ mã để tiếp tục. & FR-PRO-001 & DERIVED \\
BR-044 & Mã hợp lệ khi tồn tại, đã đến hạn, chưa hết hạn và tạm tính của đơn đạt giá trị tối thiểu. Mã được chuẩn hóa (cắt khoảng trắng, không phân biệt hoa thường). Mã có hai loại: giảm theo phần trăm (có mức giảm tối đa) và giảm số tiền cố định. & FR-PRO-001; Techstack §7 & CONFIRMED (kiểm tra mã hợp lệ); DERIVED (điều kiện cụ thể) \\
BR-045 & Giảm giá áp lên toàn bộ tạm tính (gồm cả topping), tính ở server. Số tiền giảm không âm, không vượt tạm tính; tổng tiền không âm. Mã đã dùng và số tiền giảm được lưu vào đơn và bất biến. & FR-PRO-002 & DERIVED \\
BR-046 & Tổng lượt sử dụng mã giảm giá và số lượt sử dụng của mỗi Customer không được vượt giới hạn, kể cả khi đặt đơn đồng thời.

Lượt mã được giữ (HELD) khi tạo Order và tiếp tục được giữ trong các trạng thái PENDING hoặc PAYMENT\_FAILED.

Khi Order chuyển sang PAID, lượt mã chuyển thành CONSUMED. Khi Order chuyển sang CANCELLED, lượt mã chuyển thành RELEASED đúng một lần.

Việc thử lại thanh toán không giữ thêm lượt mã và không được làm vượt giới hạn sử dụng. & FR-PRO-004 & DERIVED \\
BR-047 & Điểm tích lũy = phần nguyên của (tổng tiền của đơn sau giảm giá ÷ 10.000). Điểm chỉ được cộng khi đơn chuyển PENDING sang PAID, đúng một lần, trong cùng transaction; không cộng cho đơn thất bại, bị hủy hay quá hạn; gửi lại thanh toán không cộng thêm. & FR-PRO-003; Techstack §7 & CONFIRMED (cộng sau PAID, không cộng trùng); DERIVED (tỷ lệ) \\
BR-048 & Điểm là số nguyên không âm và chỉ backend thay đổi. Trong MVP điểm chỉ để tích lũy và hiển thị, chưa đổi được giảm giá. & FR-PRO-003; FR-PRO-005 & DERIVED \\
BR-049 & Trạng thái còn hay hết hàng hiển thị ở Web Client chỉ mang tính tham khảo; quyết định cuối cùng là cập nhật có điều kiện khi tạo đơn. & FR-CAT-004; FR-INV-002 & DERIVED \\
BR-050 & Khi giá trị giảm giá bằng tổng tiền tạm tính, tổng tiền phải thanh toán của Order bằng 0 VND.

Order có tổng tiền 0 VND vẫn được tạo và xử lý qua Mock Payment Service theo quy trình thanh toán thông thường.

Khi Mock Payment trả kết quả SUCCESS, Order chuyển sang PAID. Điểm tích lũy được tính theo tổng tiền thực trả, do đó Order 0 VND được cộng 0 điểm.

Hệ thống không cho phép tổng tiền âm hoặc giá trị giảm giá vượt quá tổng tiền tạm tính. & FR-PRO-002; FR-PAY-001 & DERIVED \\
BR-051 & Dữ liệu giỏ đọc từ localStorage coi là không tin cậy: khi khôi phục phải kiểm tra lại cấu trúc và giới hạn, dòng không hợp lệ bị loại bỏ. & FR-CART-004 & DERIVED \\
BR-052 & JWT hết hạn sau 60 phút. JWT\_SECRET đủ dài, ngẫu nhiên, lưu trong .env, không commit lên repo; cung cấp .env.example không chứa secret. & FR-AUTH-002; Techstack §4 & CONFIRMED (không commit secret); DERIVED (thời hạn 60 phút) \\
\end{longtable}

\textbf{Mục lục:}

\begin{longtable}[]{@{}
  >{\raggedright\arraybackslash}p{(\columnwidth - 6\tabcolsep) * \real{0.1373}}
  >{\raggedright\arraybackslash}p{(\columnwidth - 6\tabcolsep) * \real{0.3106}}
  >{\raggedright\arraybackslash}p{(\columnwidth - 6\tabcolsep) * \real{0.3451}}
  >{\raggedright\arraybackslash}p{(\columnwidth - 6\tabcolsep) * \real{0.2070}}@{}}
\toprule\noalign{}
\endhead
\bottomrule\noalign{}
\endlastfoot
\textbf{Nhóm} & \textbf{Chủ đề} & \textbf{Các mã BR} & \textbf{Số quy tắc} \\
9.1 & Nguyên tắc chung & BR-001 đến BR-004 & 4 \\
9.2 & Tài khoản và quyền truy cập & BR-005 đến BR-012, BR-052 & 9 \\
9.3 & Catalog và giá & BR-013 đến BR-019, BR-049 & 8 \\
9.4 & Giỏ hàng & BR-020 đến BR-023, BR-051 & 5 \\
9.5 & Đơn hàng & BR-024 đến BR-030 & 7 \\
9.6 & Tồn kho & BR-031 đến BR-034 & 4 \\
9.7 & Thanh toán & BR-035 đến BR-042 & 8 \\
9.8 & Khuyến mãi và điểm tích lũy & BR-043 đến BR-048, BR-050 & 7 \\
& \textbf{Tổng} & \textbf{52} & \\
\end{longtable}

\textbf{10. Domain Concept và yêu cầu dữ liệu}

\begin{longtable}[]{@{}
  >{\raggedright\arraybackslash}p{(\columnwidth - 4\tabcolsep) * \real{0.4461}}
  >{\raggedright\arraybackslash}p{(\columnwidth - 4\tabcolsep) * \real{0.2056}}
  >{\raggedright\arraybackslash}p{(\columnwidth - 4\tabcolsep) * \real{0.3483}}@{}}
\toprule\noalign{}
\begin{minipage}[b]{\linewidth}\raggedright
\textbf{Domain concept}
\end{minipage} & \begin{minipage}[b]{\linewidth}\raggedright
\textbf{Ý nghĩa nghiệp vụ}
\end{minipage} & \begin{minipage}[b]{\linewidth}\raggedright
\textbf{Bất biến/Phân tích dữ liệu}
\end{minipage} \\
\midrule\noalign{}
\endhead
\bottomrule\noalign{}
\endlastfoot
User (Customer) & Danh tính nội bộ của khách; chủ sở hữu đơn và điểm tích lũy & Email duy nhất (đã chuẩn hóa); mật khẩu chỉ lưu dạng băm bcrypt; điểm tích lũy là số nguyên không âm, chỉ backend thay đổi; không có role hay permission trong cơ sở dữ liệu; họ tên là trường tùy chọn (DERIVED). \\
Product & Đồ uống bán trong menu; đơn vị bán là một ly & Giá cơ bản (cũng là giá size S) là số nguyên VND không âm; tồn kho (stock) là số nguyên không âm, bảo đảm bằng ràng buộc ở cơ sở dữ liệu; tồn kho tính theo sản phẩm, dùng chung mọi size và topping; không xóa sản phẩm trong MVP vì đơn tham chiếu tới; đổi giá hoặc tên không ảnh hưởng đơn cũ. Nếu chọn optimistic locking thì thêm cột phiên bản cho tồn kho. \\
Size (S, M, L) & Cỡ ly kèm hệ số giá 100, 120, 140 & Là cấu hình cố định ở backend (kiểu liệt kê), không phải bảng dữ liệu; dòng hàng lưu size đã chọn; hệ số chỉ nhân với giá cơ bản của món và dùng chung toàn menu. \\
Topping & Tùy chọn thêm, giá cố định, dùng chung toàn menu & Mã duy nhất; tên; giá số nguyên không âm; không có tồn kho riêng; không xóa khi đã có đơn tham chiếu; giá chỉ cộng vào đơn giá, không bị nhân hệ số. \\
Cart/CartItem & Ý định mua sắm của khách & Chỉ tồn tại ở client (Zustand, lưu localStorage), \textbf{không có bảng} ở backend; mỗi dòng gồm mã sản phẩm, size, danh sách mã topping, số lượng; không chứa giá làm căn cứ, token hay thông tin cá nhân; không phải nguồn của giá hay tồn kho. \\
Order & Giao dịch đặt đồ uống của một Customer; đơn vị của vòng đời đơn & Thuộc đúng một User; trạng thái là kiểu liệt kê 7 giá trị (PENDING, PAID, PREPARING, READY, COMPLETED, PAYMENT\_FAILED, CANCELLED); lưu tạm tính, số tiền giảm, tổng tiền (tổng = tạm tính − giảm, không âm), mã giảm giá đã dùng, điểm đã cộng (pointsEarned), thời điểm tạo, thời điểm hết hạn khi chờ thanh toán; mã đơn duy nhất và khó đoán; tối đa một đơn PENDING mỗi User, bảo đảm ở tầng dữ liệu; không xóa. Nếu làm Idempotency-Key cho tạo đơn (FR-ORD-006), lưu khóa duy nhất theo (User, khóa) trên Order. \\
OrderItem & Dòng hàng của đơn, dạng bản chụp (snapshot) & Thuộc đúng một Order; lưu mã sản phẩm để tham chiếu và các giá trị tại thời điểm đặt: tên món, size, đơn giá một ly, số lượng (1 đến 10), thành tiền dòng; bất biến, không có API sửa hoặc xóa; tổng thành tiền các dòng bằng tạm tính của đơn. \\
OrderItemTopping & Topping đã chọn của một dòng hàng, dạng bản chụp & Thuộc đúng một OrderItem; lưu mã tham chiếu, tên và giá tại thời điểm đặt; tối đa 3 dòng mỗi OrderItem và không trùng topping; bất biến; đổi hoặc xóa topping không làm đổi dữ liệu này. \\
Payment & Bản ghi giao dịch thanh toán mock & Mỗi Order có thể có nhiều Payment FAILED nhưng tối đa một Payment SUCCESS; mỗi Payment thuộc duy nhất một Order; không đặt UNIQUE trên Payment.orderId; gắn User; khóa idempotency duy nhất theo (User, khóa); lưu giá trị băm của nội dung request và dữ liệu đủ dựng lại phản hồi gốc; số tiền bằng tổng của đơn tại thời điểm ghi; phương thức (WALLET hoặc CARD); trạng thái (SUCCESS hoặc FAILED) tách khỏi trạng thái đơn; mã tham chiếu nội bộ; không lưu số thẻ đầy đủ (nếu cần chỉ 4 số cuối); bất biến sau khi ghi. \\
Coupon & Chính sách giảm giá & Mã duy nhất (đã chuẩn hóa); loại PERCENT hoặc FIXED; giá trị; giá trị đơn tối thiểu; mức giảm tối đa (chỉ với PERCENT); thời điểm bắt đầu và hết hạn; giới hạn tổng lượt dùng và lượt mỗi người (mặc định 1); số lượt đang dùng không vượt giới hạn, cập nhật nguyên tử; chỉ nạp bằng dữ liệu mẫu. \\
CouponRedemption & Lượt dùng mã gắn với một đơn & Duy nhất theo mã đơn (mỗi đơn tối đa một mã); trạng thái HELD (giữ khi tạo đơn) → CONSUMED (khi đơn PAID) hoặc RELEASED (chỉ khi Order chuyển sang CANCELLED), đúng một lần; lượt của một người = số bản ghi HELD và CONSUMED của họ \emph{(}RELEASED không tính\emph{)}. \\
Điểm tích lũy & Điểm thưởng khi thanh toán thành công & \textbf{Không có bảng riêng} trong MVP: điểm hiện có ở User, điểm của từng đơn ở Order (pointsEarned); cộng nguyên tử khi đơn chuyển PENDING sang PAID, cùng transaction; không âm; chưa có chức năng đổi điểm. \\
InventoryMovement (SHOULD) & Nhật ký biến động tồn kho do đơn hàng & Một bản ghi cho mỗi lần giữ chỗ hoặc hoàn tồn; duy nhất theo (đơn, sản phẩm, loại biến động); ghi cùng transaction với việc đổi tồn kho; không sửa hoặc xóa; chỉ phục vụ nội bộ và đối soát. \\
\end{longtable}

\textbf{10.1 Yêu cầu dữ liệu xuyên domain}

\begin{itemize}
\item
  Mọi thực thể lịch sử hoặc tài chính (Order, OrderItem, Payment, CouponRedemption, InventoryMovement) cần mã định danh ổn định và thời điểm tạo; Order có thêm thời điểm cập nhật cuối.
\item
  Tiền là số nguyên VND (kiểu số nguyên hoặc Decimal không có phần thập phân); không dùng số thực dấu phẩy động. Quy tắc làm tròn duy nhất: half-up đến đồng, tính trên đơn giá một ly (BR-002).
\item
  Mã đơn duy nhất và khó đoán; mã dùng trong API và hiển thị cho khách không phải số tăng dần.
\item
  Không xóa dữ liệu đang được Order, Payment hoặc nhật ký tham chiếu nếu việc đó làm mất lịch sử. Trong MVP không xóa sản phẩm, topping, mã giảm giá và không có API sửa hoặc xóa dòng hàng, Payment, nhật ký.
\item
  Thông tin giao dịch của đơn là bản chụp: đổi giá, tên món, topping hoặc mã giảm giá sau khi đặt không làm đổi đơn cũ (BR-016, BR-026, BR-045).
\item
  Dữ liệu cá nhân chỉ gồm email và họ tên (tùy chọn); chỉ trả cho chính chủ theo token. Không lưu số thẻ đầy đủ, mã bảo mật hay mật khẩu gốc, và không ghi các dữ liệu đó vào log.
\item
  Thời gian lưu theo UTC; so sánh thời hạn mã giảm giá và thời hạn đơn chờ theo giờ máy chủ; giao diện hiển thị theo giờ Việt Nam.
\item
  ERD phải thể hiện chuỗi sở hữu đủ để phân quyền mà không tin dữ liệu của client: Payment, OrderItem, CouponRedemption đều truy ngược được về User qua Order.
\item
  Các ràng buộc duy nhất và kiểm tra ở mục 10.2 phải được thể hiện ở tầng cơ sở dữ liệu, không chỉ ở tầng ứng dụng, vì chúng là lớp bảo vệ cuối cùng cho các tình huống đồng thời.
\item
  Trạng thái đơn và trạng thái thanh toán dùng hai kiểu liệt kê riêng; tên trạng thái trùng với Techstack §7 từng chữ.
\item
  Khóa ngoại và quan hệ chính: User 1---n Order; Order 1---n OrderItem; OrderItem 1---n OrderItemTopping; Order 1---0..N Payment; Order 1---0..1 CouponRedemption; Coupon 1---n CouponRedemption; Product 1---n OrderItem; Product 1---n InventoryMovement.
\end{itemize}

\textbf{10.2. Ràng buộc dữ liệu tối thiểu cho ERD}

\begin{longtable}[]{@{}
  >{\raggedright\arraybackslash}p{(\columnwidth - 4\tabcolsep) * \real{0.2519}}
  >{\raggedright\arraybackslash}p{(\columnwidth - 4\tabcolsep) * \real{0.5423}}
  >{\raggedright\arraybackslash}p{(\columnwidth - 4\tabcolsep) * \real{0.2057}}@{}}
\toprule\noalign{}
\endhead
\bottomrule\noalign{}
\endlastfoot
\textbf{Thực thể} & \textbf{Ràng buộc} & \textbf{Quy tắc nguồn} \\
User & Duy nhất (email) & BR-006 \\
Product & Kiểm tra (stock ≥ 0) & BR-031 \\
Order & Duy nhất (mã đơn) & BR-030 \\
Order & Tối đa một đơn thuộc \{PENDING, PAYMENT\_FAILED\} mỗi User & BR-025 \\
Order & Kiểm tra (total ≥ 0; total = subtotal − discount) & BR-045 \\
OrderItem & Kiểm tra (qty từ 1 đến 10) & BR-015 \\
OrderItemTopping & Duy nhất (OrderItem, topping) & BR-015 \\
Payment & Mỗi Payment thuộc một Order; một Order có thể có nhiều Payment & BR-038 \\
Payment & Tối đa một Payment SUCCESS trên mỗi Order & BR-038 \\
Payment & Duy nhất (User, khóa idempotency) & BR-037 \\
Coupon & Duy nhất (mã); kiểm tra (số lượt đang dùng ≤ giới hạn tổng) & BR-044, BR-046 \\
CouponRedemption & Duy nhất (mã đơn) & BR-043 \\
InventoryMovement & Duy nhất (mã đơn, sản phẩm, loại) & FR-INV-004 \\
User & Kiểm tra (điểm tích lũy ≥ 0) & BR-048 \\
Coupon & Kiểm tra (giá trị \textgreater{} 0; PERCENT thì giá trị ≤ 100; đơn tối thiểu ≥ 0; giới hạn mỗi người ≥ 1) & BR-044, BR-046 \\
Product, Topping & Kiểm tra (giá ≥ 0) & BR-002 \\
Payment & Kiểm tra (số tiền ≥ 0) & BR-050 \\
\end{longtable}

Quy tắc tối đa 3 topping mỗi dòng hàng và `số tiền Payment bằng tổng đơn' được bảo đảm ở tầng ứng dụng (BR-015, BR-035), không biểu diễn được bằng ràng buộc cơ sở dữ liệu đơn giản.

\textbf{11. Các yêu cầu về vòng đời và trạng thái}

\textbf{11.1 Các đối tượng có vòng đời}

\begin{longtable}[]{@{}
  >{\raggedright\arraybackslash}p{(\columnwidth - 6\tabcolsep) * \real{0.2688}}
  >{\raggedright\arraybackslash}p{(\columnwidth - 6\tabcolsep) * \real{0.2202}}
  >{\raggedright\arraybackslash}p{(\columnwidth - 6\tabcolsep) * \real{0.3104}}
  >{\raggedright\arraybackslash}p{(\columnwidth - 6\tabcolsep) * \real{0.2006}}@{}}
\toprule\noalign{}
\endhead
\bottomrule\noalign{}
\endlastfoot
\textbf{Đối tượng} & \textbf{Nơi quản lý} & \textbf{Có trạng thái lưu không} & \textbf{Mục chi tiết} \\
Đơn hàng (Order) & Backend & Có (7 trạng thái) & 11.2 \\
Thanh toán (Payment) & Backend & Có (SUCCESS, FAILED) & 11.3 \\
Lượt dùng mã (CouponRedemption) & Backend & Có (HELD, CONSUMED, RELEASED), đổi theo trạng thái đơn & 11.4 \\
Tồn kho, điểm tích lũy & Backend (theo trạng thái đơn) & Không có trạng thái riêng, suy ra từ đơn & 11.4 \\
Giỏ hàng & Web Client & Có (client, không lưu ở backend) & 11.7 \\
Phiên đăng nhập & Web Client (JWT) & Không lưu ở backend & 11.8 \\
\end{longtable}

Cơ chế chuyển trạng thái nội bộ, tác vụ định kỳ và các tình huống đua nhau được mô tả ở 11.5, 11.6 và 11.9.

\textbf{11.2 Vòng đời đơn hàng (Order)}

(Tạo đơn) → PENDING → PAID → PREPARING → READY → COMPLETED

PENDING → PAYMENT\_FAILED

PAYMENT\_FAILED → PENDING (Thử lại thanh toán)

PENDING → CANCELLED (Hủy hoặc quá hạn)

PAYMENT\_FAILED → CANCELLED (Hủy hoặc quá hạn)

\textbf{Bảng trạng thái}

\begin{longtable}[]{@{}
  >{\raggedright\arraybackslash}p{(\columnwidth - 6\tabcolsep) * \real{0.2738}}
  >{\raggedright\arraybackslash}p{(\columnwidth - 6\tabcolsep) * \real{0.2973}}
  >{\raggedright\arraybackslash}p{(\columnwidth - 6\tabcolsep) * \real{0.1741}}
  >{\raggedright\arraybackslash}p{(\columnwidth - 6\tabcolsep) * \real{0.2548}}@{}}
\toprule\noalign{}
\endhead
\bottomrule\noalign{}
\endlastfoot
\textbf{Trạng thái} & \textbf{Ý nghĩa} & \textbf{Cuối} & \textbf{Nhãn hiển thị cho khách} \\
PENDING & Đơn đã tạo, tồn kho đã được giữ, chờ thanh toán & Không & Chờ thanh toán \\
PAID & Thanh toán thành công & Không & Đã thanh toán \\
PREPARING & Đang pha chế & Không & Đang pha chế \\
READY & Đã pha xong, chờ khách nhận & Không & Sẵn sàng \\
COMPLETED & Hoàn tất & Có & Hoàn tất \\
PAYMENT\_FAILED & Thanh toán thất bại, khách có thể thử lại khi đơn chưa hết hạn. & Không & Thanh toán thất bại - Thử lại \\
CANCELLED & Đã hủy (khách hủy hoặc quá hạn) & Có & Đã hủy \\
\end{longtable}

COMPLETED và CANCELLED là hai trạng thái kết thúc của đơn hàng.

PAYMENT\_FAILED là trạng thái thanh toán thất bại nhưng chưa kết thúc vòng đời đơn. Customer có thể yêu cầu thử lại thanh toán trên cùng đơn khi đơn chưa hết hạn.

Khi thử lại, đơn chuyển từ PAYMENT\_FAILED về PENDING. Hệ thống không tạo đơn mới, không trừ tồn kho thêm và không giữ thêm lượt sử dụng mã giảm giá.

Nếu Customer hủy đơn hoặc đơn hết thời hạn thanh toán, hệ thống chuyển đơn từ PENDING hoặc PAYMENT\_FAILED sang CANCELLED và giải phóng các tài nguyên đã giữ đúng một lần.

\textbf{Bảng chuyển trạng thái hợp lệ}

\begin{landscape}
\begin{longtable}[]{@{}
  >{\raggedright\arraybackslash}p{(\columnwidth - 12\tabcolsep) * \real{0.0741}}
  >{\raggedright\arraybackslash}p{(\columnwidth - 12\tabcolsep) * \real{0.2011}}
  >{\raggedright\arraybackslash}p{(\columnwidth - 12\tabcolsep) * \real{0.1429}}
  >{\raggedright\arraybackslash}p{(\columnwidth - 12\tabcolsep) * \real{0.1608}}
  >{\raggedright\arraybackslash}p{(\columnwidth - 12\tabcolsep) * \real{0.1228}}
  >{\raggedright\arraybackslash}p{(\columnwidth - 12\tabcolsep) * \real{0.1481}}
  >{\raggedright\arraybackslash}p{(\columnwidth - 12\tabcolsep) * \real{0.1503}}@{}}
\toprule\noalign{}
\endhead
\bottomrule\noalign{}
\endlastfoot
\textbf{Mã} & \textbf{Từ} & \textbf{Đến} & \textbf{Sự kiện kích hoạt} & \textbf{Ai kích hoạt} & \textbf{Điều kiện (guard)} & \textbf{Tác động đi kèm (cùng transaction)} \\
TR-01 & (tạo mới) & PENDING & Tạo đơn thành công (FR-ORD-001) & Customer, qua POST /orders & Giỏ hàng hợp lệ; Customer chưa có Order ở trạng thái PENDING hoặc PAYMENT\_FAILED; đủ tồn kho và đủ lượt sử dụng mã giảm giá (nếu áp dụng). & Trừ tồn kho; giữ lượt mã (nếu có); lưu bản chụp dòng hàng; đặt thời hạn 15 phút \\
TR-02 & PENDING & PAID & Thanh toán mock thành công (FR-PAY-002) & Customer, qua POST /payments & Đơn của người gọi; đang PENDING & Ghi Payment SUCCESS; cộng điểm một lần; lưu pointsEarned; lượt dùng mã chuyển sang CONSUMED \\
TR-03 & PENDING & PAYMENT\_FAILED & Thanh toán mock thất bại (FR-PAY-002) & Customer, qua POST /payments & Order thuộc Customer đang gọi, đang PENDING, chưa hết hạn; Mock Payment trả FAILED. & Ghi Payment FAILED; chuyển Order sang PAYMENT\_FAILED; tiếp tục giữ tồn kho và lượt dùng mã giảm giá (HELD); không cộng điểm tích lũy; giữ nguyên giỏ hàng ở client. \\
TR-04 & PENDING & CANCELLED & Khách hủy (FR-ORD-005) & Customer, qua POST /orders/:id/cancel & Đơn của người gọi; đang PENDING & Hoàn tồn kho; lượt dùng mã chuyển sang RELEASED; không có Payment \\
TR-05 & PENDING & CANCELLED & Quá hạn 15 phút (FR-ORD-007) & Hệ thống (tác vụ định kỳ) & Đang PENDING và quá thời hạn & Như TR-04 \\
TR-06 & PAID & PREPARING & Thao tác nội bộ (11.5) & Đại diện Barista, qua API Key & Đang PAID & Không \\
TR-07 & PREPARING & READY & Thao tác nội bộ (11.5) & Như trên & Đang PREPARING & Không \\
TR-08 & READY & COMPLETED & Thao tác nội bộ (11.5) & Như trên & Đang READY & Không \\
TR-09 & PAYMENT\_FAILED & PENDING & Customer yêu cầu thử lại thanh toán (FR-PAY-002) & Customer, qua POST /payments & Đơn thuộc Customer; đang PAYMENT\_FAILED; chưa hết hạn; chưa có Payment SUCCESS; sử dụng Idempotency-Key mới & Chuyển Order về PENDING để thử thanh toán; không tạo Order mới, không trừ thêm tồn kho hoặc giữ thêm lượt mã \\
TR-10 & PAYMENT\_FAILED & CANCELLED & Customer hủy đơn hoặc hệ thống phát hiện đơn quá hạn & Customer qua POST /orders/:id/cancel; hoặc hệ thống qua tác vụ định kỳ & Đơn đang PAYMENT\_FAILED; thuộc Customer nếu hủy thủ công, hoặc đã hết hạn nếu tự hủy; chưa có Payment SUCCESS & Chuyển Order sang CANCELLED; hoàn tồn kho và trả lượt mã đã giữ đúng một lần trong cùng transaction; không xóa lịch sử Payment FAILED \\
\end{longtable}
\end{landscape}

Mọi chuyển trạng thái không được khai báo là hợp lệ trong ma trận ở mục 11.2.1 đều bị backend từ chối.

Hệ thống cho phép PAYMENT\_FAILED → PENDING để thử lại thanh toán trên cùng đơn khi đơn chưa hết hạn. Đây là ngoại lệ được định nghĩa rõ trong State Machine, không phải thao tác quay lùi trạng thái tùy ý.

Không cho phép PAID → CANCELLED hoặc chuyển ngược từ các trạng thái PAID, PREPARING, READY, COMPLETED về trạng thái trước đó.

\textbf{11.2.1. Ma trận đầy đủ (dùng làm bộ test)}

Mỗi ô ``Chặn'' là một test case, kỳ vọng 409 INVALID\_TRANSITION và dữ liệu không thay đổi.

\begin{landscape}
\begin{longtable}[]{@{}
  >{\raggedright\arraybackslash}p{(\columnwidth - 14\tabcolsep) * \real{0.1795}}
  >{\raggedright\arraybackslash}p{(\columnwidth - 14\tabcolsep) * \real{0.1005}}
  >{\raggedright\arraybackslash}p{(\columnwidth - 14\tabcolsep) * \real{0.0605}}
  >{\raggedright\arraybackslash}p{(\columnwidth - 14\tabcolsep) * \real{0.1877}}
  >{\raggedright\arraybackslash}p{(\columnwidth - 14\tabcolsep) * \real{0.1313}}
  >{\raggedright\arraybackslash}p{(\columnwidth - 14\tabcolsep) * \real{0.1251}}
  >{\raggedright\arraybackslash}p{(\columnwidth - 14\tabcolsep) * \real{0.0810}}
  >{\raggedright\arraybackslash}p{(\columnwidth - 14\tabcolsep) * \real{0.1344}}@{}}
\toprule\noalign{}
\endhead
\bottomrule\noalign{}
\endlastfoot
\textbf{Từ ↓ \textbackslash{} Đến →} & \textbf{PENDING} & \textbf{PAID} & \textbf{PAYMENT\_FAILED} & \textbf{CANCELLED} & \textbf{PREPARING} & \textbf{READY} & \textbf{COMPLETED} \\
(tạo mới) & TR-01 & Chặn & Chặn & Chặn & Chặn & Chặn & Chặn \\
PENDING & Chặn & TR-02 & TR-03 & TR-04, TR-05 & Chặn & Chặn & Chặn \\
PAID & Chặn & Chặn & Chặn & Chặn & TR-06 & Chặn & Chặn \\
PREPARING & Chặn & Chặn & Chặn & Chặn & Chặn & TR-07 & Chặn \\
READY & Chặn & Chặn & Chặn & Chặn & Chặn & Chặn & TR-08 \\
COMPLETED & Chặn & Chặn & Chặn & Chặn & Chặn & Chặn & Chặn \\
PAYMENT\_FAILED & TR-09 & Chặn & Chặn & TR-10 & Chặn & Chặn & Chặn \\
CANCELLED & Chặn & Chặn & Chặn & Chặn & Chặn & Chặn & Chặn \\
\end{longtable}
\end{landscape}

TR-09 -- Thử lại thanh toán: Khi đơn ở trạng thái PAYMENT\_FAILED và chưa hết hạn, Customer có thể yêu cầu thử lại thanh toán. Hệ thống chuyển đơn về PENDING, không tạo Order mới, không trừ thêm tồn kho và không giữ thêm lượt sử dụng mã giảm giá. Lần thử thanh toán mới phải dùng Idempotency-Key mới.

TR-10 -- Hủy đơn sau khi thanh toán thất bại: Khi Customer hủy đơn PAYMENT\_FAILED hoặc đơn hết thời hạn thanh toán, hệ thống chuyển đơn sang CANCELLED, hoàn tồn kho và trả lượt sử dụng mã giảm giá đã giữ đúng một lần.

Quy ước: chuyển về chính trạng thái hiện tại (ví dụ hủy đơn đã CANCELLED) cũng bị chặn 409, kể cả khi lần đầu đã thành công. Riêng POST /payments phát lại cùng khóa idempotency được xử lý trước bước kiểm tra này (xem 11.3). Việc chuyển PAYMENT\_FAILED → PENDING chỉ được phép thực hiện thông qua luồng thử lại thanh toán hợp lệ; không cho phép client trực tiếp thay đổi trạng thái đơn hàng.

\textbf{11.2.2. Quy tắc chung của mọi chuyển trạng thái}

\begin{longtable}[]{@{}
  >{\raggedright\arraybackslash}p{(\columnwidth - 4\tabcolsep) * \real{0.3333}}
  >{\raggedright\arraybackslash}p{(\columnwidth - 4\tabcolsep) * \real{0.3333}}
  >{\raggedright\arraybackslash}p{(\columnwidth - 4\tabcolsep) * \real{0.3334}}@{}}
\toprule\noalign{}
\endhead
\bottomrule\noalign{}
\endlastfoot
\textbf{Quy tắc} & \textbf{Nội dung} & \textbf{BR} \\
Không gán tùy ý & Không có hàm hay API đặt trạng thái chung; mỗi chuyển là một hành động riêng & BR-028 \\
Cập nhật có điều kiện & Chỉ ghi khi đơn đang ở đúng trạng thái xuất phát; hai thao tác đồng thời chỉ một thao tác thành công, thao tác còn lại nhận 409 & BR-028 \\
Một transaction & Đổi trạng thái và mọi tác động đi kèm nằm trong cùng transaction & BR-004 \\
Tác động đúng một lần & Hoàn tồn, trả lượt mã, cộng điểm chỉ chạy khi cập nhật có điều kiện thành công & BR-033, BR-046, BR-047 \\
Thứ tự kiểm tra & Xác thực → sở hữu → trạng thái. Đơn của người khác luôn nhận 404, không bao giờ 409 & BR-011 \\
Kiểm soát thử lại thanh toán & Chỉ cho phép PAYMENT\_FAILED → PENDING khi đơn chưa hết hạn và Customer là chủ sở hữu. Không tạo đơn mới, không trừ tồn kho hoặc giữ lượt mã thêm. Mọi lần thử thanh toán phải tuân theo idempotency. & BR-036, BR-037, BR-039 \\
Kiểm soát thời hạn đơn & Mọi yêu cầu thanh toán mới phải kiểm tra \textbf{expiresAt} ở backend. Đơn đã hết hạn không được tiếp tục thanh toán, kể cả khi tác vụ tự hủy chưa chạy. & BR-029, FR-ORD-007 \\
\end{longtable}

\textbf{11.3 Vòng đời thanh toán (Payment)}

\textbf{Sơ đồ}

(chưa có Payment) ──► SUCCESS (đơn → PAID)

│

└──────────► FAILED (đơn → PAYMENT\_FAILED)

Payment chỉ được tạo khi Mock Payment Service đã trả về kết quả SUCCESS hoặc FAILED. Hệ thống không lưu trạng thái đang xử lý (PROCESSING) trong cơ sở dữ liệu theo thiết kế MVP.

Một Order có thể có nhiều bản ghi Payment, tương ứng với những lần thử thanh toán khác nhau. Mỗi Payment thuộc về duy nhất một Order và được định danh bằng Idempotency-Key theo phạm vi người dùng.

Một Order có thể có nhiều Payment FAILED nhưng tối đa một Payment SUCCESS. Khi khách hàng thử lại thanh toán, hệ thống tạo Payment mới thay vì ghi đè kết quả của lần thanh toán thất bại trước đó.

Các bản ghi Payment đã được lưu không được thay đổi kết quả hoặc xóa thông qua API. Backend phải ngăn hai lần thanh toán thành công trên cùng một Order, kể cả khi các yêu cầu sử dụng Idempotency-Key khác nhau và được gửi đồng thời.

\textbf{Quan hệ giữa trạng thái đơn và Payment}

\begin{longtable}[]{@{}
  >{\raggedright\arraybackslash}p{(\columnwidth - 4\tabcolsep) * \real{0.4067}}
  >{\raggedright\arraybackslash}p{(\columnwidth - 4\tabcolsep) * \real{0.2600}}
  >{\raggedright\arraybackslash}p{(\columnwidth - 4\tabcolsep) * \real{0.3333}}@{}}
\toprule\noalign{}
\endhead
\bottomrule\noalign{}
\endlastfoot
\textbf{Trạng thái đơn} & \textbf{Có Payment?} & \textbf{Trạng thái Payment} \\
PENDING (mới tạo) & Không & (chưa có) \\
PENDING (sau khi tạo lại) & Có & Có một hoặc nhiều FAILED trước đó \\
PAID, PREPARING, READY, COMPLETED & Có & Có đúng một SUCCESS, có thể có FAILED trước đó \\
PAYMENT\_FAILED & Có & Có ít nhất một FAILED, chưa có SUCCESS \\
CANCELLED & Có thể có hoặc không & Chưa có hoặc có các FAILED trước đó; không có SUCCESS \\
\end{longtable}

Việc đối soát phải xét toàn bộ các bản ghi Payment thuộc cùng một Order, không chỉ bản ghi Payment mới nhất.

Order ở trạng thái PAID, PREPARING, READY hoặc COMPLETED phải có đúng một Payment SUCCESS. Order ở trạng thái PENDING, PAYMENT\_FAILED hoặc CANCELLED không được có Payment SUCCESS.

Một Order có thể có nhiều Payment FAILED. Việc tồn tại Payment FAILED trước đó không làm mất hiệu lực của Payment SUCCESS sau cùng.

\textbf{Xử lý một request POST /payments} (khớp FR-PAY-001, FR-PAY-003):

\begin{longtable}[]{@{}
  >{\raggedright\arraybackslash}p{(\columnwidth - 4\tabcolsep) * \real{0.2613}}
  >{\raggedright\arraybackslash}p{(\columnwidth - 4\tabcolsep) * \real{0.2832}}
  >{\raggedright\arraybackslash}p{(\columnwidth - 4\tabcolsep) * \real{0.4555}}@{}}
\toprule\noalign{}
\endhead
\bottomrule\noalign{}
\endlastfoot
\textbf{Bước} & \textbf{Kiểm tra} & \textbf{Nếu không đạt} \\
1 & JWT hợp lệ & 401 \\
2 & Dữ liệu request hợp lệ: Idempotency-Key, phương thức thanh toán, dữ liệu mock & 400 \\
3 & Đơn tồn tại và thuộc Customer đang gọi & 404 \\
4 & Tra Idempotency-Key trong phạm vi người dùng, trước khi kiểm tra trạng thái hoặc thời hạn cho một lần thanh toán mới & Cùng khóa và cùng nội dung: trả kết quả cũ. Cùng khóa nhưng khác nội dung: 409 IDEMPOTENCY\_KEY\_CONFLICT \\
5 & Khóa chưa dùng: xác nhận Order đang PENDING hoặc PAYMENT\_FAILED, chưa hết hạn, chưa có Payment SUCCESS; khóa xử lý đồng thời theo Order & Trạng thái không hợp lệ: 409 INVALID\_TRANSITION. Đơn hết hạn: từ chối thanh toán theo mã lỗi hết hạn đã định nghĩa \\
6 & Nếu Order là PAYMENT\_FAILED, thực hiện TR-09 để chuyển về PENDING; gọi Mock Payment, ghi Payment mới và áp dụng TR-02 hoặc TR-03 một cách nguyên tử & Lỗi kỹ thuật: hoàn tác transaction, không lưu Payment mới, không áp dụng tác động tài nguyên \\
\end{longtable}

Trạng thái ``đang xử lý'' không được lưu trong cơ sở dữ liệu. PAYMENT\_IN\_PROGRESS chỉ xuất hiện khi hai request cùng khóa đến đồng thời và request sau không lấy được khóa dòng của đơn trong thời gian chờ quy định; thông thường request sau chờ request trước xong rồi nhận kết quả đã lưu.

Ngoài các yêu cầu sử dụng cùng Idempotency-Key, hệ thống còn phải đồng bộ những yêu cầu thanh toán đồng thời sử dụng các Idempotency-Key khác nhau nhưng cùng tham chiếu một Order. Chỉ một lần thanh toán được phép chuyển Order sang PAID; những yêu cầu còn lại phải được từ chối hoặc trả kết quả phù hợp, không tạo thêm Payment SUCCESS.

\textbf{11.4 Trạng thái tài nguyên theo trạng thái đơn}

Bảng này cho thấy tồn kho, lượt dùng mã, điểm, Payment, giỏ ở client trông như thế nào ở mỗi trạng thái. Đây là công cụ đối soát khi test.

\begin{longtable}[]{@{}
  >{\raggedright\arraybackslash}p{(\columnwidth - 10\tabcolsep) * \real{0.2744}}
  >{\raggedright\arraybackslash}p{(\columnwidth - 10\tabcolsep) * \real{0.1266}}
  >{\raggedright\arraybackslash}p{(\columnwidth - 10\tabcolsep) * \real{0.1948}}
  >{\raggedright\arraybackslash}p{(\columnwidth - 10\tabcolsep) * \real{0.1477}}
  >{\raggedright\arraybackslash}p{(\columnwidth - 10\tabcolsep) * \real{0.1526}}
  >{\raggedright\arraybackslash}p{(\columnwidth - 10\tabcolsep) * \real{0.1039}}@{}}
\toprule\noalign{}
\endhead
\bottomrule\noalign{}
\endlastfoot
\textbf{Trạng thái đơn} & \textbf{Tồn kho} & \textbf{Lượt dùng mã} & \textbf{Điểm tích lũy} & \textbf{Payment} & \textbf{Giỏ ở client} \\
PENDING & Đang giữ & Đang giữ & Chưa cộng & Chưa có hoặc có FAILED trước đó & Giữ \\
PAID & Đã bán & CONSUMED & Đã cộng (pointsEarned) & Một SUCCESS, có thể có FAILED trước đó & Xóa \\
PREPARING, READY, COMPLETED & Đã bán & CONSUMED & Đã cộng & Một SUCCESS, có thể có FAILED trước đó & Không liên quan \\
PAYMENT\_FAILED & Đang giữ & Đang giữ & Chưa cộng & Có FAILED, chưa có SUCCESS & Giữ \\
CANCELLED & Đã hoàn & Đã trả (RELEASED) & Không cộng & Chưa có hoặc có FAILED trước đó & Giữ \\
\end{longtable}

Quy tắc đối soát nhanh:

Số lượng tồn kho hiện tại của một Product bằng số lượng tồn kho ban đầu, cộng các điều chỉnh tồn kho hợp lệ ngoài luồng đơn hàng (nếu có), trừ tổng số lượng sản phẩm thuộc các Order ở trạng thái PENDING, PAYMENT\_FAILED, PAID, PREPARING, READY và COMPLETED.

Order ở trạng thái PENDING hoặc PAYMENT\_FAILED đang giữ tồn kho. Order từ PAID trở đi được xem là đã bán. Order CANCELLED không còn giữ tồn kho do số lượng đã được hoàn trả đúng một lần.

Khi PAYMENT\_FAILED chuyển về PENDING để thử thanh toán, tồn kho và lượt sử dụng mã giảm giá không bị tác động thêm.

Khi PENDING hoặc PAYMENT\_FAILED chuyển sang CANCELLED, hệ thống hoàn tồn kho và trả lượt sử dụng mã giảm giá đã giữ đúng một lần. Mọi biến động phải có khả năng đối soát và không được làm tồn kho âm.

\textbf{11.5. Cơ chế chuyển trạng thái nội bộ (TR-06, TR-07, TR-08)}

\begin{longtable}[]{@{}
  >{\raggedright\arraybackslash}p{(\columnwidth - 2\tabcolsep) * \real{0.2195}}
  >{\raggedright\arraybackslash}p{(\columnwidth - 2\tabcolsep) * \real{0.7805}}@{}}
\toprule\noalign{}
\endhead
\bottomrule\noalign{}
\endlastfoot
\textbf{Mục} & \textbf{Quy định} \\
Endpoint & POST /internal/orders/:id/advance (EP-09 ở mục 8) \\
Xác thực & Header X-Internal-Key, giá trị lấy từ biến môi trường trong .env; không dùng JWT và không dùng chung với Customer; thiếu hoặc sai khóa trả 403 \\
Body & \{ "from": "\textless trạng thái hiện tại mà người gọi đang thấy\textgreater" \}; không nhận trạng thái đích \\
Hành vi & Chỉ đưa đơn sang bước kế tiếp trong chuỗi PAID → PREPARING → READY → COMPLETED; nếu trạng thái hiện tại của đơn khác from thì trả 409 và không đổi gì \\
Phản hồi & Trả mã đơn, trạng thái trước và trạng thái sau \\
Lỗi & Đơn không tồn tại: 404. Đơn ở PENDING, PAYMENT\_FAILED, CANCELLED hoặc COMPLETED, hoặc from không khớp: 409 INVALID\_TRANSITION; thiếu from hoặc from không phải một trạng thái hợp lệ: 400. \\
Ghi nhận & Ghi log mã đơn, trạng thái trước và sau, thời điểm (không ghi khóa). \\
\end{longtable}

\begin{itemize}
\item
  Lý do cần có tham số from: Nếu endpoint chỉ ``chuyển sang bước kế tiếp'' thì bấm đúp hoặc gọi lại do mạng chậm sẽ làm đơn nhảy hai bước (PAID thành READY). Yêu cầu người gọi nêu trạng thái họ đang thấy làm thao tác an toàn khi lặp lại, mà endpoint vẫn không nhận trạng thái đích tùy ý.
\item
  Giới hạn đã biết: Khóa nội bộ là một chuỗi bí mật dùng chung, không phân biệt người dùng nội bộ nào đã gọi; chấp nhận trong MVP và ghi rõ trong tài liệu bàn giao.
\end{itemize}

\textbf{11.6. Thời hạn và tác vụ định kỳ (TR-05)}

\begin{longtable}[]{@{}
  >{\raggedright\arraybackslash}p{(\columnwidth - 2\tabcolsep) * \real{0.3453}}
  >{\raggedright\arraybackslash}p{(\columnwidth - 2\tabcolsep) * \real{0.6547}}@{}}
\toprule\noalign{}
\endhead
\bottomrule\noalign{}
\endlastfoot
\textbf{Mục} & \textbf{Quy định} \\
Thời hạn đơn chờ & 15 phút kể từ lúc tạo (FR-ORD-007) \\
Tần suất quét & Mỗi phút \\
Điều kiện chọn đơn & Order thuộc trạng thái PENDING hoặc PAYMENT\_FAILED và thời điểm hiện tại đã đạt hoặc vượt expiresAt. \\
Cách xử lý & Xử lý theo lô có giới hạn; từng Order được chuyển CANCELLED bằng cập nhật có điều kiện, hoàn tồn kho và trả lượt mã trong cùng transaction. \\
Tác vụ chạy trùng hoặc chậm & Chỉ tác vụ cập nhật trạng thái thành công mới thực hiện tác động tài nguyên; những lần chạy khác bỏ qua. \\
Đơn quá hạn nhưng tác vụ chưa chạy & Từ chối mọi lần thanh toán mới; trạng thái có thể chưa chuyển CANCELLED cho đến khi tác vụ xử lý hoặc một luồng nghiệp vụ thực hiện hủy an toàn. \\
Khi đã CANCELLED & Từ chối thanh toán mới với 409 INVALID\_TRANSITION. \\
Đơn quá hạn nhưng chưa bị hủy & Có thể tạm thời ngăn tạo đơn mới do quy tắc một đơn chờ mỗi khách. Tác vụ định kỳ có trách nhiệm giải phóng đơn; thời gian xử lý thực tế phụ thuộc thời điểm và tình trạng hoạt động của tác vụ. \\
\end{longtable}

\textbf{11.7. Vòng đời giỏ hàng (client)}

\begin{longtable}[]{@{}
  >{\raggedright\arraybackslash}p{(\columnwidth - 2\tabcolsep) * \real{0.5653}}
  >{\raggedright\arraybackslash}p{(\columnwidth - 2\tabcolsep) * \real{0.4347}}@{}}
\toprule\noalign{}
\endhead
\bottomrule\noalign{}
\endlastfoot
\textbf{Sự kiện} & \textbf{Giỏ} \\
Thêm, sửa, xóa món & Thay đổi theo thao tác của khách \\
Tải lại trang, chuyển trang & Giữ \\
Guest đăng nhập thành công & Giữ \\
Token hết hạn hoặc nhận 401 & Giữ \\
Tạo đơn thành công (PENDING) & Giữ \\
Thanh toán thất bại (PAYMENT\_FAILED) & Giữ \\
Khách hủy đơn, đơn hết hạn & Giữ \\
Thanh toán thành công (PAID) & \textbf{Xóa} \\
Đăng xuất chủ động & \textbf{Xóa} \\
\end{longtable}

\textbf{11.8. Vòng đời phiên đăng nhập (JWT)}

\begin{longtable}[]{@{}
  >{\raggedright\arraybackslash}p{(\columnwidth - 2\tabcolsep) * \real{0.5000}}
  >{\raggedright\arraybackslash}p{(\columnwidth - 2\tabcolsep) * \real{0.5000}}@{}}
\toprule\noalign{}
\endhead
\bottomrule\noalign{}
\endlastfoot
\textbf{Sự kiện} & \textbf{Kết quả} \\
Đăng nhập thành công & Cấp access token, hết hạn sau 60 phút (BR-052) \\
Token còn hạn & Gọi được các endpoint cần JWT \\
Token hết hạn hoặc sai & 401; Web Client chuyển tới trang đăng nhập và giữ giỏ \\
Đăng xuất chủ động & Client xóa token và giỏ; token cũ vẫn dùng được tới khi hết hạn (không thu hồi phía server) \\
\end{longtable}

Token hết hạn trong lúc đơn đang PENDING không làm đổi trạng thái đơn: đơn tự hết hạn sau 15 phút nếu khách không đăng nhập lại để thanh toán.

\textbf{11.9. Các tình huống đua nhau}

\begin{longtable}[]{@{}
  >{\raggedright\arraybackslash}p{(\columnwidth - 4\tabcolsep) * \real{0.3161}}
  >{\raggedright\arraybackslash}p{(\columnwidth - 4\tabcolsep) * \real{0.3657}}
  >{\raggedright\arraybackslash}p{(\columnwidth - 4\tabcolsep) * \real{0.3182}}@{}}
\toprule\noalign{}
\endhead
\bottomrule\noalign{}
\endlastfoot
\textbf{Tình huống} & \textbf{Kết quả mong đợi} & \textbf{Cơ chế} \\
Hai request tạo đơn cùng lúc của một khách & Một đơn, request còn lại 409 PENDING\_ORDER\_EXISTS & Ràng buộc tối đa một Order thuộc trạng thái PENDING hoặc PAYMENT\_FAILED cho mỗi Customer ở tầng dữ liệu, kết hợp transaction và kiểm soát đồng thời (BR-025). \\
10 khách đặt đồng thời món còn 5 & Đúng 5 đơn thành công, tồn không âm & Cập nhật nguyên tử có điều kiện (BR-032) \\
Thanh toán cùng một đơn hai lần với cùng khóa & Một Payment, kết quả giống nhau, điểm cộng một lần & Khóa idempotency (BR-037) \\
Hai request thanh toán cùng Order nhưng dùng hai Idempotency-Key khác nhau & Tối đa một Payment SUCCESS; Order chỉ chuyển PAID một lần; điểm chỉ được cộng một lần & Đồng bộ theo Order, kiểm tra và cập nhật trạng thái nguyên tử; ràng buộc tối đa một Payment SUCCESS (BR-028, BR-038) \\
Customer thử lại thanh toán khi Order đang PAYMENT\_FAILED & Không tạo Order mới; không trừ thêm tồn kho hoặc giữ thêm lượt mã; tạo Payment mới nếu xử lý có kết quả & Kiểm tra sở hữu, thời hạn, Idempotency-Key mới và chuyển trạng thái hợp lệ TR-09 (BR-039) \\
Thanh toán và khách hủy cùng lúc & Chỉ một thao tác thắng; trạng thái cuối nhất quán với tồn kho, lượt mã và điểm & Cập nhật có điều kiện theo trạng thái (BR-028) \\
Thanh toán và hết hạn cùng lúc & Như trên & Như trên \\
Bấm hủy hai lần & Lần hai 409; hoàn tồn đúng một lần & Cập nhật có điều kiện (BR-033) \\
Tác vụ hết hạn chạy hai lần & Hoàn tồn đúng một lần & Như trên \\
Gọi endpoint nội bộ hai lần với cùng from & Lần hai 409, không nhảy hai bước & Tham số from (11.5) \\
\end{longtable}

\textbf{12. External Integration Requirements}

\textbf{12.1 Mock Payment Service}

\begin{itemize}
\item
  Là module nội bộ trong backend NestJS, được gọi đồng bộ trong cùng tiến trình. Không có mạng ngoài, không có nhà cung cấp thật, không có redirect hay callback (Charter §3.2; Techstack §5).
\item
  Đầu vào: mã đơn, số tiền \textbf{do server lấy từ đơn}, phương thức (WALLET hoặc CARD) và dữ liệu test. Không nhận số tiền hay kết quả từ client (BR-035).
\item
  Đầu ra: kết quả SUCCESS hoặc FAILED, mã tham chiếu nội bộ và lý do thất bại ở mức chung.
\item
  Kết quả xác định theo dữ liệu test (thẻ kết thúc 0002 và ví kết thúc 0000 cho FAILED), không dùng yếu tố ngẫu nhiên, để kiểm thử và demo lặp lại được (BR-040, FR-PAY-006).
\item
  Vì gọi đồng bộ nên không có callback trễ, lặp hoặc sai thứ tự như với cổng thanh toán thật. Việc chống giao dịch lặp do Idempotency-Key của backend đảm nhiệm (FR-PAY-003).
\item
  Backend gọi Mock Payment qua một ranh giới rõ ràng (interface hoặc service riêng), để sau này thay bằng nhà cung cấp thật mà không phải sửa logic đơn hàng, tồn kho và điểm.
\item
  Phân biệt kết quả nghiệp vụ FAILED (thanh toán bị từ chối) với lỗi kỹ thuật của mock (ngoại lệ). Lỗi kỹ thuật làm rollback toàn bộ transaction; Order giữ nguyên trạng thái trước lần thử thanh toán (PENDING hoặc PAYMENT\_FAILED), không tạo Payment mới và không lưu Idempotency-Key của lần thử chưa hoàn tất (BR-041).
\item
  Việc gây lỗi kỹ thuật để kiểm thử chỉ bật được bằng cấu hình môi trường test, không qua API công khai.
\item
  Không ghi log số thẻ hay số ví đầy đủ; không lưu số thẻ đầy đủ (FR-PAY-004).
\end{itemize}

\textbf{12.2 Các tích hợp không có trong BrewLite}

\begin{itemize}
\item
  Không có đăng nhập bằng bên thứ ba (Google và tương tự); chỉ đăng ký, đăng nhập bằng email và mật khẩu.
\item
  Không có cổng thanh toán thật (VNPay, Momo, ví điện tử, thẻ).
\item
  Không có hệ thống vận chuyển hay định vị.
\item
  Không có email, SMS hay thông báo đẩy.
\item
  Không có thành phần hạ tầng bổ sung ngoài PostgreSQL: không API gateway, event bus hay dịch vụ độc lập (Techstack §5); không dùng Redis hay công cụ tìm kiếm riêng
\end{itemize}

\textbf{13. Yêu cầu phi chức năng (NFR)}

\begin{longtable}[]{@{}
  >{\raggedright\arraybackslash}p{(\columnwidth - 4\tabcolsep) * \real{0.2200}}
  >{\raggedright\arraybackslash}p{(\columnwidth - 4\tabcolsep) * \real{0.3000}}
  >{\raggedright\arraybackslash}p{(\columnwidth - 4\tabcolsep) * \real{0.4800}}@{}}
\toprule\noalign{}
\begin{minipage}[b]{\linewidth}\raggedright
\textbf{ID}
\end{minipage} & \begin{minipage}[b]{\linewidth}\raggedright
\textbf{Nhóm}
\end{minipage} & \begin{minipage}[b]{\linewidth}\raggedright
\textbf{Yêu cầu}
\end{minipage} \\
\midrule\noalign{}
\endhead
\bottomrule\noalign{}
\endlastfoot
NFR-SEC-001 & Security & Mọi endpoint cần đăng nhập phải xác thực JWT và kiểm tra quyền sở hữu ở server; danh tính chỉ lấy từ token (FR-AUTH-003, FR-AUTHZ-002). \\
NFR-SEC-002 & Security & Mặc định yêu cầu xác thực; endpoint công khai phải được khai báo tường minh (mục 8). \\
NFR-SEC-003 & Security & Mật khẩu lưu bằng bcrypt; không lưu hoặc log mật khẩu gốc (Charter §7; Techstack §3). \\
NFR-SEC-004 & Security & JWT\_SECRET, X-Internal-Key và thông tin kết nối cơ sở dữ liệu nằm trong .env, không commit lên repo; .env.example không chứa secret (Techstack §4). \\
NFR-SEC-005 & Security & Kiểm tra dữ liệu đầu vào ở mọi endpoint bằng DTO và danh sách trường được phép; chống mass assignment (Charter §7; Techstack §3). \\
NFR-SEC-006 & Security & Endpoint nội bộ chỉ chấp nhận X-Internal-Key, không dùng chung với JWT của Customer; thiếu hoặc sai khóa trả 403 (mục 8, 11.5). \\
NFR-SEC-007 & Security & Không lưu số thẻ đầy đủ (FR-PAY-004) \\
NFR-SEC-008 & Security & (SHOULD) Giới hạn số lần thử đăng nhập sai nếu thời gian cho phép (FR-AUTH-002). \\
NFR-SEC-009 & Security & Kiểm thử âm tính truy cập chéo bắt buộc: Customer A không đọc, thanh toán hoặc hủy được đơn của Customer B (mục 8.6). \\
NFR-DATA-001 & Data integrity & Tạo đơn, chuyển trạng thái, thanh toán, hoàn tồn, trả lượt mã và cộng điểm dùng transaction và cơ chế đồng thời bảo toàn bất biến (Charter §6 R5; Techstack §7; BR-004). \\
NFR-DATA-002 & Data integrity & Đơn giữ bản chụp dòng hàng; thay đổi catalog hoặc mã giảm giá không sửa lịch sử đơn (BR-016, BR-026). \\
NFR-DATA-003 & Data integrity & Tiền là số nguyên VND hoặc Decimal không có phần thập phân; quy tắc làm tròn duy nhất; không dùng số thực dấu phẩy động (BR-002). \\
NFR-DATA-004 & Data integrity & Các ràng buộc duy nhất và kiểm tra ở mục 10.2 phải có ở tầng cơ sở dữ liệu, không chỉ ở tầng ứng dụng. \\
NFR-DATA-005 & Data integrity & Dữ liệu mẫu tái tạo được bằng script seed chạy lại an toàn (FR-CAT-005, FR-PRO-006). Sao lưu và khôi phục dữ liệu nằm ngoài phạm vi MVP. \\
NFR-REL-001 & Reliability & Request bị gửi lặp (bấm đúp, thử lại khi mạng chậm, tác vụ chạy trùng) không tạo thêm đơn, giao dịch, lần hoàn tồn hay lần cộng điểm (Charter §6 R4; FR-PAY-003; BR-028). \\
NFR-REL-002 & Reliability & Lỗi giữa chừng không được để lại dữ liệu nửa vời. Toàn bộ transaction phải được rollback và Order giữ nguyên trạng thái trước lần thử thanh toán, bao gồm PENDING hoặc PAYMENT\_FAILED (BR-041). \\
NFR-REL-003 & Reliability & Giỏ hàng ở client có thể lỗi thời hoặc bị sửa; backend luôn kiểm tra lại giỏ khi tạo đơn (FR-CART-005). \\
NFR-REL-004 & Reliability & (SHOULD) Tác vụ hết hạn đơn chờ chạy mỗi phút, an toàn khi chạy trùng hoặc chậm (FR-ORD-007). \\
NFR-OBS-001 & Observability & (SHOULD) Backend ghi log lỗi và các chuyển trạng thái đơn quan trọng, kèm mã đơn, trạng thái trước và sau, thời điểm (11.5). \\
NFR-OBS-002 & Observability & Log không chứa mật khẩu, token, X-Internal-Key, dữ liệu thẻ hoặc thông tin cá nhân không cần thiết. \\
NFR-PERF-001 & Performance & Các luồng chính (danh sách, chi tiết, tạo đơn, thanh toán, lịch sử đơn) phản hồi dưới 500 ms với dữ liệu mẫu (Charter §7). \\
NFR-PERF-002 & Performance & Danh sách sản phẩm hỗ trợ phân trang tùy chọn, tối đa 50 phần tử mỗi trang (FR-CAT-001); (SHOULD) lịch sử đơn cùng quy ước (FR-ORD-004). \\
NFR-PERF-003 & Performance & (SHOULD) Tránh truy vấn lặp (N+1) khi dựng danh sách sản phẩm và lịch sử đơn. \\
NFR-PERF-004 & Performance & Ngoài mục tiêu 500 ms của Charter, không cam kết thêm mức tải hay thời gian đáp ứng khi chưa đo; số yêu cầu đồng thời dùng khi kiểm thử là tham số kiểm thử, không phải cam kết tải. \\
NFR-UX-001 & Usability & Giao diện có trạng thái đang tải, trống và lỗi cho menu, giỏ hàng, lịch sử đơn và thanh toán (Techstack §5, §10). \\
NFR-UX-002 & Usability & Lỗi trả kèm mã máy đọc được và thông điệp đủ để giao diện hướng dẫn khách sửa (ví dụ INVALID\_CART, INSUFFICIENT\_STOCK, COUPON\_UNAVAILABLE). \\
NFR-UX-003 & Usability & Nếu giá hoặc tồn kho thay đổi so với lúc khách xem, giao diện phải hiển thị lại tổng cuối do server trả về trước bước thanh toán, không thu số tiền khác với số khách thấy (FR-CAT-003, FR-CART-003). \\
NFR-COMP-001 & Compatibility & Hệ thống là ứng dụng web (REST/JSON); ứng dụng mobile native nằm ngoài phạm vi (Charter §3.2). \\
NFR-MNT-001 & Maintainability & Quy tắc nghiệp vụ cốt lõi chỉ được thực thi ở backend; frontend chỉ hiển thị tạm tính theo cùng công thức (BR-001). \\
NFR-MNT-002 & Maintainability & Hợp đồng API (DTO, mã lỗi) thống nhất giữa frontend và backend; thay đổi hợp đồng phải cập nhật tài liệu và test (Charter §6 R3). \\
NFR-MNT-003 & Maintainability & Một backend duy nhất chia module; không microservices, API gateway hay event bus (Techstack §1, §5); không thêm Redis hoặc công cụ tìm kiếm riêng khi chưa có nhu cầu thực tế. \\
NFR-MNT-004 & Maintainability & Tuân theo workflow và chuẩn kỹ thuật của Techstack §9: ESLint và Prettier, PR có người review, pull request đổi nghiệp vụ phải cập nhật test và tài liệu. \\
NFR-TEST-001 & Testability & Mỗi yêu cầu MUST có kiểm thử luồng chính, luồng thay thế và kiểm thử âm tính; thanh toán, tồn kho, phân quyền và chuyển trạng thái cần integration test (Techstack §10). \\
NFR-TEST-002 & Testability & Concurrency test và idempotency test bắt buộc: cùng khóa gửi hai lần chỉ một kết quả; nhiều yêu cầu đồng thời không làm tồn kho hay lượt dùng mã vượt giới hạn (Charter §7; Techstack §10). \\
NFR-TEST-003 & Testability & Unit test cho tính giá size và topping, tổng đơn, giảm giá, điểm tích lũy và ma trận chuyển trạng thái \textbf{(11.2.1)} (Techstack §10). \\
NFR-TEST-004 & Testability & Dữ liệu test cố định của Mock Payment cho cả hai nhánh thành công và thất bại, ghi trong README (FR-PAY-006). \\
NFR-TEST-005 & Testability & Có bộ Postman cho các endpoint ở mục 8 để kiểm thử thủ công và demo (Techstack §3). \\
NFR-TEST-006 & Testability & (SHOULD) CI chạy lint, test và build khi tạo hoặc cập nhật pull request (Techstack §3). \\
NFR-PRIV-001 & Privacy & Chỉ lưu và trả thông tin cá nhân cần thiết (email, họ tên tùy chọn); chỉ trả cho chính chủ theo token. \\
NFR-OPS-001 & Operability & Docker Compose chạy đủ frontend, backend và cơ sở dữ liệu bằng một lệnh; README hướng dẫn chạy, nạp dữ liệu mẫu và các dữ liệu test (Charter §7). \\
NFR-OPS-002 & Operability & Ghim phiên bản: commit lockfile, ghi tag ảnh cơ sở dữ liệu cụ thể, không dùng latest (Techstack §4). \\
NFR-OPS-003 & Operability & Các giá trị cấu hình (thời hạn JWT 60 phút, thời hạn đơn chờ 15 phút, khóa nội bộ, kết nối cơ sở dữ liệu) đặt ở cấu hình môi trường, không gắn cứng trong mã. Hệ số size là hằng số cố định trong mã, không thuộc cấu hình môi trường. \\
\end{longtable}

\textbf{14. Phân tích rủi ro trọng yếu và hướng xử lý}

\begin{longtable}[]{@{}
  >{\raggedright\arraybackslash}p{(\columnwidth - 8\tabcolsep) * \real{0.0698}}
  >{\raggedright\arraybackslash}p{(\columnwidth - 8\tabcolsep) * \real{0.1663}}
  >{\raggedright\arraybackslash}p{(\columnwidth - 8\tabcolsep) * \real{0.2669}}
  >{\raggedright\arraybackslash}p{(\columnwidth - 8\tabcolsep) * \real{0.1047}}
  >{\raggedright\arraybackslash}p{(\columnwidth - 8\tabcolsep) * \real{0.3923}}@{}}
\toprule\noalign{}
\begin{minipage}[b]{\linewidth}\raggedright
\textbf{ID}
\end{minipage} & \begin{minipage}[b]{\linewidth}\raggedright
\textbf{Rủi ro}
\end{minipage} & \begin{minipage}[b]{\linewidth}\raggedright
\textbf{Phân tích}
\end{minipage} & \begin{minipage}[b]{\linewidth}\raggedright
\textbf{Mức}
\end{minipage} & \begin{minipage}[b]{\linewidth}\raggedright
\textbf{Hướng xử lý}
\end{minipage} \\
\midrule\noalign{}
\endhead
\bottomrule\noalign{}
\endlastfoot
RA-01 & Truy cập chéo dữ liệu (IDOR) & Customer có JWT hợp lệ đổi mã đơn trong request để xem, thanh toán hoặc hủy đơn của người khác. Là rủi ro bảo mật lớn nhất của hệ thống một cửa hàng. & Critical & Danh tính chỉ lấy từ JWT; truy vấn đơn kèm userId; đơn người khác và đơn không tồn tại cùng trả 404; kiểm thử âm tính bắt buộc (FR-AUTHZ-002, mục 8.6, NFR-SEC-009). \\
RA-02 & Bán vượt tồn kho (Charter R5) & Nhiều khách đặt đồng thời món cuối cùng, hoặc một đơn có hai dòng cùng sản phẩm mà mỗi dòng riêng đều đủ nhưng tổng vượt tồn. & Critical & Cập nhật nguyên tử có điều kiện (hoặc optimistic locking) trong transaction; gộp số lượng theo sản phẩm; trừ theo thứ tự mã sản phẩm; ràng buộc stock ≥ 0 ở cơ sở dữ liệu; concurrency test (FR-INV-002, BR-032, NFR-TEST-002). \\
RA-03 & Giao dịch và tác động phụ bị lặp (Charter R4) & Bấm đúp, gửi lại khi mạng chậm hoặc request song song làm thanh toán, hoàn tồn kho, trả lượt mã hoặc cộng điểm xảy ra hai lần. & Critical & Idempotency-Key duy nhất theo (người dùng, khóa); cùng khóa khác nội dung thì từ chối; Mỗi Order có thể có nhiều Payment FAILED nhưng tối đa một Payment SUCCESS ; tác động chỉ chạy khi cập nhật trạng thái thành công (FR-PAY-003, FR-PAY-004, BR-037, BR-038). \\
RA-04 & Trạng thái đơn sai do tranh chấp & Hủy, hết hạn và thanh toán xảy ra gần như đồng thời làm đơn vừa CANCELLED vừa PAID, hoặc tồn kho và điểm lệch với trạng thái đơn. & Critical & Mỗi chuyển trạng thái là một hành động riêng, cập nhật có điều kiện theo trạng thái xuất phát, cùng một transaction với tác động đi kèm; chặn mọi chuyển ngoài ma trận; test ma trận và tình huống đua nhau (FR-ORD-003, mục 11.2.1, 11.9, BR-027, BR-028). \\
RA-05 & Tin dữ liệu từ client & Client gửi giá, tổng tiền, số tiền giảm, điểm, role hoặc kết quả thanh toán; nếu backend tin thì khách tự giảm giá hoặc tự báo thành công. & Critical & Backend tính lại mọi giá trị tiền từ catalog; DTO chỉ nhận trường được phép; kết quả thanh toán do Mock Payment quyết định; kiểm thử gửi kèm trường giả (BR-001, BR-003, FR-CART-005, FR-PAY-001). \\
RA-06 & Lộ bí mật và dữ liệu nhạy cảm & JWT\_SECRET, X-Internal-Key hoặc thông tin cơ sở dữ liệu bị commit hoặc ghi log; số thẻ, token, mật khẩu xuất hiện trong log hoặc cơ sở dữ liệu. & Critical & Bí mật chỉ nằm trong .env, có .env.example không chứa secret; mật khẩu bcrypt; không lưu số thẻ đầy đủ; không ghi log dữ liệu nhạy cảm; (SHOULD) quét secret khi tạo pull request (NFR-SEC-003, 004, 007, NFR-OBS-002). \\
RA-07 & Lạm dụng endpoint nội bộ & Endpoint chuyển trạng thái đơn bị Customer gọi được hoặc khóa dùng chung bị lộ, khiến đơn bị đẩy sang COMPLETED hoặc nhảy bước. & High & Chỉ chấp nhận X-Internal-Key, không dùng JWT; thiếu hoặc sai khóa trả 403; không nhận trạng thái đích, chỉ nhận from; ghi log mỗi lần gọi; giới hạn đã biết (khóa dùng chung) được ghi rõ trong tài liệu bàn giao (11.5, NFR-SEC-006). \\
RA-08 & Lịch sử đơn bị thay đổi & Đổi giá, tên món, topping hoặc mã giảm giá làm đơn cũ hiển thị sai; xóa sản phẩm làm mất dữ liệu đơn. & High & Bản chụp dòng hàng và topping bất biến; không xóa sản phẩm, topping, mã giảm giá trong MVP; không có API sửa dòng hàng, Payment (FR-ORD-002, FR-PAY-004, BR-026). \\
RA-09 & Vượt giới hạn hoặc đốt lượt mã giảm giá & Đặt đồng thời làm mã vượt giới hạn lượt dùng; hoặc đơn thất bại, bị hủy không trả lượt khiến mã hết dù chưa ai dùng thật. & High & Cập nhật nguyên tử có điều kiện để kiểm soát giới hạn lượt sử dụng mã, trong cùng transaction với thao tác tạo Order và giữ tồn kho.

Lượt mã tiếp tục được giữ (HELD) khi Order ở trạng thái PENDING hoặc PAYMENT\_FAILED; không giữ thêm lượt khi Customer thử lại thanh toán.

Chỉ hoàn trả lượt mã (RELEASED) đúng một lần khi Order chuyển sang CANCELLED. Khi Order chuyển sang PAID, lượt mã được xác nhận sử dụng (CONSUMED).

Kiểm thử các tình huống đặt đơn đồng thời, thanh toán thất bại, thử lại thanh toán và hủy đơn (FR-PRO-004, BR-046). \\
RA-10 & Điểm tích lũy sai & Cộng điểm khi đơn chưa PAID, cộng hai lần khi thanh toán lặp, hoặc tính trên tạm tính thay vì tổng sau giảm. & High & Cộng điểm nguyên tử trong cùng transaction với chuyển PENDING → PAID; tính trên tổng sau giảm; lưu pointsEarned trên đơn; test lặp khóa và đồng thời (FR-PRO-003, BR-047). \\
RA-11 & Lệch hợp đồng và quy tắc giữa frontend và backend (Charter R3) & Frontend tính giá tạm tính khác công thức backend, hoặc hai bên hiểu khác về DTO và mã lỗi, gây lỗi khi tích hợp hoặc demo. & High & Backend là nguồn quyết định; frontend dùng cùng công thức và cùng bộ ví dụ test chỉ để hiển thị; thống nhất DTO và mã lỗi; .env.example; kiểm thử API sớm bằng Postman (NFR-MNT-001, NFR-MNT-002, NFR-TEST-005). \\
RA-12 & Tài liệu lệch nhau (Charter R2) & Requirements, ERD, Techstack và API không khớp (tên endpoint, cột dữ liệu, tên trạng thái), làm thiết kế và code sai lệch. & High & Review Requirements trước khi chốt API và ERD; đồng bộ Techstack §6 với mục 8 và mục 10; dùng mục 8, 9, 11 làm nguồn duy nhất; kiểm tra tên trạng thái trùng Techstack §7 từng chữ. \\
RA-13 & Giả định của nhóm lệch đề bài & Nhiều quy tắc ở mục 7 là quyết định của nhóm (DERIVED), như trừ tồn khi tạo đơn, không hủy sau thanh toán, công thức giảm giá, tỷ lệ tích điểm. Nếu khác đề thì mất điểm hoặc phải làm lại. & High & Đối chiếu từng quy tắc DERIVED có ghi ``cần đối chiếu'' với đề BrewLite VER 1.0 trước khi chốt bản baseline; hỏi giảng viên khi đề chưa rõ; ghi quyết định cuối vào mục 9. \\
RA-14 & Scope creep (Charter R1) & Thêm đổi điểm lấy giảm giá, hoàn tiền, hủy sau thanh toán, giao diện Barista, đăng nhập bên thứ ba làm trễ 10 task bắt buộc. & Medium & Khóa MVP theo đề; mục 4.2 liệt kê rõ ngoài phạm vi; yêu cầu mới đi qua quy tắc kiểm soát thay đổi (Charter §3.3); các FR SHOULD chỉ làm khi MUST xong. \\
RA-15 & Dồn nghiệp vụ nâng cao vào cuối kỳ (Charter R6) & Task 10 (state machine, idempotency, đồng thời, khuyến mãi, điểm) tập trung ở Sprint 3, không còn thời gian kiểm thử. & High & Thiết kế trạng thái đơn và tồn kho từ Sprint 1 và 2; dùng mục 10, 11 làm đầu vào cho ERD; viết concurrency test từ khi làm tạo đơn; dành thời gian đệm cho Sprint 3. \\
RA-16 & Đơn chờ bỏ dở và giỏ hàng lỗi thời & Đơn PENDING bỏ dở giữ tồn kho và lượt mã, làm khách khác thấy hết hàng và khách chính bị kẹt vì quy tắc một đơn chờ; giỏ ở client cũ hoặc bị sửa làm tạo đơn thất bại. & Medium & Khách tự hủy đơn chờ, hệ thống tự hủy sau 15 phút; backend luôn kiểm tra lại giỏ, lỗi trả theo dòng, UI giữ giỏ để khách sửa; chấp nhận hạn chế đơn quá hạn chưa bị hủy tối đa khoảng một phút (FR-ORD-005, FR-ORD-007, FR-CART-005, 11.6). \\
RA-17 & Môi trường demo không tái tạo được & Không chạy được bằng Docker Compose, thiếu dữ liệu mẫu, phiên bản công nghệ lệch giữa máy làm demo không mong muốn, Mock Payment cho kết quả khác lần trước. & Medium & Docker Compose chạy đủ ba dịch vụ; seed chạy lại an toàn có mã giảm giá và tồn kho biên; ghim phiên bản và lockfile; Mock Payment cho kết quả xác định; README ghi dữ liệu test; chạy thử từ máy sạch trước nghiệm thu (FR-CAT-005, FR-PRO-006, FR-PAY-006, NFR-OPS-001, NFR-OPS-002). \\
\end{longtable}

\textbf{Thang mức:} Critical là rủi ro làm sai tiền, tồn kho hoặc quyền truy cập nên phải có test tự động; High là rủi ro gây sai dữ liệu phụ, lệch tài liệu hoặc mất điểm; Medium là rủi ro chủ yếu ảnh hưởng tiến độ và trải nghiệm.

\textbf{15. Ma trận lỗi, trường hợp ngoại lệ và sự cố}

\textbf{15.1. Bảng mã lỗi tập trung}

\begin{longtable}[]{@{}
  >{\raggedright\arraybackslash}p{(\columnwidth - 6\tabcolsep) * \real{0.1637}}
  >{\raggedright\arraybackslash}p{(\columnwidth - 6\tabcolsep) * \real{0.4487}}
  >{\raggedright\arraybackslash}p{(\columnwidth - 6\tabcolsep) * \real{0.2081}}
  >{\raggedright\arraybackslash}p{(\columnwidth - 6\tabcolsep) * \real{0.1795}}@{}}
\toprule\noalign{}
\endhead
\bottomrule\noalign{}
\endlastfoot
\textbf{Mã HTTP} & \textbf{Mã máy đọc} & \textbf{Tình huống} & \textbf{Nguồn} \\
400 & VALIDATION\_ERROR & Dữ liệu sai định dạng: đăng ký, tham số page và limit, thiếu Idempotency-Key, phương thức thanh toán lạ, dữ liệu test sai, thiếu hoặc sai from & FR-AUTH-001, FR-CAT-001, FR-PAY-001, mục 11.5 \\
400 & INVALID\_CART & Giỏ không hợp lệ hoặc trống; chỉ rõ dòng lỗi & FR-CART-005, FR-ORD-001 \\
400 & INVALID\_COUPON & Mã giảm giá không hợp lệ; kèm lý do & FR-PRO-001 \\
401 & UNAUTHORIZED & Không có token, token sai chữ ký hoặc hết hạn & FR-AUTH-003, FR-AUTHZ-001 \\
401 & INVALID\_CREDENTIALS & Đăng nhập sai email hoặc mật khẩu (một thông báo chung) & FR-AUTH-002 \\
403 & FORBIDDEN & Thiếu hoặc sai X-Internal-Key ở endpoint nội bộ & FR-AUTHZ-001, mục 11.5 \\
404 & NOT\_FOUND & Sản phẩm không tồn tại; đơn không tồn tại hoặc không phải của người gọi (hai trường hợp cùng phản hồi); đơn không tồn tại ở endpoint nội bộ & FR-CAT-002, FR-AUTHZ-002 \\
409 & EMAIL\_ALREADY\_EXISTS & Đăng ký với email đã tồn tại & FR-AUTH-001 \\
409 & INSUFFICIENT\_STOCK & Không đủ tồn kho; nêu sản phẩm thiếu, số yêu cầu và số còn có thể đặt & FR-INV-002 \\
409 & PENDING\_ORDER\_EXISTS & Khách đã có đơn PENDING; kèm mã đơn đang chờ & FR-ORD-006 \\
409 & COUPON\_UNAVAILABLE & Mã hết lượt dùng hoặc người này đã dùng đủ lượt & FR-PRO-004 \\
409 & INVALID\_TRANSITION & Hành động không hợp lệ ở trạng thái hiện tại của đơn (thanh toán, hủy, chuyển bước nội bộ) & FR-ORD-003, mục 11.2.1 \\
409 & IDEMPOTENCY\_KEY\_CONFLICT & Cùng khóa idempotency nhưng khác nội dung & FR-PAY-003 \\
409 & PAYMENT\_IN\_PROGRESS & Request cùng khóa đang được xử lý; gửi lại cùng khóa & FR-PAY-003 \\
200 & (trường trạng thái FAILED) & Thanh toán bị từ chối là kết quả nghiệp vụ, không phải lỗi HTTP & FR-PAY-002 \\
200 & (kèm dấu phát lại) & Phát lại cùng khóa và cùng nội dung; trả đúng kết quả cũ & FR-PAY-003 \\
5xx & INTERNAL\_ERROR & Lỗi kỹ thuật; thông điệp chung, không lộ chi tiết nội bộ; log ở server & FR-PAY-001, BR-041 \\
\end{longtable}

\textbf{Thứ tự ưu tiên khi một request vi phạm nhiều điều kiện:}

\begin{itemize}
\item
  POST /orders: xác thực (401) → dữ liệu giỏ và mã giảm giá (400) → đã có đơn chờ (409) → tồn kho (409 INSUFFICIENT\_STOCK) → lượt dùng mã (409 COUPON\_UNAVAILABLE) (BR-024).
\item
  POST /payments: xác thực (401) → dữ liệu request (400) → đơn thuộc người gọi (404) → tra khóa idempotency → trạng thái đơn (409) → xử lý (mục 11.3).
\item
  Sở hữu luôn được kiểm tra trước trạng thái, nên đơn của người khác không bao giờ nhận 409 (BR-011).
\end{itemize}

\textbf{15.2. Ma trận kịch bản}

\begin{longtable}[]{@{}
  >{\raggedright\arraybackslash}p{(\columnwidth - 4\tabcolsep) * \real{0.2200}}
  >{\raggedright\arraybackslash}p{(\columnwidth - 4\tabcolsep) * \real{0.3000}}
  >{\raggedright\arraybackslash}p{(\columnwidth - 4\tabcolsep) * \real{0.4800}}@{}}
\toprule\noalign{}
\begin{minipage}[b]{\linewidth}\raggedright
\textbf{Nhóm}
\end{minipage} & \begin{minipage}[b]{\linewidth}\raggedright
\textbf{Kịch bản}
\end{minipage} & \begin{minipage}[b]{\linewidth}\raggedright
\textbf{Cách xử lý dự kiến ở mức yêu cầu}
\end{minipage} \\
\midrule\noalign{}
\endhead
\bottomrule\noalign{}
\endlastfoot
Authentication & Đăng ký với email đã tồn tại & 409 EMAIL\_ALREADY\_EXISTS; không tạo tài khoản; chấp nhận để lộ việc email tồn tại ở đăng ký (FR-AUTH-001). \\
Authentication & Đăng ký với dữ liệu sai (email sai định dạng, mật khẩu dưới 8 ký tự, thiếu trường) & 400 VALIDATION\_ERROR nêu từng trường; không tạo tài khoản. \\
Authentication & Đăng ký hoặc gửi request kèm trường lạ (role, điểm tích lũy) & Bỏ qua trường lạ; tài khoản mới có điểm bằng 0 (BR-003, BR-007). \\
Authentication & Đăng nhập sai email hoặc sai mật khẩu & 401 INVALID\_CREDENTIALS; hai trường hợp cho cùng thông báo và phản hồi; không cấp token (FR-AUTH-002). \\
Authentication & Token thiếu, sai chữ ký hoặc hết hạn & 401 UNAUTHORIZED; Web Client chuyển tới trang đăng nhập và giữ giỏ hàng (FR-AUTH-003). \\
Authentication & Đăng xuất chủ động & Client xóa token và giỏ; token cũ dùng được tới khi hết hạn vì không thu hồi phía server (FR-AUTH-004). \\
Authorization & Customer A xem, thanh toán hoặc hủy đơn của Customer B & 404 NOT\_FOUND, giống hệt đơn không tồn tại; không lộ trạng thái đơn (FR-AUTHZ-002). \\
Authorization & Customer gọi endpoint nội bộ, kể cả có JWT hợp lệ & 403 FORBIDDEN (FR-AUTHZ-001). \\
Authorization & Gọi API trực tiếp (Postman) khi không đủ điều kiện mà giao diện đã ẩn nút & Vẫn bị từ chối; backend là nơi quyết định cuối cùng (FR-AUTHZ-003). \\
Authorization & Request kèm giá, tổng tiền, số tiền giảm, điểm hoặc kết quả thanh toán & Bỏ qua; dùng giá trị server tính (BR-003). \\
Authorization & Endpoint mới chưa khai báo công khai & Yêu cầu xác thực; không mở mặc định (FR-AUTHZ-001). \\
Catalog & Chi tiết sản phẩm với mã sai hoặc không tồn tại & 404 NOT\_FOUND; giao diện hiển thị ``không tìm thấy món'' (FR-CAT-002). \\
Catalog & page hoặc limit sai (0, âm, chữ, limit lớn hơn 50) & 400 VALIDATION\_ERROR; page vượt totalPages thì 200 với danh sách rỗng, giao diện có nút về trang 1 (FR-CAT-001). \\
Catalog & Menu rỗng vì chưa nạp dữ liệu mẫu & 200 danh sách rỗng; giao diện hiển thị trạng thái trống; README hướng dẫn nạp seed (FR-CAT-005). \\
Catalog & Món có tồn kho bằng 0 & Hiển thị hết hàng; không thêm vào giỏ qua giao diện; gọi trực tiếp tạo đơn thì 409 INSUFFICIENT\_STOCK (FR-CAT-004). \\
Catalog & Giá catalog đổi giữa lúc khách xem và lúc đặt & Backend dùng giá tại thời điểm tạo đơn; giao diện hiển thị lại tổng cuối từ phản hồi trước thanh toán (FR-CAT-003, NFR-UX-003). \\
Catalog & Ảnh sản phẩm lỗi hoặc thiếu & Dùng ảnh thay thế; giao diện không vỡ (FR-CAT-001). \\
Cart & Thêm cùng một cấu hình nhiều lần, cộng dồn vượt 10 & Gộp dòng; giữ ở 10 và thông báo (FR-CART-001). \\
Cart & Giỏ đã có 20 dòng & Chặn thêm cấu hình mới; thông báo rõ (FR-CART-001). \\
Cart & Giảm số lượng khi đang là 1 & Xóa dòng như nút xóa; (SHOULD) hoàn tác; không bao giờ có số lượng 0 (FR-CART-002). \\
Cart & Dữ liệu localStorage bị sửa hoặc hỏng & Khi khôi phục, dòng không hợp lệ (số lượng nhỏ hơn 1, size hoặc topping sai) bị loại bỏ (FR-CART-004, BR-051). \\
Cart & Không dùng được localStorage (chế độ riêng tư) & Giỏ hoạt động trong bộ nhớ, mất khi tải lại, không báo lỗi (FR-CART-004). \\
Cart & Đóng tab hoặc mất mạng ngay sau khi thanh toán thành công, trước khi nhận phản hồi & Đơn đã PAID nhưng giỏ còn: hạn chế chấp nhận; khách xác nhận qua lịch sử đơn (FR-CART-004). \\
Cart & Sửa giỏ sau khi tạo đơn PENDING rồi thanh toán đơn cũ & Giỏ đã sửa bị xóa theo quy tắc PAID: hạn chế chấp nhận trong MVP (BR-023). \\
Order & Giỏ không hợp lệ (sản phẩm không còn, size lạ, topping sai hoặc trùng, quá 3 topping, số lượng ngoài 1 đến 10, hơn 20 dòng, hai dòng trùng cấu hình) & 400 INVALID\_CART chỉ rõ dòng lỗi; không tạo đơn, không đổi tồn kho, không giữ lượt mã; giỏ ở client giữ nguyên (FR-CART-005). \\
Order & Giỏ trống & 400 INVALID\_CART; giao diện vô hiệu nút đặt hàng (FR-CART-002). \\
Order & Đã có đơn PENDING, khách tạo đơn mới & 409 PENDING\_ORDER\_EXISTS kèm mã đơn chờ; giao diện đưa khách tới đơn đó để thanh toán hoặc hủy (FR-ORD-006). \\
Order & Hai request tạo đơn cùng lúc của một khách & Đúng một đơn; request còn lại 409 PENDING\_ORDER\_EXISTS (BR-025). \\
Order & Client gửi lại sau timeout trong khi request cũ đã commit & Không tạo đơn trùng: 409 PENDING\_ORDER\_EXISTS kèm mã đơn cũ (hoặc trả lại đơn đã tạo nếu có Idempotency-Key) (FR-ORD-006). \\
Order & Phản hồi tạo đơn bị thất lạc & Khách tìm lại đơn PENDING ở lịch sử đơn (FR-ORD-004). \\
Order & Lỗi kỹ thuật giữa chừng khi tạo đơn & Hoàn tác toàn bộ: không có đơn, không trừ tồn, không giữ lượt mã; 5xx (BR-004). \\
Order & Token hết hạn khi đơn đang PENDING & 401 khi gọi tiếp; giữ giỏ; khách đăng nhập lại để thanh toán nếu chưa quá 15 phút, nếu không đơn tự hết hạn (mục 11.8). \\
Inventory & Số lượng đặt vượt tồn kho & 409 INSUFFICIENT\_STOCK nêu sản phẩm thiếu, số yêu cầu và số còn có thể đặt; từ chối toàn bộ đơn (FR-INV-002). \\
Inventory & Hai dòng cùng sản phẩm, tổng vượt tồn dù từng dòng đủ & 409 INSUFFICIENT\_STOCK; không đơn nào được lưu (FR-INV-002). \\
Inventory & Đơn có sản phẩm A đủ tồn nhưng B thiếu & Từ chối cả đơn; tồn kho của A không đổi (FR-INV-002). \\
Inventory & 10 khách đặt đồng thời món còn 5 & Đúng 5 đơn thành công, 5 đơn nhận 409; tồn kho cuối bằng 0, không âm (BR-032). \\
Inventory & Hai đơn cùng chứa A và B nhưng khai báo theo thứ tự ngược nhau, gửi đồng thời & Không đơn nào lỗi do deadlock (FR-INV-002). \\
Inventory & Hoàn tồn có thể xảy ra hai lần (hủy hai lần, tác vụ hết hạn chạy hai lần, thanh toán thất bại rồi hết hạn) & Hoàn đúng một lần; thao tác sau nhận 409 INVALID\_TRANSITION (FR-INV-003). \\
Inventory & Cập nhật trực tiếp làm tồn kho âm & Bị ràng buộc stock ≥ 0 ở cơ sở dữ liệu từ chối (FR-INV-001). \\
Promotion & Mã không tồn tại, chưa đến hạn, đã hết hạn, hoặc đơn dưới mức tối thiểu & 400 INVALID\_COUPON kèm lý do; không tự bỏ mã, không tạo đơn; mã không tồn tại dùng thông điệp chung (FR-PRO-001). \\
Promotion & Gửi hai mã trong một đơn & 400; mỗi đơn dùng tối đa một mã (BR-043). \\
Promotion & Mã hết lượt hoặc người này đã dùng đủ lượt & 409 COUPON\_UNAVAILABLE (FR-PRO-004). \\
Promotion & Đặt đồng thời khi mã còn lượt cuối & Số đơn thành công bằng số lượt còn; các đơn còn lại 409 COUPON\_UNAVAILABLE (FR-PRO-004). \\
Promotion & Vừa thiếu tồn kho vừa hết lượt mã & Trả INSUFFICIENT\_STOCK; không giữ lượt mã (FR-PRO-004). \\
Promotion & Đơn sử dụng mã giảm giá bị thanh toán thất bại & Tiếp tục giữ lượt sử dụng mã giảm giá (HELD), không trả lượt mã. Customer có thể thử lại thanh toán trên cùng Order nếu chưa hết hạn (FR-PRO-004, BR-046). \\
Promotion & Đơn sử dụng mã giảm giá bị hủy hoặc hết hạn và chuyển sang CANCELLED & Trả lượt sử dụng mã giảm giá đã giữ (RELEASED) đúng một lần trong cùng transaction với việc chuyển Order sang CANCELLED (FR-PRO-004, BR-046). \\
Promotion & Giảm giá lớn hơn hoặc bằng tạm tính & Số tiền giảm không vượt tạm tính; PERCENT không vượt mức giảm tối đa; tổng không âm (FR-PRO-002). \\
Promotion & Tổng tiền sau giảm bằng 0 & Vẫn đi qua thanh toán mock với số tiền 0 (BR-050, cần nhóm xác nhận). \\
Promotion & Xem trước giảm giá khác kết quả khi tạo đơn (mã hết lượt giữa hai lần gọi) & Kết quả cuối là kết quả khi tạo đơn; xem trước chỉ tham khảo, không giữ lượt (FR-PRO-001). \\
Payment & Thiếu Idempotency-Key, phương thức lạ hoặc dữ liệu test sai định dạng & 400 VALIDATION\_ERROR (FR-PAY-001). \\
Payment & Mock từ chối (thẻ kết thúc 0002, ví kết thúc 0000) & Trả HTTP 200 với kết quả Payment FAILED; Order chuyển sang PAYMENT\_FAILED.

Hệ thống tiếp tục giữ tồn kho và lượt sử dụng mã giảm giá đã giữ khi tạo Order; không cộng điểm tích lũy và không xóa giỏ hàng.

Customer được phép thử lại thanh toán trên cùng Order nếu chưa hết hạn. Lần thử mới sử dụng Idempotency-Key mới và tạo Payment mới, không tạo Order mới hoặc giữ thêm tài nguyên (FR-PAY-002, BR-039). \\
Payment & Bấm đúp hoặc gửi lại với cùng khóa và cùng nội dung & Một Payment; trả lại đúng kết quả cũ kèm dấu phát lại; tồn kho, điểm, lượt mã không đổi thêm; phát lại kết quả FAILED không hoàn tồn lần hai (FR-PAY-003). \\
Payment & Cùng khóa nhưng khác nội dung & 409 IDEMPOTENCY\_KEY\_CONFLICT; dữ liệu không đổi (FR-PAY-003). \\
Payment & 10 request song song cùng khóa & Một Payment, điểm cộng một lần; các request nhận kết quả đã lưu hoặc 409 PAYMENT\_IN\_PROGRESS rồi gửi lại (FR-PAY-003). \\
Payment & Khóa mới cho đơn đã thanh toán & 409 INVALID\_TRANSITION; không có Payment thứ hai (FR-PAY-004). \\
Payment & Thanh toán mới cho Order đã CANCELLED hoặc đã hết thời hạn thanh toán. & Từ chối lần thanh toán mới với HTTP 409; không tạo Payment mới và không thay đổi trạng thái Order, tồn kho, lượt mã giảm giá hoặc điểm tích lũy. Request phát lại cùng Idempotency-Key và cùng nội dung vẫn được xử lý theo FR-PAY-003. \\
Payment & Thanh toán đơn của người khác & 404 NOT\_FOUND (FR-AUTHZ-002). \\
Payment & Customer thử thanh toán lại trên Order PAYMENT\_FAILED chưa hết hạn & Cho phép chuyển PAYMENT\_FAILED → PENDING để thực hiện lần thanh toán mới với Idempotency-Key mới; tạo Payment mới theo kết quả Mock Payment; không trừ thêm tồn kho hoặc giữ thêm lượt mã giảm giá (FR-PAY-002, BR-039). \\
Payment & Hủy hoặc hết hạn xảy ra cùng lúc với thanh toán & Chỉ một thao tác thành công; trạng thái cuối nhất quán với tồn kho, lượt mã và điểm (BR-028). \\
Payment & Lỗi kỹ thuật khi đang thanh toán & HTTP 5xx; rollback toàn bộ transaction, không tạo Payment mới, không lưu Idempotency-Key của lần thử chưa hoàn tất và không thay đổi tồn kho, lượt mã giảm giá hoặc điểm tích lũy.

Order giữ nguyên trạng thái trước lần thử: PENDING nếu thanh toán lần đầu; PAYMENT\_FAILED nếu đang thử lại sau thất bại.

Customer có thể gửi lại request với cùng Idempotency-Key sau khi lỗi kỹ thuật được khắc phục (BR-041). \\
Payment & Mất kết nối sau khi gửi thanh toán & Khách xem lịch sử đơn; đơn còn PENDING thì gửi lại cùng khóa an toàn (FR-PAY-005). \\
Payment & Dữ liệu thẻ hoặc ví xuất hiện trong log hoặc cơ sở dữ liệu & Không được xảy ra; chỉ lưu tối đa 4 số cuối (FR-PAY-004, NFR-SEC-007). \\
Order lifecycle & Chuyển trạng thái ngoài ma trận & HTTP 409 INVALID\_TRANSITION; không thay đổi trạng thái Order hoặc dữ liệu liên quan.

Mọi chuyển trạng thái không được khai báo là hợp lệ trong State Machine tại mục 11.2 đều bị backend từ chối, bao gồm cả các yêu cầu chuyển trạng thái không đáp ứng điều kiện (guard). \\
Order lifecycle & Hủy đơn đã PAID hoặc đã CANCELLED & 409 INVALID\_TRANSITION; hủy đơn PENDING của mình thì thành công và hoàn tồn (FR-ORD-005). \\
Order lifecycle & Hủy đơn của người khác & 404 NOT\_FOUND (FR-AUTHZ-002). \\
Order lifecycle & Tác vụ hết hạn chạy hai lần hoặc chạy chậm & Tác vụ định kỳ chạy trùng hoặc chạy chậm không được làm hoàn tồn kho hay trả lượt mã giảm giá nhiều lần.

Order ở trạng thái PENDING hoặc PAYMENT\_FAILED đã hết hạn phải bị từ chối mọi lần thanh toán mới, kể cả khi tác vụ chưa chuyển Order sang CANCELLED.

Khi tác vụ thực hiện hủy thành công, hệ thống chuyển Order sang CANCELLED, hoàn tồn kho và trả lượt mã đã giữ đúng một lần trong cùng transaction.

Nếu tác vụ chạy chậm hoặc gặp lỗi, Order có thể tạm thời ngăn Customer tạo đơn mới. Hệ thống phải tiếp tục xử lý đơn quá hạn ở những lần quét tiếp theo (FR-ORD-007, mục 11.6). \\
Internal & Thiếu hoặc sai X-Internal-Key & 403 FORBIDDEN (FR-AUTHZ-001). \\
Internal & Gọi hai lần với cùng from & Lần hai 409; đơn không nhảy hai bước (mục 11.5). \\
Internal & Đơn ở PENDING, PAYMENT\_FAILED, CANCELLED, COMPLETED, hoặc from không khớp & 409 INVALID\_TRANSITION (mục 11.5). \\
Internal & Đơn không tồn tại, hoặc thiếu hay sai kiểu from & 404 NOT\_FOUND; 400 VALIDATION\_ERROR (mục 11.5). \\
Loyalty & Thanh toán thành công & Cộng đúng một lần phần nguyên của (tổng sau giảm ÷ 10.000); đơn dưới 10.000đ cho 0 điểm (FR-PRO-003). \\
Loyalty & Thanh toán thất bại, đơn hủy hoặc hết hạn & Điểm không đổi (FR-PRO-003). \\
Loyalty & Gửi lại thanh toán hoặc gửi song song cùng khóa & Không cộng thêm (FR-PAY-003). \\
Loyalty & Client gửi kèm số điểm & Bỏ qua; chỉ backend thay đổi điểm (BR-048). \\
Loyalty & Điểm hiển thị chưa cập nhật ngay sau thanh toán & Giao diện làm mới từ server khi vào trang xác nhận và trang tài khoản (FR-PRO-005). \\
Data & Đổi giá, tên món, topping hoặc mã giảm giá sau khi đặt & Đơn cũ vẫn hiển thị bản chụp; tổng không đổi (FR-ORD-002). \\
Data & Chèn Payment thứ hai cho một đơn, hoặc hai đơn PENDING của cùng một người & Cơ sở dữ liệu từ chối bằng ràng buộc (mục 10.2). \\
Data & Chạy seed hai lần & Không tạo dữ liệu trùng; không đặt lại số lượt mã đã dùng ngoài ý muốn (FR-CAT-005, FR-PRO-006). \\
Data & Mã giảm giá mẫu ``còn hạn'' hết hạn đúng ngày nghiệm thu & Thời hạn mã mẫu tính theo ngày nạp dữ liệu (FR-PRO-006). \\
System & Cơ sở dữ liệu không sẵn sàng hoặc lỗi không lường trước & 5xx INTERNAL\_ERROR thông điệp chung, không lộ chi tiết nội bộ; ghi log ở server; giao diện cho phép thử lại; không để lại dữ liệu nửa vời (NFR-REL-002). \\
System & Thiếu JWT\_SECRET hoặc X-Internal-Key trong cấu hình & Backend báo lỗi cấu hình rõ ràng và không khởi động; không dùng giá trị mặc định (NFR-SEC-004). \\
System & Chạy hệ thống từ máy sạch & Docker Compose chạy đủ frontend, backend, cơ sở dữ liệu theo README (NFR-OPS-001). \\
\end{longtable}

\end{document}
